\chapter{Artificial Neural Networks}\label{ch:neural_networks}

Artificial neural networks (\gls{ANN}) represent one of the earliest and most successful approaches to biologically inspired computation. Their development began in 1943 when Warren McCulloch and Walter Pitts proposed the first formal mathematical model of an artificial neuron~\citep{McCullochPitts1943}, establishing the conceptual foundations that endure to this day.

The central motivation of \gls{ANN} is to emulate, in a simplified manner, the information processing of biological neurons. These specialized cells receive input signals, integrate them through complex electrochemical processes, and generate output responses that propagate through interconnected networks. Crucially, biological neurons exhibit synaptic plasticity, adapting their responses based on experience and enabling learning from examples.

This analogy with the processing mechanisms of the human brain has motivated decades of research in artificial intelligence, seeking to develop computational systems capable of learning and adapting in a manner similar to biological systems. In this chapter we will examine the mathematical foundations of \gls{ANN}, from the individual neuron to modern deep architectures.

\section{Artificial Neuron}\label{sec:neurona_artificial}

\begin{figure}[h]
	\centering
	\includegraphics[width=0.6\textwidth]{figures/neuron_biological}
	\caption[Diagram of a biological neuron]{Diagram of a biological neuron showing dendrites (inputs), soma (cell body), and axon (output). This structure inspired the McCulloch-Pitts mathematical model.}
	\label{fig:neuron}
\end{figure}

The McCulloch-Pitts model established the conceptual foundations by abstracting the behavior of a biological neuron (\autoref{fig:neuron}) into simple mathematical terms. In the nervous system, neurons receive electrochemical signals from other cells through dendrites and synaptic connections. These signals are integrated in the cell soma, and if the weighted sum exceeds a critical depolarization threshold, the neuron generates an action potential that propagates through the axon toward other neurons.

Inspired by this mechanism, an artificial neuron constitutes a mathematical model that:
\begin{enumerate}
	\item Receives multiple numerical input signals
	\item Weights them through adjustable parameters (synaptic weights)
	\item Applies a nonlinear activation function that models the firing threshold
	\item Produces an output signal that can serve as input to other neurons
\end{enumerate}

It is important to recognize that real biological neurons exhibit a complexity that goes far beyond this initial abstraction. Neuronal processes involve complex temporal dynamics, specific chemical neurotransmitters, time-dependent synaptic plasticity, and intricate three-dimensional dendritic architectures. Nevertheless, the simplification of the McCulloch-Pitts model has proved to be extraordinarily fruitful, providing the conceptual basis for the development of modern \gls{ANN}.

\subsection{Mathematical formulation}

The standard mathematical model of an artificial neuron can be formulated as a function that maps an input vector \(\mathbf{x} = (x_1, x_2, \ldots, x_n)^T\) to a scalar output \(y\). The fundamental operation consists of:

\begin{equation}
	y = \sigma\left(\sum_{i=1}^{n} w_i x_i + b\right)
	\label{eq:neuron}
\end{equation}

where:
\begin{itemize}
	\item \(\mathbf{w} = (w_1, w_2, \ldots, w_n)^T\) are the synaptic weights that modulate the importance of each input
	\item \(b\) is the bias or activation threshold
	\item \(\sigma\) is the nonlinear activation function
\end{itemize}

The activation function \(\sigma\) plays a fundamental role by introducing the essential nonlinearity that allows neural networks to approximate complex functions and capture nonlinear relationships in the data. The synaptic weights \(\mathbf{w}\) and the bias \(b\) constitute the adjustable parameters of the model. These are updated during the training process to approximate the desired target function.

\subsubsection{Spiking Neural Networks}

It is important to mention that this classical formulation, although extremely successful, represents a simplified abstraction of neuronal behavior. There exist alternative models that are closer to biology, such as \textbf{Spiking Neural Networks} or \gls{SNN}, considered the ``third generation'' of neural networks \citep{Maass1997spiking}. Unlike the classical model where the output is a continuous value, \glspl{SNN} process information through discrete pulses (spikes) that propagate over time, thus capturing the temporal dynamics inherent to biological neurons. These networks use models such as \textit{Integrate-and-Fire} or \textit{Hodgkin-Huxley}, where the neuron accumulates membrane potential until reaching a threshold and emitting a spike. \glspl{SNN} offer significant computational advantages in terms of energy efficiency and are especially promising for neuromorphic hardware \citep{Tavanaei2018deep}. However, for the purposes of this work, we will focus on the classical artificial neuron model described above.

\subsubsection{Formal definition}

To establish a rigorous mathematical framework for this work, we first formalize the concept of activation function and, subsequently, that of artificial neuron.

\begin{definition}[Activation function]\label{def:funcion_activacion}
	An \textit{activation function} is a function $\sigma: \mathbb{R} \to \mathbb{R}$ that satisfies the following properties:
	\begin{itemize}
		\item \textbf{Nonlinearity}: $\sigma$ is not an affine map,
		\item \textbf{Continuity}: $\sigma$ is continuous,
		\item \textbf{Differentiability almost everywhere}: $\sigma$ is differentiable except possibly on a set of Lebesgue measure zero\footnote{A set $A \subset \mathbb{R}$ has Lebesgue measure zero if it can be covered by intervals whose total length is arbitrarily small. See Definitions~\ref{def:medida_lebesgue} and \ref{def:casi_por_doquier} for the formal definition.}.
	\end{itemize}
\end{definition}

These properties ensure that $\sigma$ can introduce the nonlinearities necessary for the expressive capacity of the network, while allowing the computation of gradients for training through backpropagation. \autoref{tab:funciones_activacion} summarizes the main activation functions used in practice along with their most relevant properties, and \autoref{subsubsec:funciones_activacion_capas} describes each of them in detail.

\begin{definition}[Artificial neuron]\label{def:neurona_artificial}
	Let $n \in \mathbb{N}$ and let $\sigma: \mathbb{R} \to \mathbb{R}$ be an activation function. For each $\theta = (\mathbf{w}, b) \in \mathbb{R}^{n+1}$ with $\mathbf{w} \in \mathbb{R}^n$ (weights) and $b \in \mathbb{R}$ (bias), an \textit{artificial neuron} is a function $f_{\theta}: \mathbb{R}^n \to \mathbb{R}$ expressed as the composition of an affine map with an activation function:
	\begin{equation}
		f_{\theta}(\mathbf{x}) = \sigma \circ T_{\theta}(\mathbf{x}) = \sigma\Big( \langle \mathbf{w}, \mathbf{x} \rangle + b \Big)
		\label{eq:neuron_formal}
	\end{equation}
	where $T_{\theta}: \mathbb{R}^n \to \mathbb{R}$ is the \textit{affine map} defined by
	\[
	T_{\theta}(\mathbf{x}) = \langle \mathbf{w}, \mathbf{x} \rangle + b.
	\]
\end{definition}

The set of parameters $\theta$ is trainable, which means that during the learning process it is adjusted to minimize a loss function. Thus, an artificial neuron defines a \textit{family of functions} $\{f_{\theta} : \theta \in \mathbb{R}^{n+1}\}$ parameterized by the weights and the bias.

\autoref{fig:neurona_artificial} illustrates the structure of an artificial neuron, showing the flow from the weighted inputs to the activated output.

\begin{figure}[h]
\centering
\begin{tikzpicture}[
    node distance=0.8cm,
    input/.style={circle, draw, minimum size=0.6cm, fill=gdlblue!15},
    weight/.style={rectangle, draw, minimum width=0.6cm, minimum height=0.4cm, fill=gdlgray!15, font=\scriptsize},
    sum/.style={circle, draw, minimum size=0.8cm, fill=gdlyellow!20},
    activation/.style={rectangle, draw, rounded corners, minimum width=0.8cm, minimum height=0.6cm, fill=gdlgreen!20},
    output/.style={circle, draw, minimum size=0.6cm, fill=gdlred!15},
    arrow/.style={->, thick, >=stealth}
]

% Entradas
\node[input] (x1) at (0, 1.5) {$x_1$};
\node[input] (x2) at (0, 0) {$x_2$};
\node[input] (xn) at (0, -1.5) {$x_n$};
\node at (0, -0.75) {$\vdots$};

% Pesos
\node[weight] (w1) at (1.5, 1.5) {$w_1$};
\node[weight] (w2) at (1.5, 0) {$w_2$};
\node[weight] (wn) at (1.5, -1.5) {$w_n$};

% Nodo sumador
\node[sum] (suma) at (3.5, 0) {$\sum$};
\node[below=0.1cm of suma, font=\scriptsize] {$+b$};

% Función de activación
\node[activation] (sigma) at (5.5, 0) {$\sigma$};

% Salida
\node[output] (y) at (7, 0) {$y$};

% Conexiones
\draw[arrow] (x1) -- (w1);
\draw[arrow] (x2) -- (w2);
\draw[arrow] (xn) -- (wn);
\draw[arrow] (w1) -- (suma);
\draw[arrow] (w2) -- (suma);
\draw[arrow] (wn) -- (suma);
\draw[arrow] (suma) -- node[above, font=\scriptsize] {$z$} (sigma);
\draw[arrow] (sigma) -- (y);

% Etiquetas de fases
\node[above=1.5cm of x1, font=\small\bfseries] {Inputs};
\node[above=0.3cm of w1, font=\small\bfseries] {Weights};
\node[above=0.5cm of suma, font=\small\bfseries] {Sum};
\node[above=0.5cm of sigma, font=\small\bfseries] {Activation};

\end{tikzpicture}
\caption[Structure of an artificial neuron]{Structure of an artificial neuron. The inputs $x_i$ are multiplied by weights $w_i$, summed together with the bias $b$, and the result $z = \sum w_i x_i + b$ is transformed by the activation function $\sigma$ to produce the output $y = \sigma(z)$.}
\label{fig:neurona_artificial}
\end{figure}

The nonlinearity of $\sigma$ is crucial: the composition of multiple affine maps remains an affine map, so that without nonlinear activation functions, a deep neural network with multiple layers would be equivalent to a single linear transformation, losing all its expressive capacity. It is this nonlinearity that allows neural networks to approximate arbitrarily complex functions, as established by the universal approximation \autoref{thm:aproximacion_universal} (see \autoref{sec:aproximacion_universal}).

\subsubsection{Desirable properties of activation functions}

Although \autoref{def:funcion_activacion} establishes the minimum properties necessary for a function to serve as an activation function, in practice it is desirable that it satisfy additional properties that facilitate training and improve the behavior of neural networks. Below we present a detailed analysis of these desirable properties.

\begin{definition}[Boundedness or controlled growth]\label{prop:acotacion_activacion}
	An activation function $\sigma: \mathbb{R} \to \mathbb{R}$ has \textit{controlled growth} if it satisfies one of the following conditions:
	\begin{enumerate}
		\item \textbf{Boundedness}: There exists $M > 0$ such that $|\sigma(z)| \leq M$ for all $z \in \mathbb{R}$.
		\item \textbf{Linear growth}: There exist constants $C_1, C_2 > 0$ such that $|\sigma(z)| \leq C_1|z| + C_2$ for all $z \in \mathbb{R}$.
	\end{enumerate}
\end{definition}

This property is crucial for numerical stability during training. The first condition (boundedness) guarantees that activations remain within a finite range regardless of the magnitude of the input, preventing any numerical overflow. The second condition (linear growth) is more permissive: it allows activations to grow, but at a rate no greater than that of the input itself. In both cases, the explosion of values during forward propagation is prevented, which otherwise could cause numerical overflows or excessively large gradients. \autoref{tab:funciones_activacion} shows that the usual activation functions satisfy one or the other condition; for detailed descriptions of each function, see \autoref{subsubsec:funciones_activacion_capas}.

\begin{definition}[Monotonicity]\label{prop:monotonia_activacion}
	A function $\sigma: \mathbb{R} \to \mathbb{R}$ is \textit{monotonic} (increasing) if for any $z_1, z_2 \in \mathbb{R}$ with $z_1 < z_2$ it holds that:
	\[
	\sigma(z_1) \leq \sigma(z_2)
	\]
\end{definition}

Monotonicity provides three fundamental advantages:
\begin{enumerate}
	\item \textbf{Order preservation}: Activations maintain the relative order of the inputs, facilitating interpretability.
	\item \textbf{Uniqueness of local optima}: In combination with convexity of the loss function, it simplifies the optimization landscape.
	\item \textbf{Smoothness of the loss surface}: It reduces the presence of spurious local minima in the parameter space.
\end{enumerate}

Most classical activation functions are monotonic (see \autoref{subsubsec:funciones_activacion_capas}). However, monotonicity is not a necessary condition for good empirical performance: there exist non-monotonic activation functions that exhibit competitive results in practice, which suggests that the theoretical advantages of monotonicity may be offset by other favorable properties.

\begin{definition}[Zero-centered]\label{prop:zero_centered}
	An activation function $\sigma: \mathbb{R} \to \mathbb{R}$ is \textit{zero-centered} if it satisfies:
	\begin{enumerate}
		\item $\sigma(0) = 0$ or $\sigma(0) \approx 0$
		\item It is symmetric with respect to the origin: $\sigma(-z) = -\sigma(z)$ for all $z \in \mathbb{R}$ (odd function)
	\end{enumerate}
\end{definition}

This property significantly accelerates the convergence of training. When activations are zero-centered, the gradients have a mean close to zero, which avoids systematic biases in the weight updates.

When an activation function does not satisfy this property (for example, if its range is not centered at the origin, if $\sigma(0) \neq 0$), the outputs of each layer present a systematic bias. This causes all the gradients with respect to the weights of the same layer to have the same sign, causing the ``zigzagging'' phenomenon during gradient descent: the parameter updates oscillate between orthogonal directions instead of converging directly toward the optimum.

\begin{definition}[Non-saturating gradient]\label{prop:gradiente_no_saturante}
	An activation function $\sigma: \mathbb{R} \to \mathbb{R}$ has a \textit{non-saturating gradient} if there exists $\delta > 0$ and a set $S \subseteq \mathbb{R}$ of infinite measure such that:
	\[
	|\sigma'(z)| \geq \delta \quad \text{for all } z \in S
	\]
\end{definition}

This property is fundamental to avoid the \textit{vanishing gradient} problem, one of the main obstacles in training deep networks (see the formal analysis in \autoref{subsubsec:backpropagation}).

\begin{corollary}
Every bounded activation function with horizontal asymptotes necessarily saturates.
\end{corollary}

\begin{proof}
If $\lim_{z \to \pm\infty} \sigma(z) = c_{\pm}$ with $c_+, c_-$ finite, then $\lim_{z \to \pm\infty} \sigma'(z) = 0$. For any $\delta > 0$, there exists $M > 0$ such that $|\sigma'(z)| < \delta$ for all $|z| > M$, so that the set $\{z : |\sigma'(z)| \geq \delta\} \subseteq [-M, M]$ has finite measure. By \autoref{prop:gradiente_no_saturante}, $\sigma$ cannot have a non-saturating gradient and, therefore, necessarily saturates.
\end{proof}

When $|z|$ is large, $|\sigma'(z)| \approx 0$, which causes the gradients to become exponentially small in deep layers during backpropagation. Mathematically, if we consider a network with $L$ layers where each layer has a saturated activation, the gradient of the loss with respect to the weights of the first layers will be proportional to $\prod_{\ell=1}^{L} \sigma'(z^{(\ell)})$, which can become numerically zero.

A function that maintains $|\sigma'(z)| \geq \delta > 0$ on a set of infinite measure allows the gradient to flow without attenuation through multiple layers. The optimal case corresponds to functions whose derivative is constant equal to 1 on a half-line, which completely eliminates saturation in that active region and facilitates the training of very deep networks. Moreover, if one requires $|\sigma'(z)| > 0$ for all $z \in \mathbb{R}$ (that is, a strictly nonzero gradient over the entire domain), the risk that certain neurons stop contributing to learning (a phenomenon known as ``dead neurons'') is eliminated.

\autoref{fig:vanishing_gradient} (see \autoref{subsubsec:backpropagation}) visually illustrates how the gradient vanishes exponentially in deep networks with saturating activations, compared to the stable behavior of non-saturating activations.

\begin{definition}[Computational efficiency]\label{prop:eficiencia_computacional}
	An activation function is \textit{computationally efficient} if both $\sigma(z)$ and its derivative $\sigma'(z)$ can be evaluated in $O(1)$ time with basic arithmetic operations.
\end{definition}

In deep neural networks with millions of neurons, the activation function is evaluated billions of times during training. Therefore, its computational efficiency has a direct impact on the total training time.

Taxonomy by type of operation:
\begin{itemize}
	\item \textbf{Comparisons and selections} (functions defined piecewise through $\max$): They require only comparison and selection operations, which makes them extremely efficient on modern hardware.
	\item \textbf{Transcendental functions} (exponentials, logarithms): They are evaluated in $O(1)$ but with a significantly larger constant than the previous ones, as they require internal polynomial approximations.
	\item \textbf{Vector functions} (normalization over the entire output vector): They have $O(n)$ complexity in the dimension of the vector, since they require computing global statistics before normalizing each component.
\end{itemize}

Although all these categories are polynomial in complexity, the differences in constants are significant in practice. The computationally simplest activation functions allow deeper networks to be trained in less time, which has been a relevant factor in the widespread adoption of comparison-based activations.

\begin{definition}[Lipschitz continuity]\label{prop:lipschitz_continuidad}
	A function $\sigma: \mathbb{R} \to \mathbb{R}$ is \textit{Lipschitz continuous} if there exists a constant $L > 0$ (called the \textit{Lipschitz constant}) such that for any $z_1, z_2 \in \mathbb{R}$:
	\[
	|\sigma(z_1) - \sigma(z_2)| \leq L |z_1 - z_2|
	\]
\end{definition}

Lipschitz continuity is stronger than uniform continuity and provides quantitative control over how perturbations in the input propagate to the output. This has certain advantages in the training of neural networks:

\begin{enumerate}
	\item \textbf{Numerical stability}: Small changes in the weights produce bounded changes in the activations.
	\item \textbf{Gradient control}: If $\sigma$ is Lipschitz with constant $L$, then $|\sigma'(z)| \leq L$ almost everywhere (where the derivative exists).
	\item \textbf{Convergence guarantees}: Many convergence theorems for optimization algorithms assume Lipschitz continuity of the objective function.
\end{enumerate}

For activation functions that are differentiable almost everywhere, the Lipschitz constant can be computed as $L = \sup_{z} |\sigma'(z)|$. Bounded functions with horizontal asymptotes tend to have small Lipschitz constants (since their derivative vanishes at the extremes and reaches a moderate maximum), while functions with linear growth typically present $L = 1$. \autoref{tab:funciones_activacion} collects the concrete values for the most common activation functions.

It should be noted that a function can be Lipschitz continuous without being of class $C^1$: it suffices that it be continuous and have a bounded derivative almost everywhere. At points of non-differentiability, subgradients are used, which is sufficient for the gradient-based optimization algorithms employed in the training of neural networks.

The properties presented are not mutually exclusive, and the most successful activation functions in practice usually combine several of them. \autoref{tab:funciones_activacion} summarizes the properties of the most widely used activation functions.

\begin{table}[h]
\centering
\caption{Comparison of common activation functions}
\label{tab:funciones_activacion}
\begin{tabular}{@{}lccccc@{}}
\toprule
\textbf{Function} & \textbf{Formula} & \textbf{Range} & \textbf{Centered} & \textbf{Non-saturating} & \textbf{Lipschitz} \\
\midrule
Sigmoid & $\frac{1}{1+e^{-z}}$ & $(0,1)$ & No & No & $L=\frac{1}{4}$ \\
\addlinespace
Tanh & $\frac{e^z-e^{-z}}{e^z+e^{-z}}$ & $(-1,1)$ & Yes & No & $L=1$ \\
\addlinespace
ReLU & $\max(0,z)$ & $[0,\infty)$ & No & Yes ($z>0$) & $L=1$ \\
\addlinespace
Leaky ReLU & $\max(\alpha z,z)$ & $\mathbb{R}$ & No & Yes & $L=1$ \\
\addlinespace
Swish & $z \cdot \sigma(z)$ & $\approx(-0.3,\infty)$ & No & Partial & $L\approx 1.1$ \\
\bottomrule
\end{tabular}
\end{table}

\autoref{fig:activation_functions} graphically shows the most common activation functions together with their derivatives, allowing one to visualize the saturation zones and the behavior of the gradient.

\begin{figure}[htbp]
\centering
\includegraphics[width=0.95\textwidth]{figures/activation_functions}
\caption[Plots of activation functions]{Common activation functions and their derivatives. The shaded zones indicate regions of saturation where the gradient tends to zero. Note how ReLU maintains a constant gradient for $z > 0$, while Sigmoid and Tanh saturate at the extremes.}
\label{fig:activation_functions}
\end{figure}

A fundamental theoretical result that justifies the use of activation functions with these properties is the \textbf{universal approximation theorem} (see \autoref{sec:aproximacion_universal}), which establishes that neural networks with a single hidden layer and appropriate activation functions can approximate any continuous function on a compact set with arbitrary precision. This theorem will be presented and proved in detail later in this chapter, together with the functional analysis preliminaries necessary for its complete understanding.

The choice of the appropriate activation function depends on the specific problem, the network architecture, and practical training considerations. In the following sections we will see how these properties influence architecture design.

\section{Data and its structure}\label{sec:datos_estructura}

In the previous section we defined the artificial neuron as a function that processes input data and generates output data. However, in this formulation we have not spoken about the information or the data that the neural network processes; we have not stopped to think about what the space $\mathbb{R}^n$ really represents. It is necessary to understand whether the data in the space $\mathbb{R}^n$ are simply numbers or whether there is some kind of relationship between them that contains relevant information.

This distinction is crucial. When we represent a $28 \times 28$ pixel image as a vector in $\mathbb{R}^{784}$, we are not dealing with $784$ independent values: each component corresponds to a specific pixel in a two-dimensional grid, and spatially neighboring pixels are naturally related. The grid structure is not arbitrary, but intrinsic to the nature of the image.

In this section we will formalize this distinction through two complementary conceptual frameworks that will allow us to unify the treatment of classical Euclidean data (vectors, images, sequences) with complex geometric data (graphs, manifolds, point clouds). In the following subsections we will present two mathematically equivalent frameworks but with distinct conceptual emphases: the \textit{domain-structure-signal} triple $(\Omega, \mathcal{S}, f)$ and the \textit{geometry-attributes} pair $(\mathcal{D}, f)$.

\subsection{Signals over structured domains}\label{sec:senales_dominios}

The first conceptual framework is based on three fundamental components that clearly separate the notions of positions, relations, and values. This tripartite decomposition constitutes the basis of the unifying paradigm of \gls{GDL} developed in \autoref{ch:gdl}, where the ``five G's'' of \citet{bronstein2017geometricdeeplearning} (geometry, groups, graphs, geodesics, and gauges) extend and enrich these fundamental concepts.

\begin{definition}[Data domain]\label{def:dominio_datos}
A \textit{data domain} is a set $\Omega$ whose elements represent the positions or locations where the data are observed. It can be:
\begin{itemize}
	\item \textbf{Discrete and finite}: $\Omega = \{p_1, \ldots, p_n\}$ (vertices of a graph, pixels of an image, indices of a vector)
	\item \textbf{Discrete and infinite}: $\Omega = \mathbb{Z}^d$ (infinite grid)
	\item \textbf{Continuous}: $\Omega \subseteq \mathbb{R}^d$ or a differentiable manifold $\mathcal{M}$
\end{itemize}
\end{definition}

\begin{definition}[Domain structure]\label{def:estructura_dominio}
The \textit{structure} $\mathcal{S}$ over a domain $\Omega$ defines the intrinsic relations between its elements. It can be of a:
\begin{itemize}
	\item \textbf{Topological} nature: A topology $\tau$ that defines open sets and notions of neighborhood
	\item \textbf{Metric} nature: A distance function $d: \Omega \times \Omega \to \mathbb{R}_{\geq 0}$ that quantifies separations between points
	\item \textbf{Algebraic} nature: Symmetries represented through a group $G$ that acts on $\Omega$
	\item \textbf{Combinatorial} nature: A graph $(V, E)$ that defines explicit connections between points
	\item \textbf{Geometric} nature: A differentiable manifold with atlas, Riemannian metric, connection, etc.
\end{itemize}

The pair $(\Omega, \mathcal{S})$ constitutes a \textit{structured domain}.
\end{definition}

Each type of structure over the domain corresponds to a branch of mathematics that provides specific tools for its analysis. In \autoref{ch:gdl} we will delve deeper into these structures and into how the symmetries of the domain (formalized through groups) determine the equivariance constraints that neural architectures must satisfy.

\begin{definition}[Signal]\label{def:senal}
A \textit{signal} over a structured domain $(\Omega, \mathcal{S})$ is a function
\[
f: \Omega \to \mathcal{C}
\]
that assigns to each point $p \in \Omega$ a feature vector $f(p) \in \mathcal{C}$, where $\mathcal{C}$ is a vector space called the \textit{channel space}. Typically $\mathcal{C} = \mathbb{R}^d$, where $d$ is the number of channels or attributes.

When $\Omega$ is finite, $\Omega = \{p_1, \ldots, p_n\}$, and $\mathcal{C} = \mathbb{R}^d$, we can represent the signal through the matrix:
\[
F = \begin{bmatrix}
	f(p_1)^T \\
	f(p_2)^T \\
	\vdots \\
	f(p_n)^T
\end{bmatrix} \in \mathbb{R}^{n \times d}
\]
where each row $f(p_i) \in \mathbb{R}^d$ contains the features observed at the position $p_i$.
\end{definition}

\begin{remark}\label{rem:espacio_canales_general}
The choice of $\mathcal{C} = \mathbb{R}^d$ is sufficient for most of the applications we will treat in this chapter. However, in later chapters the channel space acquires a richer structure: in \autoref{ch:gdl}, $\mathcal{C}$ can be a representation space of a symmetry group $G$ (for example, spherical harmonics for $SO(3)$), and in the context of gauge-equivariant networks, $\mathcal{C}$ corresponds to the fiber of an associated vector bundle over the base manifold.
\end{remark}

\begin{definition}[Signal space]\label{def:espacio_senales}
Given a structured domain $(\Omega, \mathcal{S})$ and a channel space $\mathcal{C}$, the \textit{signal space} $\mathcal{F}(\Omega, \mathcal{C})$ is the set of all signals over $\Omega$ with values in $\mathcal{C}$:
\[
\mathcal{F}(\Omega, \mathcal{C}) = \{f: \Omega \to \mathcal{C}\}.
\]
When $\Omega$ is finite with $|\Omega| = n$ and $\mathcal{C} = \mathbb{R}^d$, the signal space is isomorphic to a matrix space: $\mathcal{F}(\Omega, \mathbb{R}^d) \cong \mathbb{R}^{n \times d}$, where each signal is identified with its matrix representation $F$. When $\Omega$ is continuous, the space $\mathcal{F}(\Omega, \mathcal{C})$ is restricted to a function space with appropriate regularity conditions (for example, $L^2(\Omega, \mathcal{C})$ or $C^k(\Omega, \mathcal{C})$).
\end{definition}

A neural network can then be viewed as an operator between signal spaces. A layer transforms input signals into output signals, possibly over different domains and channel spaces:
\[
T: \mathcal{F}(\Omega, \mathcal{C}) \to \mathcal{F}(\Omega', \mathcal{C}').
\]
This functional perspective will be fundamental in the chapters on convolutional networks (\autoref{ch:cnn}), where convolution is an equivariant linear operator between signal spaces, and on attention mechanisms (\autoref{ch:attention}), where attention implements an adaptive nonlinear operator.

In summary:
\begin{itemize}
	\item The \textbf{domain} $\Omega$: defines \textit{where} the data are (the positions)
	\item The \textbf{structure} $\mathcal{S}$: defines \textit{how} the positions \textit{relate} to one another
	\item The \textbf{signal} $f$: defines the \textit{observed values} at each position
\end{itemize}

\autoref{fig:dominio_estructura_senal} visually illustrates this conceptual framework and how the three components combine to form complete data.

\begin{figure}[h]
\centering
\begin{tikzpicture}[
    scale=0.75,
    transform shape,
    box/.style={rectangle, draw, rounded corners, minimum width=3cm, minimum height=1.2cm, align=center},
    arrow/.style={->, thick, >=stealth},
    label/.style={font=\footnotesize, align=center}
]

% Dominio
\node[box, fill=gdlblue!15] (dom) at (0, 2) {\textbf{Domain} $\Omega$\\{\scriptsize Positions/locations}};
\node[label, left=0.1cm of dom] {Pixels, vertices, points};

% Estructura
\node[box, fill=gdlgreen!15] (str) at (6, 2) {\textbf{Structure} $\mathcal{S}$\\{\scriptsize Relations between positions}};
\node[label, right=0.1cm of str] {Grid, graph, metric};

% Geometría
\node[box, fill=gdlyellow!20] (geo) at (3, 0) {\textbf{Geometry} $(\Omega, \mathcal{S})$\\{\scriptsize Structured domain}};

% Señal
\node[box, fill=gdlorange!18] (sig) at (8, 0) {\textbf{Signal} $f: \Omega \to \mathcal{C}$\\{\scriptsize Observed values}};
\node[label, right=0.1cm of sig] {RGB, features, attributes};

% Datos completos
\node[box, fill=gdlred!18, minimum width=4cm] (dat) at (5.5, -2) {\textbf{Complete data}\\{\scriptsize $(\Omega, \mathcal{S}, f)$}};

% Flechas
\draw[arrow] (dom) -- (geo);
\draw[arrow] (str) -- (geo);
\draw[arrow] (geo) -- (dat);
\draw[arrow] (sig) -- (dat);

\end{tikzpicture}
\caption[Domain-structure-signal conceptual framework]{Conceptual framework for the representation of data. The domain $\Omega$ defines the positions, the structure $\mathcal{S}$ defines the relations between positions, and the signal $f$ assigns values to each position. Together they form the complete data that a neural network processes.}
\label{fig:dominio_estructura_senal}
\end{figure}

A vector $\mathbf{x} \in \mathbb{R}^n$ is not merely a collection of $n$ real numbers: it represents a signal $f$ over a domain $\Omega$ with a specific structure $\mathcal{S}$. The relations between components $x_i$ and $x_j$ are not arbitrary, but are determined by the structure of the underlying domain.

\begin{example}[Audio signal (temporal)]\label{ex:senal_audio}
Let us consider an audio signal represented formally as:
\begin{itemize}
	\item \textbf{Domain}: $\Omega = [0, T] \subset \mathbb{R}$ (continuous temporal interval)
	\item \textbf{Structure}: $\mathcal{S} = $ total order of $\mathbb{R}$ with the standard topology
	\item \textbf{Signal}: $f(t) \in \mathbb{R}$ (sound amplitude at instant $t$)
\end{itemize}

At each time instant $t$, we observe an amplitude $f(t)$. The structure is the temporal order: $t_1 < t_2$ has causal meaning. Audio signals respect temporal locality: values at nearby instants are typically correlated.
\end{example}

\begin{example}[Digital image on a grid]\label{ex:imagen_grilla}
A digital image is structured as:
\begin{itemize}
	\item \textbf{Domain}: $\Omega = \{1,\ldots,H\} \times \{1,\ldots,W\}$
	\item \textbf{Structure}: $\mathcal{S} = $ rectangular 2D grid with local neighborhood (4- or 8-connected)
	\item \textbf{Signal}: $f(i,j) \in \mathbb{R}^3$ (RGB intensities of pixel $(i,j)$)
\end{itemize}

If we arbitrarily permute the order of pixels when vectorizing the image, the vector representation changes, but the image (as a signal over a 2D grid) remains invariant. The grid structure is intrinsic to the geometric object, not to its particular vector representation.
\end{example}

\begin{example}[Social network (graph)]\label{ex:red_social}
A social network is modeled through:
\begin{itemize}
	\item \textbf{Domain}: $\Omega = V = \{u_1, \ldots, u_n\}$ (set of users)
	\item \textbf{Structure}: $\mathcal{S} = (V, E)$ where $E \subseteq V \times V$ represents connections or friendships
	\item \textbf{Signal}: $f(u) \in \mathbb{R}^d$ (profile of user $u$: age, interests, activity, etc.)
\end{itemize}

There is no canonical order of users. The ``position'' of a user is not a numerical index, but its abstract identity in the graph. The relations are defined by the edges $E$, not by the proximity of indices in some arbitrary vectorization. Two users $u_i, u_j$ are ``neighbors'' if and only if $(u_i, u_j) \in E$.
\end{example}

\begin{example}[Molecule (3D geometric graph)]\label{ex:molecula}
A molecule combines discrete structure and continuous geometry:
\begin{itemize}
	\item \textbf{Domain}: $\Omega = \{a_1, \ldots, a_n\}$ (set of atoms)
	\item \textbf{Structure}: $\mathcal{S} = $ (graph of chemical bonds) $+$ (positions in $\mathbb{R}^3$) $+$ (symmetries under the Euclidean group $E(3)$)
	\item \textbf{Signal}: $f(a) \in \mathbb{R}^d$ (atomic type, partial charge, atomic mass, hybridization, etc.)
\end{itemize}

The molecule combines discrete connectivity (covalent chemical bonds) with continuous geometry (three-dimensional positions of atoms). The physicochemical properties are invariant under transformations of the Euclidean group $E(3)$: translations, rotations, and reflections do not change the molecular identity.
\end{example}

\begin{example}[3D point cloud]\label{ex:nube_puntos}
A three-dimensional point cloud is characterized by:
\begin{itemize}
	\item \textbf{Domain}: $\Omega = \{p_1, \ldots, p_n\} \subset \mathbb{R}^3$ (finite set of points in space)
	\item \textbf{Structure}: $\mathcal{S} = $ Euclidean metric $d(p_i, p_j) = \|p_i - p_j\|_2$
	\item \textbf{Signal}: $f(p) \in \mathbb{R}^d$ (surface normals, curvatures, colors, reflectance, LIDAR intensity, etc.)
\end{itemize}

There is no natural order of the points. Permuting the list $\{p_1,\ldots,p_n\}$ does not change the point cloud as a three-dimensional geometric object. The structure is purely metric: only the relative distances between points matter, not their order of enumeration.
\end{example}

\begin{example}[3D surface (Riemannian manifold)]\label{ex:superficie_riemanniana}
A 3D surface as a Riemannian manifold is defined through:
\begin{itemize}
	\item \textbf{Domain}: $\Omega = \mathcal{M}$ (2D differentiable manifold embedded in $\mathbb{R}^3$, for example, the surface of a protein or a 3D object)
	\item \textbf{Structure}: $\mathcal{S} = $ Riemannian metric $g$ induced by the embedding in $\mathbb{R}^3$
	\item \textbf{Signal}: $f: \mathcal{M} \to \mathbb{R}^d$ (scalar or vector fields over the surface: temperature, hydrophilicity, Gaussian curvature, etc.)
\end{itemize}

The geodesic distances over $\mathcal{M}$ and the curvature are intrinsic properties of the manifold, independent of the particular embedding in $\mathbb{R}^3$. Two different parameterizations of the same surface define the same intrinsic geometry.
\end{example}

\begin{example}[Text sequence (NLP)]\label{ex:secuencia_texto}
A text sequence for natural language processing is represented as:
\begin{itemize}
	\item \textbf{Domain}: $\Omega = \{1, 2, \ldots, L\}$ (discrete positions in a sequence of length $L$)
	\item \textbf{Structure}: $\mathcal{S} = $ total order with bidirectional sequential structure (each position $i$ can attend to previous and subsequent positions)
	\item \textbf{Signal}: $f(i) \in \mathbb{R}^d$ (vector embedding of the token at position $i$: semantic representation of words, subwords, or characters)
\end{itemize}

Unlike one-dimensional audio with strict causality, text allows bidirectional dependencies: the meaning of a word depends on both the preceding and the subsequent context. \gls{Transformer} architectures exploit this structure through attention mechanisms that allow each position to interact with all the others, respecting permutation invariance except for the positional encodings that preserve the sequential order.
\end{example}

\begin{remark}[Permutation invariance]\label{rem:invarianza_permutacion}
The previous examples reveal a fundamental distinction. In some domains (audio, text, image) there exists a canonical or partial order of the positions, so that the matrix representation $F \in \mathbb{R}^{n \times d}$ is essentially unique. In others (graphs, point clouds), no such order exists: if $P \in \{0,1\}^{n \times n}$ is a permutation matrix, then $F$ and $PF$ represent the same signal over the same domain, simply with the positions enumerated in a different order. Formally, the signal $f: \Omega \to \mathcal{C}$ is intrinsic to the domain, but its matrix representation depends on an arbitrary choice of enumeration $\Omega = \{p_1, \ldots, p_n\}$.

This observation has profound consequences for the design of neural architectures: a network that processes data without a canonical order must produce results that are invariant (for classification tasks) or equivariant (for per-node prediction tasks) under permutations of the input. This is one of the central ideas of \gls{GDL} that we will formalize in \autoref{ch:gdl}.
\end{remark}

\autoref{tab:resumen_ejemplos_senales} summarizes the characteristics of the previous examples, highlighting the diversity of domains and the presence or absence of a canonical order.

\begin{table}[htbp]
\centering
\caption[Summary of examples of signals over structured domains]{Comparative summary of the examples of signals over structured domains. The column ``Canonical order?'' indicates whether the positions of the domain have a natural ordering, which determines whether the matrix representation is unique up to permutations.}
\label{tab:resumen_ejemplos_senales}
\begin{tabular}{lllcl}
\toprule
\textbf{Example} & $\Omega$ & \textbf{Type of} $\mathcal{S}$ & $\dim(\mathcal{C})$ & \textbf{Order?} \\
\midrule
Audio & $[0,T]$ & Total order & 1 & Yes \\
Image & 2D grid & Rectangular grid & 3 (RGB) & Yes (partial) \\
Social network & Vertices & Graph & $d$ & No \\
Molecule & Atoms & Graph $+ E(3)$ & $d$ & No \\
3D cloud & Points $\subset \mathbb{R}^3$ & Euclidean metric & $d$ & No \\
Surface & $\mathcal{M}$ & Riemannian & $d$ & No \\
Text & $\{1,\ldots,L\}$ & Sequential & $d$ & Yes \\
\bottomrule
\end{tabular}
\end{table}

\subsection{Geometry and attributes}\label{sec:geometria_atributos}

The previous framework separates domain, structure, and signal into three distinct components. However, it can be conceptually useful to unify the domain and its structure into a single entity: the \textit{geometry} of the problem.

\begin{definition}[Domain geometry]\label{def:geometria_dominio}
The \textit{geometry} of a data problem is the pair $\mathcal{D} = (\Omega, \mathcal{S})$ that combines the domain $\Omega$ with its structure $\mathcal{S}$ into a single entity that completely characterizes the spatial, topological, or algebraic relations between positions.
\end{definition}

\begin{definition}[Attributes]\label{def:atributos}
The \textit{attributes} (or features) are the image, the signal $f: \Omega \to \mathcal{C}$ that assigns observed values to each position of the domain.
\end{definition}

The two conceptual frameworks are mathematically equivalent\label{rem:equivalencia_marcos}. The equivalence is direct: $(\Omega, \mathcal{S}, f) \leftrightarrow (\mathcal{D}, f)$ where $\mathcal{D} = (\Omega, \mathcal{S})$.

The geometry/attributes framework offers an important conceptual advantage by distinguishing between intrinsic structure and variable observations:

\begin{itemize}
	\item The \textbf{geometry} $\mathcal{D} = (\Omega, \mathcal{S})$ is typically \textit{fixed} or \textit{intrinsic to the individual datum}: it defines the ``shape'' of the data space.

	\item The \textbf{attributes} $f$ are \textit{variable}: different data instances share the same geometry but have different attributes.
\end{itemize}

This distinction is fundamental for the design of neural architectures: the operations must respect the geometry (which is intrinsic) while transforming the attributes (which are variable). Ignoring the geometry, as a simple neuron does with a global inner product, wastes valuable structural information.

\subsection{Neural processing of signals}\label{sec:neurona_senales}

Having established the framework of signals over structured domains, we can now reinterpret the artificial neuron defined earlier and analyze its limitations from a geometric perspective.

An artificial neuron $f_\theta: \mathbb{R}^n \to \mathbb{R}$ can be reinterpreted as an operator that processes signals over a discrete domain:

\begin{itemize}
	\item \textbf{Input domain}: $\Omega = \{1, 2, \ldots, n\}$ (finite set of indices with no particular geometric meaning)
	\item \textbf{Structure}: None significant (set without geometric, topological, or algebraic relations)
	\item \textbf{Input signal}: $x: \Omega \to \mathbb{R}$ where $x(i) = x_i$ represents the value at index $i$
	\item \textbf{Processing}: $y = \sigma(\langle \mathbf{w}, \mathbf{x} \rangle + b) = \sigma\left(\sum_{i=1}^n w_i \cdot x(i) + b\right)$
\end{itemize}

The linear transformation $\langle \mathbf{w}, \mathbf{x} \rangle = \sum_{i=1}^n w_i x_i$ combines all the components globally through a weighted sum. There is no notion of:

\begin{itemize}
	\item \textbf{Locality}: All positions $i \in \Omega$ contribute equally to the output; there is no concept of ``local processing'' based on neighborhood

	\item \textbf{Structured neighborhood}: No distinction is made between ``nearby'' and ``distant'' indices; the connection $w_1 x_1$ has the same structural character as $w_{100} x_{100}$

	\item \textbf{Geometric invariance}: Permuting the indices through $\pi: \{1,\ldots,n\} \to \{1,\ldots,n\}$ produces a different signal $x_\pi$ where $x_\pi(i) = x(\pi(i))$, and typically $f_\theta(x_\pi) \neq f_\theta(x)$, even though semantically the permutation may be irrelevant (for example, in graphs)
\end{itemize}

The neuron treats the domain $\Omega = \{1,\ldots,n\}$ as a set without structure, ignoring any geometry that may exist implicitly in the nature of the data.

\glspl{ANN} focus on transformations of the signal or attributes between layers, leaving the geometry in the background, or using it solely to optimize computational calculations. In \gls{GDL} we will see that the geometry of the data takes on greater relevance and is one of the fundamental pillars. The concrete examples of different types of geometries (Euclidean, of graphs, manifolds) and their relationship with neural architectures are developed in depth in the Geometric Deep Learning chapter (\autoref{ch:gdl}), where the complete mathematical framework is presented.

\section{Layers of a neural network}

In \autoref{def:neurona_artificial}, we presented an \textit{artificial neuron} as a function with a single scalar output:
\[
f(x) = \sigma(w^T x + b): \mathbb{R}^n \to \mathbb{R}
\]
This choice is not arbitrary, but responds to several criteria:

\begin{enumerate}
    \item \textbf{Biological analogy}: Biological neurons emit a single output signal (firing rate), which propagates through the axon toward other neurons. The scalar definition captures this fundamental characteristic.

    \item \textbf{Mathematical simplicity}: A function $\mathbb{R}^n \to \mathbb{R}$ is conceptually simpler than a function $\mathbb{R}^n \to \mathbb{R}^m$, facilitating the initial analysis of the properties of an individual neuron (separation by hyperplane, nonlinear activation, etc.).
\end{enumerate}

Mathematically, we could have defined a neuron as a vector function $f: \mathbb{R}^n \to \mathbb{R}^m$ with multiple outputs. However, this generalization does not provide additional expressive power: a neuron with $m$ outputs is functionally equivalent to $m$ scalar neurons operating in parallel, each with its own weights and bias. Therefore, we adopt the scalar definition as the primitive, and we will build more complex architectures through the composition and parallelization of these basic units.

To obtain multiple output values, we use several neurons that share the same input domain. Each neuron independently transforms the input signal into a single scalar value, and the set of the $m$ outputs forms a new signal that, in general, resides over a domain different from the input that may have a different structure. This organization leads us to formalize the notion of a \textit{neural layer}.

\begin{definition}[Neural layer]
\label{def:capa_neuronal}
A \textit{neural layer} is a set of $m$ neurons that share the same input domain and operate in parallel over the input signal. Given an input signal $\mathbf{x} \in \mathbb{R}^n$ (where $n$ is the dimension of the input domain), the layer produces an output signal $\mathbf{y} \in \mathbb{R}^m$ through:
\[
L(\mathbf{x}) = \sigma(\mathbf{W}\mathbf{x} + \mathbf{b})
\]
where:
\begin{itemize}
    \item $\mathbf{W} \in \mathbb{R}^{m \times n}$ is the \textit{weight matrix}, whose $i$-th row contains the weights $w_i^T$ of neuron $i$
    \item $\mathbf{b} \in \mathbb{R}^m$ is the \textit{bias vector}
    \item $\sigma$ is the activation function, applied component-wise
\end{itemize}
Each neuron $i$ computes $y_i = \sigma(w_i^T \mathbf{x} + b_i)$, transforming the attributes of the input signal into a single output attribute.
\end{definition}

\paragraph{Connectivity and sparse matrices.}

 The matrix representation $\mathbf{W}\mathbf{x} + \mathbf{b}$ is completely general. If a neuron $i$ only processes a subset of the positions of the input domain, this is expressed through zero coefficients: $W_{ij} = 0$ when neuron $i$ does not receive a signal from position $j$. Although in practical implementations optimizations for sparse matrices are used, mathematically every neural layer admits this unified representation. This observation will be relevant when we study architectures with structured connectivity, such as convolutional networks.

\paragraph{Shared weights.}

\autoref{def:capa_neuronal} assumes that each neuron has its own independent parameters $(w_i, b_i)$. However, in many architectures the weights are \textit{shared} between neurons: multiple rows of the matrix $\mathbf{W}$ contain the same values. Mathematically, this imposes constraints of the form $W_{i,:} = W_{j,:}$ for certain pairs $(i,j)$, drastically reducing the number of free parameters.

\subsection{Composition of layers}

A neuron by itself is not capable of solving complex problems, but by combining many of them into a network one can learn more complex functions. As we have seen in \autoref{def:neurona_artificial}, a neuron is a function of the form $y = \sigma\left(\sum_{i=1}^n w_i x_i + b\right)$, where $\sigma$ is a nonlinear activation function. Geometrically, the linear part $\sum_{i=1}^n w_i x_i + b$ defines a hyperplane in $\mathbb{R}^n$, and the activation function introduces a nonlinearity that transforms this hyperplane. An individual neuron, therefore, can only separate the input space through a single nonlinearly transformed hyperplane. This drastically limits its \textit{expressive capacity}: it cannot approximate functions with multiple complex decision regions, hierarchical patterns, or highly nonlinear structures. For example, an individual neuron cannot solve the classic XOR problem, which requires a non-linearly-separable partition of the space.

When we connect multiple neurons organized into layers, we obtain a \textit{\gls{composicion_funcional}} of nonlinear transformations. We will denote by $L_i$ the $i$-th layer of the network, where each layer $L_i: \mathbb{R}^{n_{i-1}} \to \mathbb{R}^{n_i}$ transforms a signal of dimension $n_{i-1}$ into a signal of dimension $n_i$. A neural network of $k$ layers computes the composition:
\[
\mathcal{N} = L_k \circ L_{k-1} \circ \cdots \circ L_2 \circ L_1: \mathbb{R}^{n_0} \to \mathbb{R}^{n_k}
\]
where $n_0$ is the dimension of the input domain and $n_k$ that of the output domain.

This composition has fundamental mathematical properties:
\begin{enumerate}
    \item \textbf{Hierarchical representations}: Each composition $L_i \circ L_{i-1} \circ \cdots \circ L_1$ can be interpreted as a transformation that extracts increasingly abstract features from the data. The first layers detect simple patterns (edges, textures), while later layers combine these patterns into representations of higher semantic level (shapes, objects).

    \item \textbf{Complex partitioning of the space}: Multiple neurons in a layer generate multiple hyperplanes, allowing $\mathbb{R}^n$ to be partitioned into arbitrarily complex polyhedral regions. The composition of layers recursively refines these partitions, creating highly nonlinear decision boundaries.
\end{enumerate}

% ============ Versión TikZ con colores ============
\begin{figure}[h]
\centering
\begin{tikzpicture}[
    x=2.1cm, y=0.85cm,
    neuron/.style={circle, minimum size=5.5mm, line width=0.6pt},
    input/.style={neuron, draw=gdlblue!70, fill=gdlblue!15},
    hidden1/.style={neuron, draw=gdlgreen!60!black, fill=gdlgreen!15},
    hidden2/.style={neuron, draw=gdlgreen!70, fill=gdlgreen!18},
    hidden3/.style={neuron, draw=gdlpurple!70, fill=gdlpurple!18},
    output/.style={neuron, draw=gdlorange!70, fill=gdlorange!18},
    conn/.style={->, >=stealth, draw=gdlgray!35, line width=0.2pt},
    lbl/.style={font=\small\bfseries}
]

% Capa de entrada (3 neuronas)
\node[input] (I1) at (0, 2) {};
\node[input] (I2) at (0, 1) {};
\node[input] (I3) at (0, 0) {};
\node[left=1mm of I1, font=\small] {$x_1$};
\node[left=1mm of I2, font=\small] {$x_2$};
\node[left=1mm of I3, font=\small] {$x_3$};
\node[lbl] at (0, 3.5) {Input};

% Primera capa oculta (4 neuronas)
\node[hidden1] (H1-1) at (1, 2.5) {};
\node[hidden1] (H1-2) at (1, 1.5) {};
\node[hidden1] (H1-3) at (1, 0.5) {};
\node[hidden1] (H1-4) at (1, -0.5) {};
\node[lbl] at (1, 3.5) {Hidden 1};

% Segunda capa oculta (4 neuronas)
\node[hidden2] (H2-1) at (2, 2.5) {};
\node[hidden2] (H2-2) at (2, 1.5) {};
\node[hidden2] (H2-3) at (2, 0.5) {};
\node[hidden2] (H2-4) at (2, -0.5) {};
\node[lbl] at (2, 3.5) {Hidden 2};

% Tercera capa oculta (4 neuronas)
\node[hidden3] (H3-1) at (3, 2.5) {};
\node[hidden3] (H3-2) at (3, 1.5) {};
\node[hidden3] (H3-3) at (3, 0.5) {};
\node[hidden3] (H3-4) at (3, -0.5) {};
\node[lbl] at (3, 3.5) {Hidden 3};

% Capa de salida (2 neuronas)
\node[output] (O1) at (4, 1.5) {};
\node[output] (O2) at (4, 0.5) {};
\node[right=1mm of O1, font=\small] {$y_1$};
\node[right=1mm of O2, font=\small] {$y_2$};
\node[lbl] at (4, 3.5) {Output};

% Conexiones: entrada -> oculta 1
\foreach \i in {1,2,3}
    \foreach \j in {1,...,4}
        \draw[conn] (I\i) -- (H1-\j);

% Conexiones: oculta 1 -> oculta 2
\foreach \i in {1,...,4}
    \foreach \j in {1,...,4}
        \draw[conn] (H1-\i) -- (H2-\j);

% Conexiones: oculta 2 -> oculta 3
\foreach \i in {1,...,4}
    \foreach \j in {1,...,4}
        \draw[conn] (H2-\i) -- (H3-\j);

% Conexiones: oculta 3 -> salida
\foreach \i in {1,...,4}
    \foreach \j in {1,2}
        \draw[conn] (H3-\i) -- (O\j);

% Peso ejemplo (conexión destacada)
\draw[->, >=stealth, draw=gdlblue!60, line width=0.5pt] (I1) -- node[above, font=\scriptsize, sloped, pos=0.4] {$w_{11}^{(1)}$} (H1-1);

% Corchete para capas ocultas
\draw[decorate, decoration={brace, amplitude=4pt, raise=2pt}, gdlgray!60, thick]
    (1, 3.7) -- (3, 3.7) node[midway, above=5pt, font=\small\itshape] {Hidden layers};

\end{tikzpicture}
\caption[Structure of a layered neural network]{Structure of a feedforward neural network organized into layers. The nodes represent individual neurons with colors differentiated by type of layer: input (blue), hidden (green/teal/violet), and output (orange). The arrows represent weighted connections $w_{ij}^{(\ell)}$.}
\label{fig:neural_network}
\end{figure}

In an \gls{ANN} we have a first layer of the network called the input layer that receives the raw data, that is, the signal of the problem domain; and the last layer is the output layer that produces the final response of the network. The intermediate layers are known as hidden layers, since they are not directly seen at the input or the output of the neural network.

Let us look in more detail at why it is important to have more than one layer and more than one neuron in some layer, in the XOR problem.

\subsubsection{The XOR problem}

The XOR problem consists of computing the logical function:
\[
\text{XOR}(x_1, x_2) = \begin{cases}
1 & \text{if exactly one input is 1} \\
0 & \text{otherwise}
\end{cases}
\]
This problem is \textit{not linearly separable}: there is no single hyperplane that separates the classes $(0,0), (1,1)$ (output 0) from $(0,1), (1,0)$ (output 1) in $\mathbb{R}^2$.

\autoref{fig:xor_geometrico} geometrically illustrates how two neurons in the hidden layer, each defining a hyperplane, allow the XOR problem to be solved.

\begin{figure}[h]
\centering
\begin{tikzpicture}[scale=2.5]

% Ejes
\draw[->] (-0.15, 0) -- (1.25, 0) node[right, font=\small] {$x_1$};
\draw[->] (0, -0.15) -- (0, 1.25) node[above, font=\small] {$x_2$};

% Puntos clase 0 (azul) - XOR=0
\fill[gdlblue!70] (0, 0) circle (0.06) node[below left, font=\small] {$(0,0)$};
\fill[gdlblue!70] (1, 1) circle (0.06) node[above right, font=\small] {$(1,1)$};

% Puntos clase 1 (rojo) - XOR=1
\fill[gdlred!70] (0, 1) circle (0.06) node[above left, font=\small] {$(0,1)$};
\fill[gdlred!70] (1, 0) circle (0.06) node[below right, font=\small] {$(1,0)$};

% Hiperplano 1 (neurona 1): x1 + x2 = 0.5
\draw[gdlgreen!60!black, thick] (-0.1, 0.6) -- (0.6, -0.1);
\node[gdlgreen!60!black, font=\small, rotate=-45] at (0.05, 0.15) {$h_1$};

% Hiperplano 2 (neurona 2): x1 + x2 = 1.5
\draw[gdlgreen!70!black, thick] (0.4, 1.1) -- (1.1, 0.4);
\node[gdlgreen!70!black, font=\small, rotate=-45] at (0.95, 0.85) {$h_2$};

% Leyenda
\node[anchor=west, font=\small] at (1.4, 0.9) {$h_1$: neuron 1};
\node[anchor=west, font=\small] at (1.4, 0.7) {$h_2$: neuron 2};
\node[anchor=west, font=\small, gdlblue!70] at (1.4, 0.4) {XOR = 0};
\node[anchor=west, font=\small, gdlred!70] at (1.4, 0.2) {XOR = 1};

\end{tikzpicture}
\caption[Geometric solution of the XOR problem]{Geometric solution of the XOR problem through two neurons. Each neuron of the hidden layer defines a hyperplane ($h_1$ and $h_2$) that divides the space. The region between the two hyperplanes corresponds to XOR=1 (red points), while the outer regions correspond to XOR=0 (blue points).}
\label{fig:xor_geometrico}
\end{figure}

Which architectures can solve XOR?
\begin{itemize}
    \item \textbf{1 layer with 1 neuron}: \textit{Impossible}. A neuron computes $\sigma(w_1 x_1 + w_2 x_2 + b)$, which defines a single hyperplane. XOR requires a non-convex decision boundary.

    \item \textbf{Many layers with 1 neuron each}: Also \textit{impossible}. If each layer applies a scalar function $f_i: \mathbb{R} \to \mathbb{R}$, the composition $f_k \circ \cdots \circ f_1: \mathbb{R}^2 \to \mathbb{R}$ collapses all the two-dimensional information to a scalar immediately, losing the ability to distinguish the four regions of the input space.

    \item \textbf{2 layers with $\{2, 1\}$ neurons}: \textit{Sufficient}. The hidden layer with 2 neurons creates two hyperplanes that partition $\mathbb{R}^2$ into regions, and the output neuron linearly combines these partitions to generate the XOR boundary.
\end{itemize}

This counterexample reveals that \textit{neither width alone nor depth alone suffices}: we need $m \geq 2$ neurons in at least one layer \textit{and} at least $k \geq 2$ layers. Extreme architectures (deep-narrow or shallow-wide with $m=1$) fail on fundamental problems.


\subsection{Width vs depth}\label{subsec:anchura_profundidad}

The analysis of the XOR problem leads us to formalize two fundamental concepts that characterize the architecture of a neural network: \textit{width} and \textit{depth}. These parameters determine the expressive capacity of the network and have profound implications, both theoretical and practical.

\begin{definition}[Width of a layer]\label{def:anchura}
The \textit{width} of a layer $L_i: \mathbb{R}^{n_{i-1}} \to \mathbb{R}^{n_i}$ is the dimension $n_i$ of its output space, that is, the number of neurons it contains.
\end{definition}

\begin{definition}[Depth of a network]\label{def:profundidad}
The \textit{depth} of a neural network $\mathcal{N} = L_k \circ \cdots \circ L_1$ is the number $k$ of layers in the composition.
\end{definition}

These formal definitions allow us to rigorously analyze the balance between width and depth.

Width provides \textit{diversity of features}: each neuron captures a different aspect of the input signal, and the output signal $h \in \mathbb{R}^m$ simultaneously encodes multiple geometric properties of the datum.

Depth provides \textit{hierarchical abstraction}: instead of attempting to solve the complete problem at once, each layer performs an intermediate transformation, delegating the final refinement to the subsequent layers. The layers implement a hierarchy of representations where each level extracts features of increasing complexity:
\[
x \xrightarrow{L_1} h^{(1)} \xrightarrow{L_2} h^{(2)} \xrightarrow{L_3} \cdots \xrightarrow{L_k} y
\]
Empirically, it is observed that:
\begin{itemize}
    \item \textbf{Early layers} ($L_1, L_2$): Detect low-level features (edges, textures, basic colors in images; phonemes in audio).
    \item \textbf{Intermediate layers} ($L_3, \ldots, L_{k-2}$): Combine low-level features into more complex patterns (geometric shapes, textural motifs).
    \item \textbf{Final layers} ($L_{k-1}, L_k$): Represent high-level concepts specific to the task (objects, scenes; words, semantic concepts).
\end{itemize}
This hierarchy reflects the compositional structure of many real-world problems and allows the network to learn increasingly abstract and semantically meaningful representations.

As we will see in the Universal Approximation Theorem (\autoref{sec:aproximacion_universal}), neural networks with a single sufficiently wide hidden layer can approximate any continuous function on a compact set with arbitrary precision. However, this existential guarantee does not specify the efficiency of such an approximation in terms of the number of neurons required.

\paragraph{Approximation theory and depth.}

Recent theoretical results demonstrate that depth (number of layers) provides exponential advantages in expressive capacity. For certain classes of functions, a deep network of moderate width can approximate functions that would require an exponentially larger number of neurons if limited to a single layer~\citep{telgarsky2016benefitsdepthneuralnetworks}. Mathematically, this is because the composition of functions allows the construction of functions with combinatorial complexity that grows exponentially with depth.

We have seen that, although a single hidden layer suffices theoretically to approximate continuous functions, deep architectures offer exponential advantages in efficiency. In practice, deep networks with different types of specialized layers are preferred, each adapted to specific aspects of the problem, in order to optimize both the expressive capacity and the computational cost. This specialization reflects a fundamental distinction: while some layers transform only the values of the signal preserving its spatial structure, others modify the geometric organization itself of the domain. This functional taxonomy will allow us to better understand the difference between traditional deep learning architectures and the \gls{GDL} approach.

\subsection{Types of layers}\label{subsec:tipos_de_capas}

Neural layers can transform the signal, the domain, or both simultaneously. This distinction is fundamental to understand the philosophical and mathematical difference between traditional \glspl{ANN} and \gls{GDL}: while the former typically ignore or destroy the original geometry of the datum, the latter preserve and systematically exploit it.

Formally, we classify layers into three categories according to which dimension they transform:

\begin{definition}[Pure signal transformation]\label{def:signal_transform}
	A layer $L$ is a \textit{pure signal transformation} if it preserves the structure of the input domain. Formally, if the input is a signal $f: \Omega \to \mathbb{R}^{c_{\text{in}}}$, the output is $g: \Omega \to \mathbb{R}^{c_{\text{out}}}$ where the domain $\Omega$ remains invariant.

	\textbf{Characteristics:}
	\begin{itemize}
		\item Point-wise (element-wise) operations: $g(x) = \phi(f(x))$ for some function $\phi: \mathbb{R}^{c_{\text{in}}} \to \mathbb{R}^{c_{\text{out}}}$
		\item Preserve the shape of the input tensor
		\item Do not require information about neighborhoods, connectivity, or geometric structure
	\end{itemize}
\end{definition}

\begin{definition}[Pure geometric transformation]\label{def:geometry_transform}
	A layer $L$ is a \textit{pure geometric transformation} if it modifies the geometric structure of the domain without processing the values of the signal in a non-trivial way. Formally, if $f: (\Omega, \mathcal{G}) \to \mathbb{R}^c$ is a signal over a domain $\Omega$ with geometric structure $\mathcal{G}$ (metric, connection, topology), then $L(f) = g: (\Omega', \mathcal{G}') \to \mathbb{R}^c$ where the geometry $\mathcal{G}' \neq \mathcal{G}$, but the values of the signal are preserved or reindexed according to a correspondence $\pi: \Omega' \to \Omega$.
\end{definition}

In \textbf{traditional neural networks}, pure geometric transformations typically reduce to reindexings (reshape, transpose, flatten), since the underlying domain is a discretized Euclidean space without rich geometric structure. However, in \gls{GDL}, these transformations include more sophisticated operations: changes of coordinate system, projections between manifolds, modifications of the metric or connection, and conformal or isometric transformations. The distinction is fundamental: while in traditional \gls{ANN} the geometry is ``flat'' and purely structural transformations are trivial, in \gls{GDL} the geometry is intrinsic to the problem and geometric transformations can be highly non-trivial.

\begin{definition}[Coupled signal-geometry transformation]\label{def:coupled_transform}
	A layer $L$ is a \textit{coupled transformation} if it simultaneously and intrinsically modifies both the signal and the domain. Formally, $L: (f: \Omega \to \mathbb{R}^{c_{\text{in}}}) \mapsto (g: \Omega' \to \mathbb{R}^{c_{\text{out}}})$ where typically $\Omega' \neq \Omega$ and $g$ cannot be expressed as a reindexing of $f$.
\end{definition}

\textbf{Traditional neural networks} (multilayer perceptrons, dense networks) are based mainly on layers that \textit{destroy} or completely ignore the original geometry of the datum. It is very typical to flatten the geometry and transform it into a vector, losing all information about spatial neighborhoods and local structure, or not using the intrinsic properties of the geometry such as symmetries to improve neural networks. In contrast, as we will see in \autoref{ch:gdl}, \gls{GDL} preserves and systematically exploits the underlying geometry.

\subsection{Traditional layers}

By way of introduction, we will briefly present here some of the layers traditionally used in Deep Learning.

\subsubsection{Activation functions}\label{subsubsec:funciones_activacion_capas}

Activation functions introduce essential nonlinearities in \gls{ANN}, allowing them to approximate complex functions. Without these nonlinearities, any neural network would reduce to a simple linear transformation, drastically limiting its expressive capacity. Earlier we saw that the activation function is an intrinsic part of the neuron, but for practical reasons it is usually implemented as an independent layer that transforms the output of each neuron component-wise. We consider this type of layers among those that only perform transformations of the signal.

Below we present the most widely used activation functions, analyzing in each case the desirable properties established in \autoref{prop:acotacion_activacion} (see also \autoref{tab:funciones_activacion} for a summarized comparison):

\paragraph{Sigmoid}

The sigmoid function is one of the classical activation functions. Its range of values is $(0, 1)$, useful for probabilistic interpretations:
\[
\sigma(x) = \frac{1}{1 + e^{-x}}
\]

It is bounded (\autoref{prop:acotacion_activacion}), monotonic (\autoref{prop:monotonia_activacion}) and Lipschitz continuous with constant $L = 1/4$ (\autoref{prop:lipschitz_continuidad}). However, it is not zero-centered (\autoref{prop:zero_centered}), since $\sigma(z) \in (0,1)$ and $\sigma(0) = 1/2 \neq 0$, which can cause the zigzagging phenomenon in gradient descent. In addition, its horizontal asymptotes ($\lim_{z\to\pm\infty}\sigma(z) \in \{0,1\}$) imply gradient saturation at both extremes (\autoref{prop:gradiente_no_saturante}).

\paragraph{Hyperbolic Tangent (tanh)}

The tanh function has a range $(-1, 1)$:
\[
\tanh(x) = \frac{e^x - e^{-x}}{e^x + e^{-x}} = \frac{e^{2x} - 1}{e^{2x} + 1}
\]

Unlike the sigmoid, the hyperbolic tangent is zero-centered (\autoref{prop:zero_centered}): it is an odd function ($\tanh(-z) = -\tanh(z)$) and $\tanh(0) = 0$. It shares with the sigmoid boundedness (\autoref{prop:acotacion_activacion}), monotonicity (\autoref{prop:monotonia_activacion}) and Lipschitz continuity with $L = 1$ (\autoref{prop:lipschitz_continuidad}).

However, both the sigmoid and the hyperbolic tangent are bounded functions with horizontal asymptotes, which makes them saturating (\autoref{prop:gradiente_no_saturante}). As analyzed in \autoref{subsubsec:backpropagation}, this causes the \textit{vanishing gradient} problem: in deep networks, the gradient is attenuated exponentially in the first layers, which hardly ``learn''.

\paragraph{ReLU (Rectified Linear Unit)}

The \gls{ReLU} function is one of the most widely used activation functions in intermediate layers:
\[
\text{ReLU}(x) = \max(0, x) = \begin{cases}
x & \text{if } x > 0 \\
0 & \text{if } x \leq 0
\end{cases}
\]

\gls{ReLU} solves the vanishing gradient problem: for $z > 0$, $\text{ReLU}'(z) = 1$, satisfying \autoref{prop:gradiente_no_saturante} on the positive half-line. It has linear growth (\autoref{prop:acotacion_activacion}, condition 2), is monotonic (\autoref{prop:monotonia_activacion}) and Lipschitz continuous with $L = 1$ (\autoref{prop:lipschitz_continuidad}). In addition, since it requires only a comparison and a selection, it is computationally very efficient (\autoref{prop:eficiencia_computacional}). Nevertheless, it is not zero-centered (\autoref{prop:zero_centered}), and its zero gradient for $z \leq 0$ can cause the ``dead neurons'' problem: neurons whose input remains negative stop contributing to learning permanently.

\paragraph{LeakyReLU}

A variant of \gls{ReLU} that allows a small gradient for negative values:
\[
\text{LeakyReLU}(x) = \begin{cases}
x & \text{if } x > 0 \\
\alpha x & \text{if } x \leq 0
\end{cases}
\]

where $\alpha > 0$ is a small parameter (typically $\alpha = 0.01$). LeakyReLU fully satisfies \autoref{prop:gradiente_no_saturante}: $|\text{LeakyReLU}'(z)| \geq \alpha > 0$ for all $z \neq 0$, eliminating both the saturation and the dead neurons problem of ReLU. It is monotonic (\autoref{prop:monotonia_activacion}), has linear growth (\autoref{prop:acotacion_activacion}), is Lipschitz with $L = 1$ (\autoref{prop:lipschitz_continuidad}) and shares the computational efficiency of ReLU (\autoref{prop:eficiencia_computacional}).

\paragraph{Softmax}

This function is normally used to obtain probabilistic values in the last layer of neural networks, and is fundamental for the attention mechanisms that we will study in \autoref{ch:attention}.

\begin{definition}[Softmax function]\label{def:softmax}
The \textit{softmax function} $\sigma: \mathbb{R}^n \to (0,1)^n$ transforms a vector of real values into a probability distribution. For a vector $\mathbf{z} = (z_1, \ldots, z_n) \in \mathbb{R}^n$:
\begin{equation}\label{eq:softmax}
\sigma(\mathbf{z})_i = \frac{e^{z_i}}{\sum_{j=1}^{n} e^{z_j}}, \quad i = 1, \ldots, n
\end{equation}
\end{definition}

\begin{proposition}[Properties of the Softmax function]\label{prop:softmax_propiedades}
The softmax function satisfies the following properties:
\begin{enumerate}
	\item \textbf{Normalization}: $\sum_{i=1}^n \sigma(\mathbf{z})_i = 1$ for all $\mathbf{z} \in \mathbb{R}^n$.
	\item \textbf{Positivity}: $\sigma(\mathbf{z})_i > 0$ for all $i$ and all $\mathbf{z} \in \mathbb{R}^n$.
	\item \textbf{Translation invariance}: $\sigma(\mathbf{z} + c\mathbf{1}) = \sigma(\mathbf{z})$ for all $c \in \mathbb{R}$.
	\item \textbf{Differentiability}: The Jacobian has the form $\frac{\partial \sigma_i}{\partial z_j} = \sigma_i(\delta_{ij} - \sigma_j)$.
	\item \textbf{Permutation equivariance}: For any permutation matrix $P$, $\sigma(P\mathbf{z}) = P\sigma(\mathbf{z})$.
\end{enumerate}
\end{proposition}

\begin{proof}
Properties (1) and (2) are verified directly from the definition. Property (3) results from the fact that $e^{z_i + c}/\sum_j e^{z_j + c} = e^{z_i}e^c/(e^c\sum_j e^{z_j}) = e^{z_i}/\sum_j e^{z_j}$. For (4), applying the quotient rule:
\[
\frac{\partial \sigma_i}{\partial z_j} = \frac{\delta_{ij}e^{z_i}\sum_k e^{z_k} - e^{z_i}e^{z_j}}{(\sum_k e^{z_k})^2} = \sigma_i\delta_{ij} - \sigma_i\sigma_j = \sigma_i(\delta_{ij} - \sigma_j)
\]
Property (5) follows from the fact that the sum in the denominator is invariant under permutation of indices.
\end{proof}

\subsubsection{Normalization}

Normalization techniques have revolutionized the training of deep neural networks by addressing fundamental problems such as internal covariate shift and gradient instability. These techniques normalize the distributions of activations at different points of the network. It is a type of pure signal transformation layer; in particular, normalization applies statistics over the signal (mean, variance) without modifying the structure of the domain.

\begin{definition}[Normalization in neural networks]\label{def:normalizacion}
	A normalization technique is a transformation that standardizes the activations of a layer by applying a function of the form:
	\[
	\text{Norm}(\mathbf{x}) = \gamma \odot \frac{\mathbf{x} - \boldsymbol{\mu}}{\boldsymbol{\sigma}} + \boldsymbol{\beta}
	\]
	where $\boldsymbol{\mu}$ and $\boldsymbol{\sigma}$ are centering and scaling statistics, and $\gamma$, $\boldsymbol{\beta}$ are learnable parameters that allow the identity transformation to be recovered if necessary\footnote{In practice, $\sqrt{\boldsymbol{\sigma}^2 + \epsilon}$ with $\epsilon \approx 10^{-8}$ is used instead of $\boldsymbol{\sigma}$ to avoid divisions by zero, as observed in the specializations that follow.}.
\end{definition}

There are different normalization layers that are used depending on how one wants to group the neurons of a layer.

\paragraph{\gls{BN}}

Introduced by Ioffe and Szegedy~\citep{ioffe2015batch}, batch normalization computes statistics over the mini-batch:

\[
\text{BN}(\mathbf{x}_i) = \gamma \frac{\mathbf{x}_i - \boldsymbol{\mu}_{\mathcal{B}}}{\sqrt{\boldsymbol{\sigma}^2_{\mathcal{B}} + \epsilon}} + \boldsymbol{\beta}
\]

where for a mini-batch $\mathcal{B} = \{\mathbf{x}_1, \ldots, \mathbf{x}_m\}$:
\begin{align}
\boldsymbol{\mu}_{\mathcal{B}} &= \frac{1}{m} \sum_{i=1}^m \mathbf{x}_i \\
\boldsymbol{\sigma}^2_{\mathcal{B}} &= \frac{1}{m} \sum_{i=1}^m (\mathbf{x}_i - \boldsymbol{\mu}_{\mathcal{B}})^2
\end{align}

\paragraph{Layer Normalization}

Proposed by Ba, Kiros and Hinton~\citep{ba2016layer}, it normalizes across the features of each individual sample:

\[
\text{LN}(\mathbf{x}) = \gamma \frac{\mathbf{x} - \boldsymbol{\mu}_L}{\sqrt{\boldsymbol{\sigma}^2_L + \epsilon}} + \boldsymbol{\beta}
\]

where for a sample $\mathbf{x} \in \mathbb{R}^d$:
\begin{align}
\mu_L &= \frac{1}{d} \sum_{i=1}^d x_i \\
\sigma^2_L &= \frac{1}{d} \sum_{i=1}^d (x_i - \mu_L)^2
\end{align}

\paragraph{Group Normalization}

Introduced by Wu and He~\citep{wu2018group}, it divides the channels into groups and normalizes within each group:

For an input tensor with shape $(N, C, H, W)$, it divides $C$ channels into $G$ groups of $C/G$ channels each, and applies normalization within each group.

\[
\text{GN}(\mathbf{x}) = \gamma \frac{\mathbf{x} - \boldsymbol{\mu}_G}{\sqrt{\boldsymbol{\sigma}^2_G + \epsilon}} + \boldsymbol{\beta}
\]

\paragraph{Instance Normalization}

It normalizes each channel of each sample independently, common in style transfer tasks~\citep{ulyanov2016instance}:

\[
\text{IN}(\mathbf{x}) = \gamma \frac{\mathbf{x} - \boldsymbol{\mu}_{IN}}{\sqrt{\boldsymbol{\sigma}^2_{IN} + \epsilon}} + \boldsymbol{\beta}
\]

where the statistics are computed per channel and per sample.

\subsubsection{Dropout}

Dropout~\citep{srivastava2014dropout}, is a fundamental regularization technique that addresses the problem of overfitting in deep neural networks in an elegant and computationally efficient way. It is a type of layer that we can consider as pure signal, since it only applies a stochastic point-wise masking without altering the geometry of the domain.

\begin{definition}[Dropout]\label{def:dropout}
	Dropout is a regularization technique that during training randomly deactivates neurons with probability $p_{\text{drop}}$, setting their outputs to zero. Formally, for a layer with input $\mathbf{x}$:
	\[
	\mathbf{y} = \frac{1}{1-p_{\text{drop}}} \mathbf{r} \odot \mathbf{x}
	\]
	where $\mathbf{r} \sim \text{Bernoulli}(1-p_{\text{drop}})$ is a vector of random binary masks, and $\odot$ denotes the element-wise product.
\end{definition}

Dropout provides several regularization benefits:
\begin{itemize}
	\item \textbf{Reduces co-adaptation}: Prevents neurons from developing complex dependencies among themselves
	\item \textbf{Improves generalization}: Forces the network not to depend excessively on specific neurons
	\item \textbf{Robustness}: Makes the representations less fragile to perturbations
\end{itemize}

\subsubsection{Dense layers}\label{def:capa_densa}

A \textit{dense layer} (also called fully connected layer) is precisely a neural layer according to \autoref{def:capa_neuronal}, without any structural restriction on the weight matrix: each output neuron receives a signal from all the input positions, that is, $W$ can have all of its coefficients nonzero. It is the most basic architecture and constitutes the basis on which more specialized architectures are built.

From a geometric perspective, a dense layer $L(\mathbf{x}) = \sigma(W\mathbf{x} + \mathbf{b})$ performs three successive operations:
\begin{enumerate}
	\item \textbf{Linear transformation}: $W\mathbf{x}$ represents a change of basis and scaling in $\mathbb{R}^n \to \mathbb{R}^m$
	\item \textbf{Translation}: $+ \mathbf{b}$ displaces the result in $\mathbb{R}^m$
	\item \textbf{Nonlinear deformation}: $\sigma(\cdot)$ introduces curvature in the feature space
\end{enumerate}

By not imposing sparse connectivity restrictions or shared weights, the dense layer completely destroys the original geometry of the input domain through an affine projection to a space of arbitrary dimension. This is a coupled signal-geometry transformation (\autoref{def:coupled_transform}).

\subsubsection{Embedding}

Embedding layers constitute a fundamental component for the processing of categorical and discrete data in neural networks. They allow the transformation of high-dimensional and sparse categorical variables into dense lower-dimensional representations that capture semantic relationships. Embeddings map discrete tokens to learned continuous spaces, fundamentally changing both the domain (discrete $\to$ continuous) and the representation of the signal, that is, they are a type of layer that transforms the signal and the geometry in a coupled way.

\begin{definition}[Embedding layer]\label{def:capa_embedding}
	An embedding layer is a transformation that maps discrete indices to dense vectors. Formally, for a vocabulary of size $V_{\text{vocab}}$ and embedding dimension $d_{\text{emb}}$:
	\[
	\text{Embedding}: \{1, 2, \ldots, V_{\text{vocab}}\} \to \mathbb{R}^{d_{\text{emb}}}
	\]
	implemented through a lookup matrix $E_{\text{emb}} \in \mathbb{R}^{V_{\text{vocab}} \times d_{\text{emb}}}$, where each row $(E_{\text{emb}})_i$ is the embedding vector for token $i$.
\end{definition}

The most popular application of embeddings is to transform text into vectors for its subsequent processing in the neural network.

Embeddings represent a fundamental case of geometric representation learning: they transform a discrete and disconnected space (categories) into a continuous manifold where Euclidean geometry captures semantic relationships. This transformation is essential for neural networks to effectively process symbolic and categorical data.

\subsubsection{Structure transformations}

Structure transformations modify the geometry and organization of the data without fundamentally changing their informational content. These operations are crucial for adapting the data to different network architectures, exploiting domain-specific patterns, and optimizing computation.

\paragraph{Reshape}

\begin{definition}[Reshape operations]\label{def:reshape}
	A reshape operation is a transformation that reorganizes a tensor by changing its dimensions while preserving the total number of elements:
	\[
	\text{reshape}: \mathbb{R}^{d_1 \times d_2 \times \cdots \times d_n} \to \mathbb{R}^{d_1' \times d_2' \times \cdots \times d_m'}
	\]
	where $\prod_{i=1}^n d_i = \prod_{j=1}^m d_j'$.
\end{definition}

\paragraph{Flatten} is a special case of reshape that converts multidimensional tensors into one-dimensional vectors:
\[
\text{flatten}: \mathbb{R}^{d_1 \times d_2 \times \cdots \times d_n} \to \mathbb{R}^{\prod_{i=1}^n d_i}
\]

A very common application of these layers is to serve as a transition between convolutional layers and dense layers.

\paragraph{Pooling}

Pooling reduces the spatial dimensionality by combining local information, providing limited translational invariance, reducing noise and computational cost.

\begin{definition}[Pooling operation]\label{def:pooling}
	A pooling operation aggregates values in local windows using an aggregation function $f$:
	\[
	y_{i,j} = f\left(\{x_{si+k, sj+l} : k,l \in W_{\text{pool}}\}\right)
	\]
	where $s$ is the stride and $W_{\text{pool}}$ defines the size of the pooling window.
\end{definition}

\textbf{Max Pooling}: Selects the maximum value in each window:
\[
\text{MaxPool}(X) = \max_{(i,j) \in W_{\text{pool}}} X_{i,j}
\]
\textbf{Average Pooling}: Computes the arithmetic mean in each window:
\[
\text{AvgPool}(X) = \frac{1}{|W_{\text{pool}}|} \sum_{(i,j) \in W_{\text{pool}}} X_{i,j}
\]

These transformations are fundamental for building hierarchical architectures that process information at multiple scales and resolutions.

\subsubsection{Convolutions}

Convolutional layers are explained in greater detail in \autoref{ch:cnn}. They are used fundamentally for grid data such as images or sequences such as audio. In traditional \glspl{ANN}, convolutions are introduced mainly as a computational optimization that takes advantage of the regular grid structure of the domain. As we will see in \autoref{ch:cnn}, there exist more sophisticated variants of convolutional networks that exploit more deeply the geometric properties and symmetries of the domain.

\begin{definition}[Discrete convolution]\label{def:convolucion_discreta}
	For an input function $f: \mathbb{Z}^d \to \mathbb{R}$ and a kernel $g: \mathbb{Z}^d \to \mathbb{R}$, the discrete convolution is defined as:
	\[
	(f * g)[n] = \sum_{m \in \mathbb{Z}^d} f[m] g[n - m]
	\]
\end{definition}

\subsubsection{Recurrent}\label{subsubsec:rnn}

Recurrent neural networks (\glspl{RNN}) extend feedforward architectures by incorporating cyclic connections that allow the processing of sequential data. This ability to maintain \textit{memory} across time makes them fundamental for tasks involving temporal sequences. \glspl{RNN} process temporal geometry (sequential chains) while aggregating stateful information (memory) that progressively transforms the signal.

\begin{definition}[Recurrent neural network]\label{def:rnn}
	An \textit{\gls{RNN}} is a neural network with at least one cycle in its connectivity graph. For the processing of sequences, the hidden state evolves according to:
	\[
	\mathbf{h}_t = f(\mathbf{h}_{t-1}, \mathbf{x}_t; \boldsymbol{\theta})
	\]
	where $\mathbf{h}_t$ is the hidden state at time $t$, $\mathbf{x}_t$ is the current input, and $\boldsymbol{\theta}$ are the parameters shared across time.
\end{definition}

The simplest form of an \gls{RNN} uses:
\begin{gather*}
    \mathbf{h}_t = \tanh(W_h \mathbf{h}_{t-1} + W_x \mathbf{x}_t + \mathbf{b}_h)\\
    \mathbf{y}_t = W_y \mathbf{h}_t + \mathbf{b}_y\\
\end{gather*}

where:
\begin{itemize}
	\item $W_h \in \mathbb{R}^{d_h \times d_h}$ are the recurrent weights
	\item $W_x \in \mathbb{R}^{d_h \times d_x}$ are the input weights
	\item $W_y \in \mathbb{R}^{d_y \times d_h}$ are the output weights
	\item $\mathbf{b}_h, \mathbf{b}_y$ are the bias vectors
\end{itemize}

A key principle of \gls{RNN} is that the same parameters $\boldsymbol{\theta} = \{W_h, W_x, W_y, \mathbf{b}_h, \mathbf{b}_y\}$ are applied at each time step. This allows:
\begin{itemize}
	\item \textbf{Temporal generalization}: The same model works for sequences of any length
	\item \textbf{Parametric efficiency}: $O(d_h^2 + d_h d_x)$ parameters independent of the sequence length
	\item \textbf{Temporal inductive bias}: It assumes that the same transition rules apply at all times
\end{itemize}

\autoref{fig:rnn_unrolled} shows how an RNN with a cycle (compact) ``unrolls'' over time to form a \gls{DAG} that allows backpropagation to be applied.

\begin{figure}[h]
\centering
\begin{tikzpicture}[
    cell/.style={rectangle, draw, rounded corners, minimum width=1.2cm, minimum height=0.8cm, fill=gdlblue!15},
    input/.style={circle, draw, minimum size=0.5cm, fill=gdlgreen!20},
    output/.style={circle, draw, minimum size=0.5cm, fill=gdlred!20},
    hidden/.style={circle, draw, minimum size=0.5cm, fill=gdlyellow!20},
    arrow/.style={->, thick, >=stealth}
]

% --- RNN compacta (izquierda) ---
\node[font=\bfseries\small] at (-4.0, 3.0) {Compact RNN};
\node[font=\scriptsize] at (-4.0, 2.6) {(cyclic graph)};

\node[cell] (rnn) at (-4.0, 0.5) {RNN};
\node[input] (xc) at (-4.0, -0.8) {$\mathbf{x}$};
\node[output] (yc) at (-4.0, 1.8) {$\mathbf{y}$};

\draw[arrow] (xc) -- (rnn);
\draw[arrow] (rnn) -- (yc);
% Ciclo de recurrencia: bucle claro a la derecha
\draw[arrow] (rnn.east) to[out=40, in=-40, distance=1.1cm] node[pos=0.25, above, font=\scriptsize] {$\mathbf{h}$} (rnn.east);

% --- Flecha de transformación ---
\draw[->, very thick, dashed] (-2.3, 0.5) -- (-1.3, 0.5) node[midway, above, font=\scriptsize] {unroll};

% --- RNN desplegada (derecha) ---
\node[font=\bfseries\small] at (3.5, 3.0) {RNN unrolled over $T$ steps};
\node[font=\scriptsize] at (3.5, 2.6) {(DAG for BPTT)};

% Paso t-1
\node[cell] (r1) at (1.0, 0.5) {RNN};
\node[input] (x1) at (1.0, -0.8) {$\mathbf{x}_{t-1}$};
\node[output] (y1) at (1.0, 1.8) {$\mathbf{y}_{t-1}$};
\node[hidden] (h0) at (-0.5, 0.5) {};
\node[font=\scriptsize] at (-0.5, 0.1) {$\mathbf{h}_{t-2}$};

\draw[arrow] (x1) -- (r1);
\draw[arrow] (r1) -- (y1);
\draw[arrow] (h0) -- (r1);

% Paso t
\node[cell] (r2) at (3.5, 0.5) {RNN};
\node[input] (x2) at (3.5, -0.8) {$\mathbf{x}_t$};
\node[output] (y2) at (3.5, 1.8) {$\mathbf{y}_t$};

\draw[arrow] (x2) -- (r2);
\draw[arrow] (r2) -- (y2);
\draw[arrow] (r1) -- node[above, font=\scriptsize] {$\mathbf{h}_{t-1}$} (r2);

% Paso t+1
\node[cell] (r3) at (6, 0.5) {RNN};
\node[input] (x3) at (6, -0.8) {$\mathbf{x}_{t+1}$};
\node[output] (y3) at (6, 1.8) {$\mathbf{y}_{t+1}$};
\node[hidden] (hf) at (8, 0.5) {};

\draw[arrow] (x3) -- (r3);
\draw[arrow] (r3) -- (y3);
\draw[arrow] (r2) -- node[above, font=\scriptsize] {$\mathbf{h}_t$} (r3);
\draw[arrow] (r3) -- node[above, font=\scriptsize] {$\mathbf{h}_{t+1}$} (hf);

% Nota sobre parámetros compartidos
\draw[decorate, decoration={brace, amplitude=5pt, mirror}] (0.5, -1.5) -- (6.5, -1.5);
\node[font=\scriptsize, align=center] at (3.5, -2.0) {Same parameters $\boldsymbol{\theta}$ at all steps};

\end{tikzpicture}
\caption[Compact vs time-unrolled RNN]{Left: Compact representation of an RNN as a cyclic graph. Right: The same RNN ``unrolled'' over time, showing how the hidden state $\mathbf{h}_t$ propagates between time steps. This representation as a DAG allows Backpropagation Through Time (BPTT) to be applied.}
\label{fig:rnn_unrolled}
\end{figure}

The \gls{RNN} like the one we have just described, which only has a recurrent connection without any extra complexity, suffer from several limitations:
\begin{itemize}
	\item \textbf{Limited memory}: Difficulty in capturing very long-term dependencies
	\item \textbf{Recency bias}: Tendency to give more weight to recent information
	\item \textbf{Limited parallelization}: Sequential processing prevents temporal parallelization
	\item \textbf{Activation saturation}: Sigmoidal activation functions cause saturation
\end{itemize}

In addition, this type of networks require a variation of the learning algorithm that is normally used to train neural networks (see~\ref{sec:learning_ann}): \gls{BPTT}~\citep{werbos1990bptt}

\subsubsection{LSTMs}\label{subsubsec:lstm}

Long Short-Term Memory networks (\gls{LSTM}), introduced by Hochreiter and Schmidhuber~\citep{hochreiter1997long}, represent a crucial evolution of \gls{RNN} that specifically addresses the vanishing gradient problem and allows the learning of long-term dependencies. \glspl{LSTM} are more sophisticated variants of \glspl{RNN} that process temporal geometry through gating mechanisms that control the flow of information. The original architecture consisted of two gates (input and output); the forget gate was added later by \citet{gers2000learning}.

\begin{definition}[LSTM Cell]\label{def:lstm_cell}
	An \textit{\gls{LSTM} cell} maintains two states: a cell state $\mathbf{C}_t$ and a hidden state $\mathbf{h}_t$, controlled by three gates:
	\begin{itemize}
		\item \textbf{Forget gate}: $\mathbf{f}_t = \sigma(W_f \cdot [\mathbf{h}_{t-1}, \mathbf{x}_t] + \mathbf{b}_f)$
		\item \textbf{Input gate}: $\mathbf{i}_t = \sigma(W_i \cdot [\mathbf{h}_{t-1}, \mathbf{x}_t] + \mathbf{b}_i)$
		\item \textbf{Output gate}: $\mathbf{o}_t = \sigma(W_o \cdot [\mathbf{h}_{t-1}, \mathbf{x}_t] + \mathbf{b}_o)$
	\end{itemize}
	The cell state $\mathbf{C}_t$ is updated through:
	\begin{gather*}
		\tilde{\mathbf{C}}_t = \tanh(W_C \cdot [\mathbf{h}_{t-1}, \mathbf{x}_t] + \mathbf{b}_C)\\
		\mathbf{C}_t = \mathbf{f}_t \odot \mathbf{C}_{t-1} + \mathbf{i}_t \odot \tilde{\mathbf{C}}_t\\
		\mathbf{h}_t = \mathbf{o}_t \odot \tanh(\mathbf{C}_t)
	\end{gather*}
	where $\odot$ denotes the element-wise product.
\end{definition}

Each gate has a specific function:

\begin{itemize}
	\item \textbf{Forget gate} ($\mathbf{f}_t$): Decides what information to discard from the previous cell state
	\begin{itemize}
		\item $f_{t,i} \approx 0$: Completely forgets the information of component $i$
		\item $f_{t,i} \approx 1$: Completely retains the information of component $i$
	\end{itemize}

	\item \textbf{Input gate} ($\mathbf{i}_t$): Controls what new information to store in the cell state
	\begin{itemize}
		\item Works together with $\tilde{\mathbf{C}}_t$ to determine which values to update
		\item Allows selective writing into the memory
	\end{itemize}

	\item \textbf{Output gate} ($\mathbf{o}_t$): Controls which parts of the cell state to expose as output
	\begin{itemize}
		\item Filters the internal state to produce the hidden state
		\item Allows selective reading of the memory
	\end{itemize}
\end{itemize}

The key is that the forget gate can learn to stay close to 1 for important information, creating \textit{gradient paths} that persist across long temporal periods.

\autoref{fig:lstm_arquitectura} illustrates the internal architecture of an LSTM cell, showing the flow of information through the three gates and the cell state.

\begin{figure}[htbp]
\centering
\begin{tikzpicture}[
    % Estilos mejorados para mejor legibilidad
    gate/.style={
        rectangle, draw, rounded corners=4pt,
        minimum width=1.6cm, minimum height=1cm,
        fill=#1!25, font=\small\bfseries,
        line width=0.6pt
    },
    op/.style={
        rectangle, draw, rounded corners=3pt,
        minimum width=1.3cm, minimum height=0.9cm,
        fill=gdlyellow!30, font=\normalsize\bfseries,
        line width=0.6pt
    },
    state/.style={
        rectangle, draw, rounded corners=3pt,
        minimum width=1.4cm, minimum height=0.8cm,
        fill=gdlgray!15, font=\small,
        line width=0.6pt
    },
    conn/.style={->, >=stealth, draw=gdlgray!65, line width=0.7pt},
    main/.style={->, >=stealth, line width=0.8pt}
]

% ===== ESTADO CELULAR (línea superior) =====
\draw[->, >=stealth, line width=1.5pt, gdlblue!60] (-1.5,5) -- (13,5);
\node[above, font=\small, gdlblue!70] at (5.25,5.5) {Cell state $\mathbf{C}_t$};
\node[left, font=\small] at (-1.5,5) {$\mathbf{C}_{t-1}$};

% ===== REGIÓN 1: FORGET GATE =====
\node[gate=gdlred] (fg) at (1.5,2) {$\sigma$};
\node[below=0.15cm of fg, font=\scriptsize, gdlred!60!black] {Forget Gate};
\node[op] (fgmult) at (1.5,5) {$\odot$};

% ===== REGIÓN 2: INPUT GATE + CANDIDATO =====
\node[gate=gdlgreen] (ig) at (4.5,2) {$\sigma$};
\node[below=0.15cm of ig, font=\scriptsize, gdlgreen!40!black] {Input Gate};

\node[gate=gdlblue] (cand) at (7,2) {tanh};
\node[below=0.15cm of cand, font=\scriptsize, gdlblue!50!black] {Candidate $\tilde{\mathbf{C}}_t$};

\node[op] (igmult) at (5.75,3.5) {$\odot$};
\node[op] (add) at (5.75,5) {$\oplus$};

% ===== REGIÓN 3: OUTPUT GATE =====
\node[gate=gdlpurple] (og) at (9.5,2) {$\sigma$};
\node[below=0.15cm of og, font=\scriptsize, gdlpurple!60!black] {Output Gate};

\node[gate=gdlblue] (tanho) at (9.5,4) {tanh};
\node[op] (ogmult) at (11,3) {$\odot$};

% ===== ENTRADAS =====
\node[state] (xt) at (5.25,-0.8) {$\mathbf{x}_t$};
\node[state] (ht1) at (-0.8,0.5) {$\mathbf{h}_{t-1}$};

% ===== SALIDA =====
\node[state, fill=gdlorange!15] (ht) at (14,3) {$\mathbf{h}_t$};

% ===== CONEXIONES DE ENTRADA (más separadas y tenues) =====
% Concatenación visual [h_{t-1}, x_t]
\coordinate (concat) at (5.25, 0.5);

% h_{t-1} hacia punto de concatenación
\draw[conn] (ht1) -- (concat);

% x_t hacia punto de concatenación
\draw[conn] (xt) -- (concat);

% Desde concatenación hacia cada compuerta (rutas separadas)
\draw[conn] (concat) -- (1.5, 0.5) -- (fg);
\draw[conn] (concat) -- (4.5, 0.5) -- (ig);
\draw[conn] (concat) -- (7, 0.5) -- (cand);
\draw[conn] (concat) -- (9.5, 0.5) -- (og);

% Etiqueta de concatenación
\node[font=\scriptsize] at (5.25, 0.15) {$[\mathbf{h}_{t-1}, \mathbf{x}_t]$};

% ===== FLUJO INTERNO =====
% Forget gate → multiplicación con estado
\draw[main] (fg) -- (fgmult);

% Input gate y candidato → multiplicación
\draw[main, gdlgreen!60!black] (ig) |- (igmult);
\draw[main, gdlgreen!60!black] (cand) |- (igmult);

% Multiplicación input → suma con estado
\draw[main] (igmult) -- (add);

% Estado anterior pasa por forget
\draw[main, gdlblue!50] (fgmult) -- (add);

% Nuevo estado hacia output tanh
\draw[main, gdlblue!50] (8.2,5) -- (8.2,4) -- (tanho.west);

% Output gate y tanh → multiplicación final
\draw[main, gdlpurple!60!black] (og) -| (ogmult);
\draw[main, gdlgreen!50!black] (tanho) -| (ogmult);

% Multiplicación final → salida h_t
\draw[main] (ogmult) -- (ht);

% ===== RECURRENCIA (hacia siguiente paso) =====
\draw[conn, dashed, line width=0.8pt] (ht.north) -- (14,6.3) -- (-0.8,6.3) -- (-0.8,0.9)
    node[pos=0.5, above, font=\scriptsize] {$\mathbf{h}_t$ recurrent};

% ===== SEPARADORES DE REGIONES (opcional, muy sutil) =====
\draw[gdlgray!20, dashed, line width=0.3pt] (3,-.5) -- (3,5.5);
\draw[gdlgray!20, dashed, line width=0.3pt] (8.25,-.5) -- (8.25,5.5);

\end{tikzpicture}
\caption[LSTM cell architecture]{Internal architecture of an LSTM cell. The cell state $\mathbf{C}_t$ (top blue line) flows horizontally, being modulated by the three gates: the \textit{forget gate} (red) decides what information to discard, the \textit{input gate} (green) together with the candidate $\tilde{\mathbf{C}}_t$ determine what new information to store, and the \textit{output gate} (purple) controls which part of the cell state to expose as output $\mathbf{h}_t$. The operators $\odot$ represent element-wise multiplication and $\oplus$ the sum.}
\label{fig:lstm_arquitectura}
\end{figure}

\textbf{Gated Recurrent Unit (GRU)}: Proposed by Cho et al.~\citep{cho2014learning}, it is a simplification of the \gls{LSTM} that combines forget and input gates into a single update gate. Empirical studies have shown that the performance of \gls{GRU} is comparable to that of \gls{LSTM} on many tasks~\citep{chung2014empirical}:
\begin{align}
\mathbf{z}_t &= \sigma(W_z \cdot [\mathbf{h}_{t-1}, \mathbf{x}_t] + \mathbf{b}_z) \\
\mathbf{r}_t &= \sigma(W_r \cdot [\mathbf{h}_{t-1}, \mathbf{x}_t] + \mathbf{b}_r) \\
\tilde{\mathbf{h}}_t &= \tanh(W_h \cdot [\mathbf{r}_t \odot \mathbf{h}_{t-1}, \mathbf{x}_t] + \mathbf{b}_h) \\
\mathbf{h}_t &= (1 - \mathbf{z}_t) \odot \mathbf{h}_{t-1} + \mathbf{z}_t \odot \tilde{\mathbf{h}}_t
\end{align}

\gls{LSTM} have been fundamental for advances in machine translation, speech recognition, and language modeling, establishing the foundations for more modern architectures such as \glspl{Transformer} that address the limitation of temporal parallelization.

\subsubsection{Residual}

Residual connections, popularized by He et al.~\citep{he2016deep} in ResNets, represent a fundamental architectural advance that enabled the effective training of extremely deep neural networks. Their impact transcends convolutional networks, influencing practically all modern deep learning architectures. They are not a layer in the strict sense, but a way of connecting existing layers.

\begin{definition}[Residual connection]\label{def:conexion_residual}
	A residual connection (or skip connection) is a type of architecture where the output of a layer is added directly to the output of a later layer. For a block of layers $\mathcal{F}$:
	\[
	\mathbf{y} = \mathcal{F}(\mathbf{x}) + \mathbf{x}
	\]
	where $\mathbf{x}$ is the input of the block and $\mathcal{F}(\mathbf{x})$ represents the learned transformation.
\end{definition}

\autoref{fig:skip_connection} compares a traditional block with a residual block, showing how the direct connection allows the gradient to flow without attenuation.

\begin{figure}[h]
\centering
\begin{tikzpicture}[
    node distance=1cm,
    block/.style={rectangle, draw, minimum width=1.8cm, minimum height=0.7cm, fill=gdlblue!15},
    sum/.style={circle, draw, minimum size=0.5cm, fill=gdlyellow!20},
    arrow/.style={->, thick, >=stealth},
    skip/.style={->, thick, >=stealth, dashed, color=gdlred!70}
]

% Bloque tradicional (izquierda)
\node[font=\bfseries] at (-2, 2.5) {Without residual};
\node[block] (f1) at (-2, 1) {$\mathcal{F}_1$};
\node[block] (f2) at (-2, -0.5) {$\mathcal{F}_2$};
\draw[arrow] (-2, 2) -- (f1);
\draw[arrow] (f1) -- (f2);
\draw[arrow] (f2) -- (-2, -1.5);
\node[font=\scriptsize] at (-2, 2.2) {$\mathbf{x}$};
\node[font=\scriptsize] at (-2, -1.7) {$\mathbf{y}$};

% Bloque residual (derecha)
\node[font=\bfseries] at (2, 2.5) {With residual};
\node[block] (r1) at (2, 1) {$\mathcal{F}_1$};
\node[block] (r2) at (2, -0.5) {$\mathcal{F}_2$};
\node[sum] (plus) at (2, -1.5) {$+$};

\draw[arrow] (2, 2) -- (r1);
\draw[arrow] (r1) -- (r2);
\draw[arrow] (r2) -- (plus);
\draw[arrow] (plus) -- (2, -2.3);
\node[font=\scriptsize] at (2, 2.2) {$\mathbf{x}$};
\node[font=\scriptsize] at (2, -2.5) {$\mathbf{y} = \mathcal{F}(\mathbf{x}) + \mathbf{x}$};

% Skip connection
\draw[skip] (2, 1.8) -- (3.5, 1.8) -- (3.5, -1.5) -- (plus);
\node[font=\scriptsize, color=gdlred!70] at (3.8, 0) {skip};

% Gradiente
\node[font=\footnotesize, align=center] at (-2, -2.5) {Gradient:\\attenuates};
\node[font=\footnotesize, align=center] at (2, -3.3) {Gradient:\\flows directly};

\end{tikzpicture}
\caption[Comparison of blocks without and with residual connection]{Comparison of blocks without (left) and with (right) a residual connection. The skip connection (red dotted line) allows the gradient to flow directly during backpropagation, avoiding vanishing in deep networks.}
\label{fig:skip_connection}
\end{figure}

Before ResNets, it was empirically observed that deeper networks did not necessarily generalize better, even when the training error was lower. The use of residual connections reduced the vanishing gradient problem.
\section{Neural networks as graphs}\label{sec:redes_grafos}

In the previous sections we have defined neural layers and their composition algebraically, as functions that transform signals. However, there exists an equally important complementary perspective: the representation of neural networks as \textit{graphs}, where nodes represent computational units and edges represent the flow of information.

This graph perspective is fundamental for several reasons:
\begin{itemize}
    \item It provides a precise mathematical language to describe the \textit{architecture} of a network
    \item It allows distinguishing between different families of networks (feedforward, recurrent, etc.) by means of topological properties
    \item It is the basis of the \textit{automatic differentiation} used in modern deep learning frameworks
    \item It connects naturally with the field of \gls{GDL}, where the very structure of the graph encodes information about the domain of the data
\end{itemize}

\subsection{Graphs}

So far we have presented neural networks as compositions of layers in a somewhat informal way. To rigorously formalize their structure, we turn to graph theory, which provides the precise mathematical language to describe the connectivity relationships between neurons and layers.

\begin{definition}[\Gls{grafo}]\label{def:grafo_neuronal}
A \textit{\gls{grafo}} is a mathematical structure $\mathcal{G} = (V, E)$ consisting of:
\begin{itemize}
    \item A set $V$ of \textit{vertices} (also called nodes)
    \item A set $E$ of \textit{edges} (also called links or connections), where each edge connects a pair of vertices
\end{itemize}
Formally, for an undirected graph, $E \subseteq \{\{u, v\} : u, v \in V, u \neq v\}$, whereas for a directed graph, $E \subseteq V \times V$\footnote{Note that the definition of a directed graph allows self-loops $(v, v)$, which is relevant for recurrent networks where a neuron can feed back its own output.}.
\end{definition}

Graphs provide a natural abstraction for modeling relationships between objects. In the context of neural networks, the vertices represent neurons (or layers, depending on the level of abstraction) and the edges represent weighted connections that transmit signals.

\begin{definition}[\Gls{grafo_dirigido}]\label{def:grafo_dirigido}
A \textit{\gls{grafo_dirigido}} is a graph $\mathcal{G} = (V, \vec{E})$ where each edge has a specific direction. A directed edge $(u, v) \in \vec{E}$ indicates that there is a connection from vertex $u$ (source) to vertex $v$ (target). The direction matters: $(u, v) \neq (v, u)$.
\end{definition}

\begin{definition}[Undirected graph]\label{def:grafo_no_dirigido}
An \textit{undirected graph} is a graph $\mathcal{G} = (V, E)$ where the edges are symmetric: if $(u, v) \in E$, then $(v, u) \in E$. Equivalently, $E$ can be viewed as a set of unordered pairs $\{\{u, v\} : u, v \in V\}$.
\end{definition}

In an \gls{ANN}, the directed edges represent the flow of information: the output of a neuron $u$ becomes the input to neuron $v$. This directional flow is fundamental to understanding how signals propagate during inference (forward pass) and how gradients are backpropagated during training (backward pass).

\begin{definition}[Path in a directed graph]\label{def:camino}
Given a \gls{grafo_dirigido} $\mathcal{G} = (V, \vec{E})$, a \textit{path} of length $k$ is a sequence of vertices $v_0, v_1, \ldots, v_k$ such that $(v_i, v_{i+1}) \in \vec{E}$ for all $i = 0, \ldots, k-1$.
\end{definition}

\begin{definition}[Cycle in a directed graph]\label{def:ciclo}
Given a \gls{grafo_dirigido} $\mathcal{G} = (V, \vec{E})$, a \textit{cycle} is a path $v_0, v_1, \ldots, v_k$ where $v_0 = v_k$, $k \geq 1$, and the vertices $v_0, v_1, \ldots, v_{k-1}$ are pairwise distinct.
\end{definition}

A cycle in a neural network indicates feedback: the output of a neuron (or layer) can, directly or indirectly, influence its own input. This is fundamental in recurrent networks that process temporal sequences.

\begin{definition}[\Gls{grafo_aciclico_dirigido} (DAG)]\label{def:dag}
A \textit{\gls{grafo_aciclico_dirigido}} (DAG, for Directed Acyclic Graph) is a \gls{grafo_dirigido} that contains no cycles.
\end{definition}

\glspl{DAG} have especially useful computational properties:

\begin{theorem}[Topological ordering]\label{thm:ordenamiento_topologico}
Every directed acyclic graph $\mathcal{G} = (V, \vec{E})$ admits at least one \textit{topological ordering}: a numbering of its vertices $v_1, v_2, \ldots, v_{|V|}$ such that if $(v_i, v_j) \in \vec{E}$, then $i < j$.
\end{theorem}

\begin{proof}
We proceed by induction on the number of vertices $n = |V|$.

\textbf{Base case} ($n = 1$): A graph with a single vertex and no edges (being acyclic) trivially has a topological ordering.

\textbf{Inductive step}: Suppose that every \gls{DAG} with $n$ vertices has a topological ordering. Consider a \gls{DAG} $\mathcal{G} = (V, \vec{E})$ with $n+1$ vertices.

We claim that $\mathcal{G}$ must have at least one vertex $v$ with no incoming edges (in-degree zero). Indeed, if every vertex had at least one incoming edge, we could construct an infinite path by going backwards indefinitely, which by the pigeonhole principle (there are finitely many vertices) would imply a cycle, contradicting that $\mathcal{G}$ is acyclic.

We select a vertex $v \in V$ with in-degree zero and place it first in the ordering. Consider the reduced graph $\mathcal{G}' = (V \setminus \{v\}, \vec{E}')$ where $\vec{E}'$ is $\vec{E}$ restricted to $V \setminus \{v\}$. This graph has $n$ vertices and remains acyclic (removing a vertex does not introduce cycles). By the induction hypothesis, $\mathcal{G}'$ has a topological ordering. Concatenating $v$ at the beginning of this ordering, we obtain a topological ordering of $\mathcal{G}$.
\end{proof}

The existence of a topological ordering guarantees that we can compute the outputs of all neurons in a sequential order without needing to solve systems of implicit equations, which makes feedforward networks computationally efficient.

\subsection{Neural networks as graphs of neurons}

An \gls{ANN} in its most general form can be represented as a \gls{grafo_dirigido} $\mathcal{G}_{\text{neuronas}} = (V_N, E_N)$ where:
\begin{itemize}
    \item $V_N$: set of neurons of the network
    \item $E_N$: set of weighted connections. An edge $(n_i, n_j) \in E_N$ with weight $w_{ij} \in \mathbb{R}$ indicates that the output of neuron $n_i$ is multiplied by $w_{ij}$ and contributes to the input of neuron $n_j$
\end{itemize}

\subsubsection[Feedforward networks: DAGs of neurons]{Feedforward networks: \glspl{DAG} of neurons}

Feedforward, or forward-propagation, neural networks are those whose graph of neurons is a \gls{DAG}. The absence of cycles guarantees that:

\begin{enumerate}
    \item \textbf{Deterministic single-pass computation}: Given an input, we can compute the output of each neuron exactly once by following a topological ordering. If $\sigma_i$ denotes the topological ordering of neuron $i$, then:
    \[
    y_i = \sigma_i\left(\sum_{(j,i) \in E_N} w_{ji} y_j + b_i\right)
    \]
    where $y_j$ has already been computed for all predecessors $j$ of $i$.

    \item \textbf{Efficient parallelization}: All neurons at the same ``level'' of the \gls{DAG} (same distance from the input) can be computed in parallel, which allows leveraging modern hardware architectures (GPUs, TPUs).

    \item \textbf{Efficient backpropagation}: The computation of gradients via the chain rule can be performed in a single backward pass following the reverse topological order. We will see more about learning in \glspl{ANN} in Section~\ref{sec:learning_ann}.
\end{enumerate}

\subsubsection{Recurrent networks: cyclic graphs}

\glspl{RNN} and other architectures with feedback are modeled by directed graphs \textit{with cycles}. A cycle represents memory or internal state: the output at time step $t$ depends not only on the current input, but also on internal states from previous steps.

Mathematically, this translates into equations of the form:
\[
h_t = f(h_{t-1}, x_t; \theta)
\]
where $h_t$ is the hidden state at time $t$, $x_t$ is the current input, and $\theta$ are the parameters. The implicit cycle $(h_{t-1} \to h_t \to h_{t+1} \to \cdots)$ allows the network to capture temporal dependencies.

Although cyclic graphs (recurrent networks) have greater theoretical expressiveness for modeling temporal dynamics, modern architectures tend to favor \glspl{DAG} for practical reasons:

\begin{enumerate}
    \item \textbf{Computational efficiency}: \glspl{DAG} allow massive parallelization, crucial for training models with millions or billions of parameters.

    \item \textbf{Numerical stability}: Recurrent networks suffer from exploding or vanishing gradient problems over long sequences, hindering their training.

    \item \textbf{Conceptual simplicity}: \glspl{DAG} are easier to analyze theoretically and to debug in practice.
\end{enumerate}

Nevertheless, mechanisms such as attention and \gls{Transformer} architectures (see \autoref{ch:attention}) have managed to capture long-range dependencies while maintaining the \gls{DAG} structure, combining the best of both worlds.

\subsection{Layers as graphs of transformations}\label{subsec:capas_grafos_transformaciones}

In addition to understanding \gls{ANN} as graphs where the \textit{nodes are individual neurons}, there exists a complementary perspective at a higher level of abstraction that is fundamental to understanding modern architectures: thinking of networks as \textit{graphs of layers}, where each node represents a complete layer (potentially with thousands of neurons) and the edges represent the flow of data between layers.

\begin{definition}[\Gls{grafo_de_capas}]\label{def:grafo_de_capas}
Given a neural network composed of $k$ layers $L_1, L_2, \ldots, L_k$, its \textit{\gls{grafo_de_capas}} is a \gls{grafo_dirigido} $\mathcal{G}_{\mathcal{N}} = (V_L, E_L, V_{\text{in}}, V_{\text{out}})$ where:
\begin{itemize}
    \item $V_L = \{L_1, L_2, \ldots, L_k\}$: the set of layers
    \item $(L_i, L_j) \in E_L$ if and only if the output of layer $L_i$ is (part of) the input of layer $L_j$
    \item $V_{\text{in}} \subseteq V_L$: the \textit{input layers}, those with in-degree zero (they receive no output from other layers)
    \item $V_{\text{out}} \subseteq V_L$: the \textit{output layers}, those with out-degree zero (their output does not feed other layers)
\end{itemize}
\end{definition}

This level of abstraction simplifies the visualization of complex architectures and facilitates the analysis of the flow of information through the network. Moreover, it naturally reflects the algebraic structure: each layer $L_i$ can be decomposed as a linear transformation (representable as a matrix) followed by a nonlinear activation function, $L_i(x) = \sigma(W_i x + b_i)$, and the edges of the graph correspond to these compositions of transformations between vector spaces.

\subsubsection[Feedforward architectures as DAGs of layers]{Feedforward architectures as \glspl{DAG} of layers}

Classical feedforward architectures (such as multilayer perceptrons, standard convolutional \glspl{CNN}) have graphs of layers that are \glspl{DAG} with a particularly simple structure: a linear path.

\begin{example}[Feedforward neural network]
A feedforward neural network with $k$ layers is represented as:
\[
\mathcal{N}(x) = L_k \circ L_{k-1} \circ \cdots \circ L_2 \circ L_1(x)
\]
Its graph of layers is the directed path:
\[
L_1 \to L_2 \to L_3 \to \cdots \to L_{k-1} \to L_k
\]
This is the simplest possible \gls{DAG}: completely linear, with no branches or skip connections.
\end{example}

\subsubsection[Modern architectures with nonlinear DAGs]{Modern architectures with nonlinear \glspl{DAG}}

Modern architectures have abandoned the restriction of linear graphs, exploring more complex \glspl{DAG} that allow richer flows of information:

Residual connections, introduced in ResNet~\citep{he2016deep}, add edges that ``skip'' layers:
\[
h^{(\ell+1)} = h^{(\ell)} + F(h^{(\ell)}; \theta_\ell)
\]
where $F$ is a learned function. In the graph of layers, this creates a direct edge $L_\ell \to L_{\ell+1}$ that coexists with the path through the intermediate layers, resulting in a nonlinear \gls{DAG}.

Networks such as Inception~\citep{szegedy2015going} process the input through multiple parallel paths that are then concatenated or summed:
\[
h^{(\ell+1)} = \text{Concatenate}(F_1(h^{(\ell)}), F_2(h^{(\ell)}), F_3(h^{(\ell)}))
\]
The graph of layers has an inverted-tree structure: one node (input) branches into multiple paths that are subsequently merged.

\subsubsection{Recurrent networks as cyclic graphs of layers}

\glspl{RNN} and \glspl{LSTM} introduced in Section~\ref{subsubsec:rnn}, by contrast, have graphs of layers with cycles:

\begin{example}[Basic recurrent neural network]
An \gls{RNN} processes a sequence $x_1, x_2, \ldots, x_T$ by maintaining a hidden state $h_t$ that is updated recursively:
\[
h_t = f(W_{hh} h_{t-1} + W_{xh} x_t + b)
\]
The graph of layers contains a self-directed cycle: the recurrent layer has an edge to itself, representing the temporal dependency $h_{t-1} \to h_t$.
\end{example}

\subsubsection{Connection between the graph-of-neurons and graph-of-layers perspectives}

The two perspectives (graph of neurons and graph of layers) are related hierarchically:
\begin{itemize}
    \item The \textbf{graph of neurons} $\mathcal{G}_{\text{neuronas}} = (V_N, E_N)$ provides the complete microscopic description of the network
    \item The \textbf{graph of layers} $\mathcal{G}_{\mathcal{N}} = (V_L, E_L)$ is a macroscopic abstraction where each node $L_i \in V_L$ represents a complete subgraph of neurons
\end{itemize}

Formally, there exists a projection function $\pi: V_N \to V_L$ that assigns each neuron to its corresponding layer. The edges in $E_L$ summarize the connectivity between groups of neurons:
\[
(L_i, L_j) \in E_L \iff \exists n \in \pi^{-1}(L_i), m \in \pi^{-1}(L_j) : (n, m) \in E_N
\]

This two-level abstraction is fundamental: it allows reasoning about global architectures (graphs of layers) without getting lost in the details of millions of individual connections (graph of neurons), while preserving the ability to compute exact gradients at the level of each individual parameter.

\section{Learning in neural networks}\label{sec:learning_ann}

The learning or training of a neural network is an iterative process that involves optimizing its parameters \(\theta\) or weights to minimize an error. The \gls{ANN} is provided with training datasets containing examples of the expected output of the \gls{ANN} given some input data. With the error at each step, an optimization algorithm for the parameters \(\theta\) is run; in particular, gradient descent or one of its variants such as stochastic gradient descent (\gls{SGD}) is used.

\subsection{Loss function}

The error of the \gls{ANN} is a measure of the difference between the expected output and the network's output. This measure of the error is performed with a function called the loss function (\emph{loss function}) or \textit{cost function}.

\begin{definition}[Loss function]\label{def:funcion_perdida}
	A \textit{loss function} is a function \(\mathcal{L}: \mathcal{Y} \times \mathcal{Y} \to \mathbb{R}_{\geq 0}\) that measures the discrepancy between the prediction \(\hat{y} \in \mathcal{Y}\) and the true value \(y \in \mathcal{Y}\), satisfying:
	\begin{enumerate}
		\item \textbf{Non-negativity}: \(\mathcal{L}(y, \hat{y}) \geq 0\) for all \(y, \hat{y} \in \mathcal{Y}\)
		\item \textbf{Consistency}: \(\mathcal{L}(y, y) = 0\) for all \(y \in \mathcal{Y}\) (the perfect prediction incurs no loss)
		\item \textbf{Differentiability}: \(\mathcal{L}\) is differentiable with respect to \(\hat{y}\) to enable gradient-based optimization
	\end{enumerate}
\end{definition}

\begin{remark}
Some loss functions satisfy the stronger condition \(\mathcal{L}(y, \hat{y}) = 0 \iff y = \hat{y}\) (for example, MSE and MAE), while others attain the minimum of zero over a wider region: the hinge loss (\autoref{def:hinge_loss}) satisfies \(\mathcal{L}(y, \hat{y}) = 0\) when \(y \hat{y} \geq 1\), not only when \(y = \hat{y}\).
\end{remark}

The loss function can be used in more than one layer of the \gls{ANN}, several can be combined, and it can even be used with the parameters \(\theta\) of the \gls{ANN} to prevent them from growing excessively\footnote{Applying a loss function over the weights of a network is a process called regularization and serves to avoid or reduce cumulative errors due to the representation capacity of computers, and also to force greater generalization of the learning.}. In the context of this work we will consider only a single loss function applied over the last layer of the \gls{ANN}, which will consist of more than one output.


\subsubsection{Common loss functions}

We now present the most widely used loss functions in practice, organized by the type of task.

\paragraph{Regression}

\textit{Regression} problems consist of predicting a continuous numerical value from the input features. Mathematically, we seek to approximate a function \(f: \mathcal{X} \to \mathbb{R}^d\) where the output space is a subset of \(\mathbb{R}^d\). Typical examples include predicting house prices, estimating future temperatures, or regressing spatial coordinates. Loss functions for regression measure the distance between the predicted value \(\hat{y}\) and the actual value \(y\) in Euclidean space.

\begin{definition}[Mean squared error (MSE)]\label{def:mse}
	The \textit{mean squared error} (\gls{MSE}) for a set of \(n\) predictions is:
	\[
		\mathcal{L}_{\text{MSE}}(y, \hat{y}) = \frac{1}{n} \sum_{i=1}^{n} (y_i - \hat{y}_i)^2
	\]
	For a single prediction: \(\mathcal{L}(y, \hat{y}) = (y - \hat{y})^2\).
\end{definition}

The \gls{MSE} heavily penalizes large errors due to the quadratic term, which makes it sensitive to outliers. Its gradient is \(\nabla_{\hat{y}} \mathcal{L} = -2(y - \hat{y})\), proportional to the error.

\begin{definition}[Mean absolute error (MAE)]\label{def:mae}
	The \textit{mean absolute error} (\gls{MAE}) is:
	\[
		\mathcal{L}_{\text{MAE}}(y, \hat{y}) = \frac{1}{n} \sum_{i=1}^{n} |y_i - \hat{y}_i|
	\]
\end{definition}

The \gls{MAE} is more robust to outliers than the \gls{MSE}, but it has a problem: it is not differentiable at \(\hat{y} = y\). In practice, the subgradient \(\text{sign}(y - \hat{y})\) is used.

\begin{definition}[Huber loss]\label{def:huber}
	The \textit{Huber loss} combines the advantages of \gls{MSE} and \gls{MAE}:
	\[
		\mathcal{L}_{\delta}(y, \hat{y}) = \begin{cases}
			\frac{1}{2}(y - \hat{y})^2 & \text{if } |y - \hat{y}| \leq \delta \\
			\delta |y - \hat{y}| - \frac{1}{2}\delta^2 & \text{otherwise}
		\end{cases}
	\]
	where \(\delta > 0\) is a hyperparameter that controls the transition.
\end{definition}

The Huber loss is quadratic for small errors (providing smooth gradients) and linear for large errors (robustness to outliers). It is differentiable everywhere, including at the transition at \(|y - \hat{y}| = \delta\).

\paragraph{Classification}

\textit{Classification} problems consist of assigning each input to one of \(K\) predefined discrete categories. Mathematically, we seek to approximate a function \(f: \mathcal{X} \to \{1, \ldots, K\}\), although in practice the network produces a vector of probabilities \(\hat{y} \in \Delta^{K-1}\) (the simplex of dimension \(K-1\)). When \(K = 2\) we speak of \textit{binary classification} (for example, detecting \textit{spam} email); when \(K > 2\), of \textit{multiclass classification} (for example, recognizing handwritten digits from 0 to 9, or classifying images into categories). Loss functions for classification penalize the discrepancies between the predicted distribution and the actual class.

\begin{definition}[Cross-entropy]\label{def:cross_entropy}
	The \textit{cross-entropy} measures the divergence between the true class distribution \(y\) and the estimated distribution \(\hat{y}\):
	\begin{equation}
		\mathcal{L}(y, \hat{y}) = -\sum_{i=1}^{K} y_i \log(\hat{y}_i)
		\label{eq:cross_entropy}
	\end{equation}
	where $K$ is the number of classes, \(y_i\) is the true probability and \(\hat{y}_i\) the probability estimated by the network.
\end{definition}

For the values of the last layer to be probabilities, the softmax function (\autoref{def:softmax}) is applied.

\begin{definition}[Binary logistic loss]\label{def:binary_cross_entropy}
	For binary classification with \(y \in \{0, 1\}\) and \(\hat{y} \in (0, 1)\):
	\[
		\mathcal{L}(y, \hat{y}) = -y \log(\hat{y}) - (1-y) \log(1 - \hat{y})
	\]
	This is the cross-entropy specialized to the binary case, also known as \textit{log loss}.
\end{definition}

\begin{definition}[Hinge loss]\label{def:hinge_loss}
	The \textit{hinge loss}, used in support vector machines (SVM), for \(y \in \{-1, +1\}\):
	\[
		\mathcal{L}(y, \hat{y}) = \max(0, 1 - y \cdot \hat{y})
	\]
\end{definition}

The hinge loss does not penalize correct predictions with margin greater than 1, which induces maximum-margin solutions. It is not differentiable at \(y \cdot \hat{y} = 1\).

\paragraph{For probability distributions}

In addition to the standard regression and classification problems, some applications require the network to produce a complete probability distribution as output, not just a point or a class. This is common in generative models, where the network learns to generate data following a target distribution, and in knowledge distillation, where a student network learns to mimic the probability distributions of a teacher network. The following loss functions measure the divergence or distance between probability distributions.

\begin{definition}[Kullback-Leibler divergence]\label{def:kl_divergence}
	The \textit{KL divergence} between two distributions \(P\) and \(Q\) over a discrete space \(\mathcal{X}\) is:
	\[
		D_{\text{KL}}(P \| Q) = \sum_{x \in \mathcal{X}} P(x) \log \frac{P(x)}{Q(x)}
	\]
	For continuous distributions, the sum is replaced by an integral.
\end{definition}

The \gls{KL} divergence measures how much information is lost when using \(Q\) to approximate \(P\). Note that \(D_{\text{KL}}(P \| Q) \neq D_{\text{KL}}(Q \| P)\) in general, so it is not symmetric. The cross-entropy is related to \gls{KL} by: \(H(P, Q) = H(P) + D_{\text{KL}}(P \| Q)\), where \(H(P)\) is the entropy of \(P\) and \(H(P, Q)\) the cross-entropy.

The \gls{KL} divergence is a special case of a more general family of divergences:

\begin{definition}[\(f\)-divergence]\label{def:f_divergence}
	Let \(f: (0, \infty) \to \mathbb{R}\) be a convex function with \(f(1) = 0\). The \textit{\(f\)-divergence} between distributions \(P\) and \(Q\) is:
	\[
		D_f(P \| Q) = \sum_{x \in \mathcal{X}} Q(x) \, f\left(\frac{P(x)}{Q(x)}\right)
	\]
	For continuous distributions with densities \(p\) and \(q\): \(D_f(P \| Q) = \int q(x) \, f\left(\frac{p(x)}{q(x)}\right) dx\).
\end{definition}

The \(f\)-divergences unify multiple measures of discrepancy between distributions:
\begin{itemize}
	\item \textbf{KL divergence}: \(f(t) = t \log t\) produces \(D_{\text{KL}}(P \| Q)\)
	\item \textbf{Reverse KL}: \(f(t) = -\log t\) produces \(D_{\text{KL}}(Q \| P)\)
	\item \textbf{Jensen-Shannon}: \(f(t) = t \log t - (t+1) \log \frac{t+1}{2}\) produces a symmetric divergence
	\item \textbf{Total variation}: \(f(t) = \frac{1}{2}|t - 1|\) produces \(\text{TV}(P, Q) = \frac{1}{2}\sum_x |P(x) - Q(x)|\)
	\item \textbf{Chi-squared}: \(f(t) = (t-1)^2\) produces the \(\chi^2\) divergence
\end{itemize}

The key property of \(f\)-divergences is that \(D_f(P \| Q) \geq 0\) with equality if and only if \(P = Q\), which makes them suitable as loss functions for comparing distributions.

\paragraph{For representation learning}

\textit{Representation learning} seeks to build a function \(f: \mathcal{X} \to \mathbb{R}^d\) that maps input data to a space of \textit{embeddings} where the geometry reflects the semantics of the problem. A particular case is \textit{metric learning}, where the objective is for the distance in the embedding space to encode similarity: semantically similar elements (same class, same object, same speaker) should be close, while dissimilar elements should be separated. Typical applications include face recognition (verifying whether two images are of the same person), similar-image search, and recommendation systems. Loss functions for metric learning operate over pairs or triplets of examples, not over individual examples.

\begin{definition}[Contrastive loss]\label{def:contrastive_loss}
	The \textit{contrastive loss}~\citep{hadsell2006dimensionality} for a pair of embeddings \((f(x_1), f(x_2))\) with label \(y \in \{0, 1\}\) (0 = similar, 1 = dissimilar):
	\[
		\mathcal{L}(y, x_1, x_2) = (1-y) \frac{1}{2} d^2 + y \frac{1}{2} \max(0, m - d)^2
	\]
	where \(d = \|f(x_1) - f(x_2)\|_2\) is the Euclidean distance and \(m > 0\) is the margin.
\end{definition}

The contrastive loss penalizes similar pairs that are far apart and dissimilar pairs that are too close (within the margin \(m\)).

\begin{definition}[Triplet loss]\label{def:triplet_loss}
	The \textit{triplet loss}~\citep{schroff2015facenet} for a triplet (anchor \(a\), positive \(p\), negative \(n\)):
	\[
		\mathcal{L}(a, p, n) = \max(0, \|f(a) - f(p)\|_2 - \|f(a) - f(n)\|_2 + m)
	\]
	where \(m > 0\) is the margin that ensures separation between classes.
\end{definition}

The triplet loss forces the anchor to be closer to the positive than to the negative by a margin \(m\). It is widely used in face recognition and similarity verification.

\subsubsection{Loss functions vs. metrics}

To understand why many loss functions are not metrics, we first establish the formal definition of a metric.

\begin{definition}[Metric]\label{def:metrica}
	A \textit{metric} (or \textit{distance}) on a set \(X\) is a function \(d: X \times X \to \mathbb{R}_{\geq 0}\) that satisfies, for all \(x, y, z \in X\):
	\begin{enumerate}
		\item \textbf{Non-negativity}: \(d(x, y) \geq 0\)
		\item \textbf{Identity of indiscernibles}: \(d(x, y) = 0 \iff x = y\)
		\item \textbf{Symmetry}: \(d(x, y) = d(y, x)\)
		\item \textbf{Triangle inequality}: \(d(x, z) \leq d(x, y) + d(y, z)\)
	\end{enumerate}
\end{definition}

It is important to emphasize that many common loss functions \textbf{are not metrics} in the sense of \autoref{def:metrica}. For example:
\begin{itemize}
	\item The \textbf{cross-entropy} does not satisfy symmetry: \(\mathcal{L}(y, \hat{y}) \neq \mathcal{L}(\hat{y}, y)\)
	\item The \textbf{mean squared error} (MSE) does define a metric when the square root is taken
	\item The \textbf{Kullback-Leibler divergence} satisfies neither symmetry nor the triangle inequality
\end{itemize}

What is crucial for machine learning is not that \(\mathcal{L}\) be a metric, but that it:
\begin{itemize}
	\item Provides informative gradients for optimization
	\item Penalizes large errors more than small ones
	\item Is computationally efficient
	\item Aligns with the objective of the task (classification, regression, etc.)
\end{itemize}

Both metrics and loss functions measure ``distance'' between values, but they have fundamentally different purposes:
\begin{itemize}
	\item \textbf{Metrics}: Define a topology over the space, allowing one to speak of convergence, continuity and completeness. They are essential in mathematical analysis.
	\item \textbf{Loss functions}: Designed to guide gradient-based optimization, they prioritize the quality of the error signal over topological properties.
\end{itemize}

\subsection{Gradient descent}\label{subsec:descenso_gradiente}

Once the loss function has been defined as a quantitative measure of the network's error, the next step is to determine how to adjust the parameters to minimize it. The most widespread method for training neural networks is supervised learning via error backpropagation\footnote{There also exist unsupervised and self-supervised training methods that fall outside the scope of this work. Nevertheless, the underlying idea is common: to have some loss function that provides a signal about the network's error. For example, \acrlong{GAN} (\gls{GAN})~\citep{goodfellow2014generative} employ an adversarial loss function in which two networks compete against each other.}. This procedure can be summarized in the following steps:

\begin{enumerate}
	\item \textbf{Parameter initialization}: The weights and biases of the network are typically initialized with small random values.
	\item \textbf{Forward propagation (forward pass)}: For a given input, the information propagates from the input layer, through the hidden layers, to the output layer. This produces a prediction.
	\item \textbf{Loss computation}: The prediction is compared with the true output and a loss value is computed.
	\item \textbf{Backpropagation (backpropagation~\citep{backpropagation})}: The gradient of the loss function with respect to each of the network's weights is computed using the chain rule. This gradient indicates the direction in which the weights should be adjusted to reduce the loss. Currently, automatic differentiation techniques~\citep{Neidinger2010IntroductionTA} are used to compute these gradients efficiently.
	\item \textbf{Weight update}: The weights are adjusted using optimization methods such as stochastic gradient descent (\gls{SGD}) or more advanced variants (\gls{Adam}, RMSProp, Adagrad). This adjustment is performed many times, iterating over the dataset, until the network converges to a set of parameters that minimize the loss.
\end{enumerate}

The result of this process is a network that, ideally, generalizes well to unseen data. This means that the network does not merely memorize the training samples, but learns underlying patterns that can be applied to new inputs.

In general it is not feasible to apply one iteration of the above algorithm over the entire training dataset; instead it is performed in batches. A small subset of the training data is taken and one iteration is applied. And this is repeated until a predefined stopping criterion is reached (a sufficiently small error or a sufficient number of iterations). This algorithm is called \gls{SGD}.

\subsubsection{Mathematical formalization}

We now formalize the fundamental optimization algorithms~\citep{robbins1951stochastic, bottou2018optimization}.

\begin{definition}[Gradient descent]\label{def:gradient_descent}
	Let \(\mathcal{L}: \mathbb{R}^p \to \mathbb{R}\) be a differentiable loss function. \textit{Gradient descent} is the iterative algorithm defined by:
	\[
		\theta_{t+1} = \theta_t - \eta \nabla_\theta \mathcal{L}(\theta_t)
	\]
	where \(\eta > 0\) is the \textit{learning rate} and \(\nabla_\theta \mathcal{L}(\theta_t)\) is the gradient of the loss evaluated at \(\theta_t\).
\end{definition}

The gradient descent algorithm follows the direction of maximum decrease of the loss function. However, computing the exact gradient over all the training data is computationally expensive for large datasets.

\autoref{fig:loss_landscapes} illustrates the fundamental difference between convex loss surfaces (with a single global minimum) and non-convex ones (typical in deep learning, with multiple local minima). This distinction has crucial implications for the convergence of optimization algorithms.

\begin{figure}[h]
\centering
\includegraphics[width=0.95\textwidth]{figures/loss_landscapes}
\caption[Convex vs non-convex loss surfaces]{Comparison of loss surfaces. Left: Convex function with a single global minimum where any gradient descent trajectory converges to the optimum. Right: Non-convex function typical of deep neural networks, with multiple local minima. Recent studies suggest that in overparameterized networks, most local minima have loss values similar to the global one.}
\label{fig:loss_landscapes}
\end{figure}

\begin{definition}[Stochastic gradient descent (SGD)]\label{def:sgd}
	Let \(\{(x_i, y_i)\}_{i=1}^N\) be a dataset and \(\mathcal{L}(\theta) = \frac{1}{N}\sum_{i=1}^N \ell(f_\theta(x_i), y_i)\) the empirical loss. \textit{Stochastic gradient descent} approximates the gradient using a random subset (\textit{mini-batch}) \(\mathcal{B} \subset \{1,\ldots,N\}\):
\[
	\theta_{t+1} = \theta_t - \eta \cdot \frac{1}{|\mathcal{B}|}\sum_{i \in \mathcal{B}} \nabla_\theta \ell(f_\theta(x_i), y_i)
\]
where \(|\mathcal{B}|\) is the mini-batch size.
\end{definition}

The gradient estimated by \gls{SGD} is an \textit{unbiased} estimator of the true gradient:
\[
	\mathbb{E}_{\mathcal{B}}\left[\frac{1}{|\mathcal{B}|}\sum_{i \in \mathcal{B}} \nabla_\theta \ell(f_\theta(x_i), y_i)\right] = \nabla_\theta \mathcal{L}(\theta)
\]
Although it has higher variance than the exact gradient, this variance can act as implicit regularization, helping to escape suboptimal local minima.

\subsubsection{Backpropagation}\label{subsubsec:backpropagation}

The backpropagation algorithm~\citep{backpropagation} is the standard method for computing the gradients needed in gradient descent. Its efficiency lies in exploiting the graph structure of the neural network.

\begin{definition}[Computational graph]\label{def:computational_graph}
	A \textit{computational graph} is a directed acyclic graph \(G = (V, E)\) where:
	\begin{itemize}
		\item Each node \(v \in V\) represents a variable or a differentiable operation \(f_v\)
		\item The input nodes (with no incoming edges) correspond to parameters \(\theta\) and input data \(x\)
		\item Each edge \((u, v) \in E\) indicates that \(v\) depends on \(u\)
		\item The output node (unique, with no outgoing edges) produces the value of the loss \(\mathcal{L}\)
	\end{itemize}
\end{definition}

The computational graph allows decomposing the gradient computation into local operations via the chain rule, as illustrated in \autoref{fig:computational_graph}.

\begin{figure}[h]
\centering
\begin{tikzpicture}[
    node distance=1.2cm,
    inputnode/.style={circle, draw, thick, minimum size=0.8cm, fill=gdlblue!20},
    opnode/.style={circle, draw, thick, minimum size=0.8cm, fill=gdlgreen!20},
    lossnode/.style={circle, draw, thick, minimum size=0.8cm, fill=gdlred!20},
    arrow/.style={->, thick, >=stealth}
]

% Nodos de entrada
\node[inputnode] (x) at (0,0) {$x$};
\node[inputnode] (w) at (0,1.5) {$w$};
\node[inputnode] (b) at (0,-1.5) {$b$};

% Operaciones intermedias
\node[opnode] (mult) at (2,0.75) {$\times$};
\node[opnode] (add) at (4,0) {$+$};
\node[opnode] (sigma) at (6,0) {$\sigma$};
\node[opnode] (diff) at (8,0) {$-$};
\node[opnode] (sq) at (10,0) {$(\cdot)^2$};

% Nodo de pérdida
\node[lossnode] (L) at (12,0) {$\mathcal{L}$};

% Target
\node[inputnode] (y) at (8,1.5) {$y$};

% Conexiones
\draw[arrow] (x) -- (mult);
\draw[arrow] (w) -- (mult);
\draw[arrow] (mult) -- node[above, font=\scriptsize] {$wx$} (add);
\draw[arrow] (b) -- (add);
\draw[arrow] (add) -- node[above, font=\scriptsize] {$z$} (sigma);
\draw[arrow] (sigma) -- node[above, font=\scriptsize] {$\hat{y}$} (diff);
\draw[arrow] (y) -- (diff);
\draw[arrow] (diff) -- node[above, font=\scriptsize] {$e$} (sq);
\draw[arrow] (sq) -- node[above, font=\scriptsize] {$\frac{1}{2}$} (L);

% Etiquetas de tipos de nodos
\node[font=\footnotesize, gdlgray] at (0,-2.3) {Inputs};
\node[font=\footnotesize, gdlgray] at (6,-1.5) {Operations};
\node[font=\footnotesize, gdlgray] at (12,-1) {Loss};

% Forward y Backward
\draw[dashed, gdlgray, ->] (0.5,-2.8) -- (11.5,-2.8) node[midway, below, font=\scriptsize] {Forward pass};
\draw[dashed, gdlgray, <-] (0.5,-3.3) -- (11.5,-3.3) node[midway, below, font=\scriptsize] {Backward pass};

\end{tikzpicture}
\caption[Computational graph for a neuron]{Computational graph for a simple neuron with MSE loss. The blue nodes represent inputs (parameters $w$, $b$ and datum $x$), the green ones represent differentiable operations ($\times$, $+$, $\sigma$, $-$, $(\cdot)^2$), and the red one represents the final loss $\mathcal{L}$. The edges indicate dependencies: $z = wx + b$, $\hat{y} = \sigma(z)$, $\mathcal{L} = \frac{1}{2}(y - \hat{y})^2$.}
\label{fig:computational_graph}
\end{figure}

The three graphs defined in this section represent the same network at different levels of abstraction:
\begin{itemize}
	\item The \textbf{graph of layers} (\autoref{def:grafo_de_capas}) is the macroscopic view: each node is a complete layer ($L_1, L_2, \ldots$)
	\item The \textbf{graph of neurons} $\mathcal{G}_{\text{neuronas}}$ is the mesoscopic view: each node is an individual neuron with its weighted connections
	\item The \textbf{computational graph} (\autoref{def:computational_graph}) is the microscopic view: each node is an elementary operation ($\times$, $+$, $\sigma$, etc.)
\end{itemize}
The computational graph ``unfolds'' each neuron of the graph of neurons into its constituent operations (multiplication by weight, addition of the bias, application of the activation), allowing the chain rule to be applied systematically during backpropagation.

\begin{theorem}[Chain rule on computational graphs]\label{thm:backprop}
	Let \(G = (V, E)\) be a computational graph with output node \(\mathcal{L}\). For each intermediate node \(v\) with successors \(\{u_1, \ldots, u_k\}\) in the graph (that is, nodes that depend on \(v\)):
	\[
		\frac{\partial \mathcal{L}}{\partial v} = \sum_{i=1}^{k} \frac{\partial \mathcal{L}}{\partial u_i} \cdot \frac{\partial u_i}{\partial v}
	\]
	The backpropagation algorithm computes all the gradients \(\frac{\partial \mathcal{L}}{\partial v}\) for \(v \in V\) in time \(O(|V| + |E|)\), that is, linear in the size of the graph.
\end{theorem}

\begin{proof}
	We prove the correctness of the algorithm and analyze its complexity.

	\textbf{Existence of a topological ordering.}
	Since \(G = (V, E)\) is a directed acyclic graph (DAG), there exists a topological ordering \(v_1, v_2, \ldots, v_n\) of the nodes such that if \((v_i, v_j) \in E\), then \(i < j\). This ordering can be computed in time \(O(|V| + |E|)\) by standard algorithms (for example, DFS or iterative removal of sources).

	\textbf{Forward pass.}
	We traverse the nodes in topological order. For each node \(v_i\):
	\begin{itemize}
		\item If \(v_i\) is an input node (parameter or datum), its value is already defined.
		\item If \(v_i\) is an operation node with predecessors \(\{u_1, \ldots, u_k\}\), we compute \(v_i = f_{v_i}(u_1, \ldots, u_k)\).
	\end{itemize}
	By the property of the topological ordering, all predecessors of \(v_i\) have smaller indices and have already been evaluated. The cost of evaluating each node is \(O(\text{in-degree})\), summing in total to \(O(|V| + |E|)\).

	\textbf{Backward pass: Correctness.}
	To prove the chain-rule formula, let \(v\) be an intermediate node and let \(u_1, \ldots, u_k\) be its direct successors (nodes that depend on \(v\)). The loss \(\mathcal{L}\) depends on \(v\) only through these successors. By the multivariable chain rule:
	\[
		\frac{\partial \mathcal{L}}{\partial v} = \sum_{i=1}^{k} \frac{\partial \mathcal{L}}{\partial u_i} \cdot \frac{\partial u_i}{\partial v}
	\]
	where \(\frac{\partial u_i}{\partial v}\) is the local partial derivative of the operation at \(u_i\) with respect to its input \(v\).

	\textbf{Backward pass: Algorithm.}
	We traverse the nodes in reverse topological order \(v_n, v_{n-1}, \ldots, v_1\):
	\begin{enumerate}
		\item \textbf{Initialization}: For the output node \(\mathcal{L}\), we set \(\frac{\partial \mathcal{L}}{\partial \mathcal{L}} = 1\).
		\item \textbf{Propagation}: For each node \(v_i\), we have already computed the gradients \(\frac{\partial \mathcal{L}}{\partial u_j}\) of all its successors \(u_j\) (they have larger indices). We compute:
		\[
			\frac{\partial \mathcal{L}}{\partial v_i} = \sum_{(v_i, u_j) \in E} \frac{\partial \mathcal{L}}{\partial u_j} \cdot \frac{\partial u_j}{\partial v_i}
		\]
	\end{enumerate}

	\textbf{Complexity analysis.}
	In the backward pass, each edge \((v_i, u_j) \in E\) contributes exactly once to the gradient sum of \(v_i\). The cost of computing each local partial derivative \(\frac{\partial u_j}{\partial v_i}\) is \(O(1)\) for elementary operations. Therefore:
	\begin{itemize}
		\item Each node is visited once: contribution \(O(|V|)\)
		\item Each edge is processed once: contribution \(O(|E|)\)
	\end{itemize}

	The total complexity of the full algorithm (forward + backward) is \(O(|V| + |E|)\).
\end{proof}

\paragraph{Automatic differentiation}
In practice, modern deep learning frameworks (PyTorch, TensorFlow, JAX) implement automatic differentiation~\citep{Neidinger2010IntroductionTA} transparently: during the forward pass, they dynamically build the computational graph by recording each operation; during the backward pass, they traverse this graph in reverse topological order applying the chain rule. This frees the user from computing derivatives manually, allowing them to concentrate on designing architectures.

\subsubsection{Numerical example: XOR}

To concretely illustrate the learning process, let us consider the XOR problem with a minimal neural network of architecture 2-2-1 (2 inputs, 2 hidden neurons, 1 output).

\textbf{Training data:}
\begin{center}
	\begin{tabular}{cc|c}
		\(x_1\) & \(x_2\) & \(y\) \\
		\hline
		0       & 0       & 0     \\
		0       & 1       & 1     \\
		1       & 0       & 1     \\
		1       & 1       & 0     \\
	\end{tabular}
\end{center}

\textbf{Architecture:}
\begin{itemize}
	\item Input layer: \(\mathbf{x} = [x_1, x_2]^T\)
	\item Hidden layer: \(\mathbf{h} = \sigma(W^{(1)}\mathbf{x} + \mathbf{b}^{(1)})\) with \(W^{(1)} \in \mathbb{R}^{2 \times 2}\), \(\mathbf{b}^{(1)} \in \mathbb{R}^2\)
	\item Output layer: \(\hat{y} = \sigma(W^{(2)}\mathbf{h} + b^{(2)})\) with \(W^{(2)} \in \mathbb{R}^{1 \times 2}\), \(b^{(2)} \in \mathbb{R}\)
	\item Activation function: \(\sigma(z) = \frac{1}{1 + e^{-z}}\) (logistic sigmoid, see \autoref{tab:funciones_activacion})
\end{itemize}

\textbf{Example of forward propagation:}

Suppose initial weights:
\begin{gather*}
    W^{(1)} = \begin{pmatrix}
		          0.5 & -0.3 \\
		          0.2 & 0.4
	\end{pmatrix}, \quad
	\mathbf{b}^{(1)} = \begin{pmatrix}
		                   -0.1 \\
		                   0.3
	\end{pmatrix}\\
    W^{(2)} = \begin{pmatrix}
		          0.7 & -0.5
	\end{pmatrix}, \quad
	b^{(2)} = 0.2\\
\end{gather*}

For the input \(\mathbf{x} = [1, 0]^T\):

1. Hidden layer:
\begin{gather*}
    \mathbf{z}^{(1)} = W^{(1)}\mathbf{x} + \mathbf{b}^{(1)} = \begin{pmatrix}
		                                                          0.5 \cdot 1 + (-0.3) \cdot 0 - 0.1 \\
		                                                          0.2 \cdot 1 + 0.4 \cdot 0 + 0.3
	\end{pmatrix} = \begin{pmatrix}
		                0.4 \\
		                0.5
	\end{pmatrix}\\
    \mathbf{h} = \sigma(\mathbf{z}^{(1)}) = \begin{pmatrix}
		                                        \sigma(0.4) \\
		                                        \sigma(0.5)
	\end{pmatrix} = \begin{pmatrix}
		                0.599 \\
		                0.622
	\end{pmatrix}\\
\end{gather*}

2. Output layer:
\begin{gather*}
    z^{(2)} = W^{(2)}\mathbf{h} + b^{(2)} = 0.7 \cdot 0.599 + (-0.5) \cdot 0.622 + 0.2 = 0.308\\
    \hat{y} = \sigma(0.308) = 0.576\\
\end{gather*}

3. Error (using \gls{MSE}):
\[
	\mathcal{L} = \frac{1}{2}(y - \hat{y})^2 = \frac{1}{2}(1 - 0.576)^2 = 0.090
\]

\textbf{Backpropagation:}

To compute the gradients systematically, we define the \textit{error signal} (or delta) of each layer. For the output layer:
\[
	\delta^{(2)} = \frac{\partial \mathcal{L}}{\partial z^{(2)}} = -(y - \hat{y}) \cdot \sigma'(z^{(2)}) = -(y - \hat{y}) \cdot \hat{y}(1 - \hat{y}) \in \mathbb{R}
\]
where we use that the derivative of the sigmoid is \(\sigma'(z) = \sigma(z)(1-\sigma(z))\).

The gradient with respect to the weights of the output layer is then:
\[
	\frac{\partial \mathcal{L}}{\partial W^{(2)}} = \delta^{(2)} \cdot \mathbf{h}^T \in \mathbb{R}^{1 \times 2}
\]
For the hidden layer, the error signal propagates backward:
\[
	\boldsymbol{\delta}^{(1)} = \left( W^{(2)T} \cdot \delta^{(2)} \right) \odot \sigma'(\mathbf{z}^{(1)}) = \left( W^{(2)T} \cdot \delta^{(2)} \right) \odot \left[ \mathbf{h} \odot (1 - \mathbf{h}) \right] \in \mathbb{R}^{2}
\]
where \(\odot\) denotes the Hadamard product (element-wise).

The gradient with respect to the weights of the hidden layer is:
\[
	\frac{\partial \mathcal{L}}{\partial W^{(1)}} = \boldsymbol{\delta}^{(1)} \cdot \mathbf{x}^T \in \mathbb{R}^{2 \times 2}
\]
This is an outer product between the vector of error signals and the input vector.

\textbf{Numerical computations:}

We compute the intermediate values step by step. First, the error signal of the output layer:
\begin{gather*}
	\delta^{(2)} = -(y - \hat{y}) \cdot \hat{y}(1 - \hat{y}) = -(1 - 0.576) \cdot 0.576 \cdot (1 - 0.576) = -0.424 \cdot 0.244 = -0.103
\end{gather*}

For the gradients of the output layer:
\begin{gather*}
	\frac{\partial \mathcal{L}}{\partial W^{(2)}} = \delta^{(2)} \cdot \mathbf{h}^T = -0.103 \cdot \begin{pmatrix} 0.599 & 0.622 \end{pmatrix} = \begin{pmatrix} -0.062 & -0.064 \end{pmatrix} \\[0.3em]
	\frac{\partial \mathcal{L}}{\partial b^{(2)}} = \delta^{(2)} = -0.103
\end{gather*}

For the hidden layer, we propagate the error backward:
\begin{gather*}
	\boldsymbol{\delta}^{(1)} = \delta^{(2)} \cdot W^{(2)T} \odot \mathbf{h} \odot (1 - \mathbf{h}) \\
	= -0.103 \cdot \begin{pmatrix} 0.7 \\ -0.5 \end{pmatrix} \odot \begin{pmatrix} 0.599 \cdot 0.401 \\ 0.622 \cdot 0.378 \end{pmatrix}
	= \begin{pmatrix} -0.072 \\ 0.052 \end{pmatrix} \odot \begin{pmatrix} 0.240 \\ 0.235 \end{pmatrix}
	= \begin{pmatrix} -0.017 \\ 0.012 \end{pmatrix}
\end{gather*}

\begin{gather*}
	\frac{\partial \mathcal{L}}{\partial W^{(1)}} = \boldsymbol{\delta}^{(1)} \cdot \mathbf{x}^T = \begin{pmatrix} -0.017 \\ 0.012 \end{pmatrix} \cdot \begin{pmatrix} 1 & 0 \end{pmatrix} = \begin{pmatrix} -0.017 & 0 \\ 0.012 & 0 \end{pmatrix}
\end{gather*}

\textbf{Weight update:}

With learning rate \(\eta = 0.5\), we update the weights of the output layer:
\begin{gather*}
	W^{(2)}_{\text{new}} = W^{(2)} - \eta \cdot \frac{\partial \mathcal{L}}{\partial W^{(2)}} = \begin{pmatrix} 0.7 & -0.5 \end{pmatrix} - 0.5 \cdot \begin{pmatrix} -0.062 & -0.064 \end{pmatrix} = \begin{pmatrix} 0.731 & -0.468 \end{pmatrix} \\[0.3em]
	b^{(2)}_{\text{new}} = b^{(2)} - \eta \cdot \frac{\partial \mathcal{L}}{\partial b^{(2)}} = 0.2 - 0.5 \cdot (-0.103) = 0.252
\end{gather*}

\textbf{Verification:}

We perform a new forward pass with the updated weights to verify that the loss has decreased:
\begin{gather*}
	z^{(2)}_{\text{new}} = W^{(2)}_{\text{new}} \cdot \mathbf{h} + b^{(2)}_{\text{new}} = 0.731 \cdot 0.599 + (-0.468) \cdot 0.622 + 0.252 = 0.399 \\
	\hat{y}_{\text{new}} = \sigma(0.399) = 0.598 \\
	\mathcal{L}_{\text{new}} = \frac{1}{2}(1 - 0.598)^2 = 0.081
\end{gather*}

Indeed, \(\mathcal{L}_{\text{new}} = 0.081 < 0.090 = \mathcal{L}_{\text{previous}}\), confirming that gradient descent reduces the loss. This process is repeated iteratively until the network learns to correctly classify all the inputs of the XOR problem.

\begin{figure}[h]
	\centering
	\begin{tikzpicture}[
		node distance=2cm,
		neuron/.style={circle, draw, thick, minimum size=1cm},
		input/.style={circle, draw, thick, minimum size=0.9cm, fill=gdlblue!20},
		hidden/.style={circle, draw, thick, minimum size=1.0cm, fill=gdlgreen!20},
		output/.style={circle, draw, thick, minimum size=1.0cm, fill=gdlred!20},
		bias/.style={circle, draw, thick, minimum size=0.7cm, fill=gdlgray!20},
		weight/.style={midway, above, sloped, font=\small},
		activation/.style={below=0.5cm, font=\small}
	]

% Input layer
		\node[input] (x1) at (0,2) {\(x_1\)};
		\node[input] (x2) at (0,0) {\(x_2\)};
		\node[bias] (b0) at (2,-2) {\scriptsize\(+1\)};

% Hidden layer (separada más de la entrada)
		\node[hidden] (h1) at (4,2) {\(h_1\)};
		\node[hidden] (h2) at (4,0) {\(h_2\)};
		\node[bias] (b1) at (6,-2) {\scriptsize\(+1\)};

% Output layer
		\node[output] (y) at (7.5,0) {\(\hat{y}\)};

% Connections with weights
		\draw[->, thick] (x1) -- (h1) node[weight, pos=0.2] {\(0.5\)};
		\draw[->, thick] (x1) -- (h2) node[weight, pos=0.2] {\(0.2\)};
		\draw[->, thick] (x2) -- (h1) node[weight, pos=0.2] {\(-0.3\)};
		\draw[->, thick] (x2) -- (h2) node[weight, pos=0.2] {\(0.4\)};
		\draw[->, thick] (h1) -- (y) node[weight, pos=0.2] {\(0.7\)};
		\draw[->, thick] (h2) -- (y) node[weight, pos=0.2] {\(-0.5\)};

% Bias connections (pesos al inicio de la flecha)
		\draw[->, dashed] (b0) -- (h1) node[weight, pos=0.2] {\(-0.1\)};
		\draw[->, dashed] (b0) -- (h2) node[weight, pos=0.2] {\(0.3\)};
		\draw[->, dashed] (b1) -- (y) node[weight, pos=0.2] {\(0.2\)};

% Layer labels
		\node[above=0.3cm of x1] {Input};
		\node[above=0.3cm of h1] {Hidden layer};
		\node[above=0.3cm of y] {Output};

% Activation values for example x=[1,0]
		\node[activation] at (x1) {\(1\)};
		\node[activation] at (x2) {\(0\)};
		\node[activation] at (h1) {\(0.599\)};
		\node[activation] at (h2) {\(0.622\)};
		\node[activation] at (y) {\(0.576\)};

	\end{tikzpicture}
	\caption[Neural network for the XOR problem]{Neural network for the XOR problem with architecture 2-2-1. The synaptic weights of the example and the activations for the input \([1,0]^T\) are shown. The bias nodes (+1) are shown with dashed connections.}
	\label{fig:xor_network}
\end{figure}

\subsubsection{The vanishing gradient problem}\label{subsubsec:vanishing_gradient}

A direct consequence of the chain rule is the \textit{vanishing gradient problem}. According to \autoref{thm:backprop}, the gradient of the loss with respect to the parameters of the first layers is proportional to the product of the derivatives of the activation functions along the entire network:
\[
\frac{\partial \mathcal{L}}{\partial \theta^{(1)}} \propto \prod_{\ell=1}^{L} \sigma'(z^{(\ell)})
\]

If each derivative $|\sigma'(z^{(\ell)})|$ is less than 1 (as happens with saturating functions such as the sigmoid, where $|\sigma'(z)| \leq 0.25$), this product decreases exponentially with the depth $L$. As a result, the gradients of the first layers become numerically zero, preventing their update during training.

This phenomenon was one of the main obstacles to training deep networks in the 1990s. The solution came with the use of non-saturating activation functions such as \gls{ReLU}, where $\sigma'(z) = 1$ for $z > 0$, allowing gradients to flow without attenuation (see \autoref{prop:gradiente_no_saturante}).

\begin{figure}[htbp]
\centering
\begin{tikzpicture}[
    layer/.style={rectangle, draw, thick, minimum width=1cm, minimum height=0.6cm, font=\small},
    gradbar/.style={rectangle, minimum height=0.25cm},
    arrow/.style={->, thick, >=stealth}
]

% === PANEL IZQUIERDO: Sigmoide ===
\begin{scope}[xshift=-4cm]
    % Título
    \node[font=\bfseries] at (0.5, 2.8) {Sigmoid (saturating)};

    % Capas con gradiente decreciente (numeración invertida: L_5 arriba, L_1 abajo)
    \foreach \i/\layer/\grad/\col in {1/5/1.0/100, 2/4/0.65/75, 3/3/0.35/50, 4/2/0.15/30, 5/1/0.05/15} {
        \node[layer, fill=gdlred!\col!white] (s\i) at (0.5, 2-\i*0.9) {$L_\layer$};
        % Barra de gradiente a la izquierda
        \fill[gdlblue!80] (-1.2, 1.85-\i*0.9) rectangle (-1.2+\grad, 2.15-\i*0.9);
        \draw[gdlgray, thin] (-1.2, 1.85-\i*0.9) rectangle (-0.2, 2.15-\i*0.9);
    }

    % Etiqueta eje gradiente
    \node[font=\scriptsize, rotate=90] at (-1.6, 0) {Magnitude $|\nabla|$};

    % Flecha backprop (de arriba hacia abajo: desde la salida hacia la entrada)
    \draw[arrow, color=gdlgray!70, line width=1.5pt] (1.3, 1.6) -- (1.3, -2.3);
    \node[font=\scriptsize, color=gdlgray!70, rotate=90] at (1.7, -0.3) {backprop};

    % Etiqueta resultado
    \node[font=\small, align=center, color=gdlred!70!black] at (0.5, -3.25) {Gradient $\to 0$};
\end{scope}

% === PANEL DERECHO: ReLU ===
\begin{scope}[xshift=4cm]
    % Título
    \node[font=\bfseries] at (0.5, 2.8) {ReLU (non-saturating)};

    % Capas con gradiente estable (numeración invertida: L_5 arriba, L_1 abajo)
    \foreach \i/\layer/\grad in {1/5/1.0, 2/4/0.95, 3/3/0.90, 4/2/0.85, 5/1/0.80} {
        \node[layer, fill=gdlgreen!50!white] (r\i) at (0.5, 2-\i*0.9) {$L_\layer$};
        % Barra de gradiente a la izquierda
        \fill[gdlblue!80] (-1.2, 1.85-\i*0.9) rectangle (-1.2+\grad, 2.15-\i*0.9);
        \draw[gdlgray, thin] (-1.2, 1.85-\i*0.9) rectangle (-0.2, 2.15-\i*0.9);
    }

    % Etiqueta eje gradiente
    \node[font=\scriptsize, rotate=90] at (-1.6, 0) {Magnitude $|\nabla|$};

    % Flecha backprop (de arriba hacia abajo: desde la salida hacia la entrada)
    \draw[arrow, color=gdlgray!70, line width=1.5pt] (1.3, 1.6) -- (1.3, -2.3);
    \node[font=\scriptsize, color=gdlgray!70, rotate=90] at (1.7, -0.3) {backprop};

    % Etiqueta resultado
    \node[font=\small, align=center, color=gdlgreen!50!black] at (0.5, -3.25) {Stable gradient};
\end{scope}

% === ECUACIÓN CENTRAL ===
\node[font=\small, align=center, draw=gdlgray!30, rounded corners, fill=gdlgray!5, inner sep=6pt] at (0, -1.2) {
    $\displaystyle\frac{\partial \mathcal{L}}{\partial \theta_1} \propto \prod_{\ell=1}^{L} \sigma'(z^{(\ell)})$
};

\end{tikzpicture}
\caption[Vanishing gradient problem]{Comparison of the gradient flow during backpropagation. \textbf{Left}: With saturating functions such as the sigmoid, where $|\sigma'(z)| \leq 0.25$, the product of derivatives decreases exponentially. \textbf{Right}: With ReLU, where $\sigma'(z) = 1$ for $z > 0$, the gradient remains stable allowing deeper networks to be trained.}
\label{fig:vanishing_gradient}
\end{figure}

\subsection{Universal approximation theorem}\label{sec:aproximacion_universal}

The universal approximation theorem guarantees that any continuous function can be approximated arbitrarily well by an \gls{ANN} with a single hidden layer. Its proof requires tools from measure theory and functional analysis: the concepts of measure and $\sigma$-algebra underpin the representation (Riesz) and separation (Hahn-Banach) theorems that constitute the core of the proof. We first present the necessary preliminaries and then the statement with its complete proof. For an in-depth treatment of these foundations, see \citet{Folland1999} or \citet{Rudin1976}.

\paragraph{Measure-theoretic preliminaries}

\begin{definition}[$\sigma$-algebra]\label{def:sigma_algebra}
	Let $X$ be a nonempty set. A \textit{$\sigma$-algebra} (or $\sigma$-field) over $X$ is a collection $\mathcal{F}$ of subsets of $X$ satisfying:
	\begin{enumerate}
		\item $X \in \mathcal{F}$
		\item If $A \in \mathcal{F}$, then $A^c = X \setminus A \in \mathcal{F}$ (closed under complements)
		\item If $\{A_n\}_{n=1}^\infty \subseteq \mathcal{F}$, then $\bigcup_{n=1}^\infty A_n \in \mathcal{F}$ (closed under countable unions)
	\end{enumerate}
	The elements of $\mathcal{F}$ are called \textit{measurable sets}.
\end{definition}

\begin{definition}[Measurable space]\label{def:espacio_medible}
	A \textit{measurable space} is a pair $(X, \mathcal{F})$ where $X$ is a nonempty set and $\mathcal{F}$ is a $\sigma$-algebra over $X$.
\end{definition}

\begin{definition}[Borel $\sigma$-algebra]\label{def:sigma_algebra_borel}
	Let $(X, \tau)$ be a topological space. The \textit{Borel $\sigma$-algebra} $\mathcal{B}(X)$ is the smallest $\sigma$-algebra that contains all the open sets of $X$. Its elements are called \textit{Borel sets}.
\end{definition}

In particular, for $\mathbb{R}^n$, the Borel $\sigma$-algebra $\mathcal{B}(\mathbb{R}^n)$ is generated by the open rectangles, or equivalently, by the half-spaces of the form $\{x \in \mathbb{R}^n : w^T x + b > 0\}$.

\begin{definition}[Measure]\label{def:medida}
	Let $(X, \mathcal{F})$ be a measurable space. A \textit{measure} is a function $\mu: \mathcal{F} \to [0, +\infty]$ satisfying:
	\begin{enumerate}
		\item $\mu(\emptyset) = 0$
		\item \textbf{$\sigma$-additivity}: For any sequence $\{A_n\}_{n=1}^\infty$ of disjoint sets in $\mathcal{F}$:
		\[
			\mu\left(\bigcup_{n=1}^\infty A_n\right) = \sum_{n=1}^\infty \mu(A_n)
		\]
	\end{enumerate}
\end{definition}

\begin{definition}[Measure space]\label{def:espacio_medida}
	A \textit{measure space} is a triple $(X, \mathcal{F}, \mu)$ where $X$ is a nonempty set, $\mathcal{F}$ is a $\sigma$-algebra over $X$, and $\mu: \mathcal{F} \to [0, +\infty]$ is a measure. If $\mu(X) < \infty$, we say that $\mu$ is a \textit{finite measure}. If $\mu(X) = 1$, we say that $\mu$ is a \textit{probability measure}.
\end{definition}

\begin{definition}[Finite signed measure]\label{def:medida_con_signo}
	Let $(X, \mathcal{F})$ be a measurable space. A \textit{finite signed measure} is a function $\mu: \mathcal{F} \to \mathbb{R}$ satisfying $\mu(\emptyset) = 0$ and which is $\sigma$-additive: for any sequence $\{A_n\}_{n=1}^\infty$ of disjoint sets in $\mathcal{F}$,
	\[
		\mu\left(\bigcup_{n=1}^\infty A_n\right) = \sum_{n=1}^\infty \mu(A_n)
	\]
	where the series converges absolutely. The \textit{total variation} of $\mu$ is the measure $|\mu|$ defined by
	\[
		|\mu|(A) = \sup \left\{ \sum_{i=1}^n |\mu(A_i)| : \{A_i\}_{i=1}^n \text{ finite partition of } A \right\}
	\]
\end{definition}

\begin{definition}[Lebesgue measure]\label{def:medida_lebesgue}
	The \textit{Lebesgue measure} $\lambda$ on $\mathbb{R}^n$ is the unique measure on $(\mathbb{R}^n, \mathcal{B}(\mathbb{R}^n))$ that extends the notion of $n$-dimensional volume and satisfies:
	\begin{itemize}
		\item For a rectangle $R = [a_1, b_1] \times \cdots \times [a_n, b_n]$: $\lambda(R) = \prod_{i=1}^n (b_i - a_i)$
		\item It is invariant under translations: $\lambda(A + x) = \lambda(A)$ for all $x \in \mathbb{R}^n$
	\end{itemize}
\end{definition}

A set $A$ has \textit{Lebesgue measure zero} (or measure zero) if $\lambda(A) = 0$. Examples include: isolated points, finite sets, countable sets such as $\mathbb{Q}$, and the Cantor set.

\begin{definition}[Measurable function]\label{def:funcion_medible}
	Let $(X, \mathcal{F})$ and $(Y, \mathcal{G})$ be measurable spaces. A function $f: X \to Y$ is \textit{measurable} (or $\mathcal{F}$-$\mathcal{G}$-measurable) if the preimage of every measurable set is measurable:
	\[
		f^{-1}(B) = \{x \in X : f(x) \in B\} \in \mathcal{F} \quad \forall B \in \mathcal{G}
	\]
\end{definition}

When $Y = \mathbb{R}$ with the Borel $\sigma$-algebra, it suffices to verify that $\{x : f(x) > \alpha\} \in \mathcal{F}$ for all $\alpha \in \mathbb{R}$. In particular, every continuous function is measurable with respect to the Borel $\sigma$-algebras.

\begin{definition}[Lebesgue integral]\label{def:integral_lebesgue}
	Let $(X, \mathcal{F}, \mu)$ be a measure space. The \textit{Lebesgue integral} is constructed in three steps:
	\begin{enumerate}
		\item \textbf{Simple functions}: For $s = \sum_{i=1}^n a_i \chi_{A_i}$ with $a_i \geq 0$ and $A_i \in \mathcal{F}$ disjoint, where $\chi_{A}$ is the characteristic function of $A$, we define
		\[
			\int_X s \, d\mu = \sum_{i=1}^n a_i \mu(A_i)
		\]
		\item \textbf{Non-negative functions}: For $f: X \to [0, +\infty]$ measurable,
		\[
			\int_X f \, d\mu = \sup\left\{\int_X s \, d\mu : s \text{ simple, } 0 \leq s \leq f\right\}
		\]
		\item \textbf{General functions}: For $f: X \to \mathbb{R}$ measurable, we write $f = f^+ - f^-$ with $f^+ = \max(f, 0)$ and $f^- = \max(-f, 0)$. We say that $f$ is \textit{integrable} (or $f \in L^1(\mu)$) if $\int_X f^+ \, d\mu < \infty$ and $\int_X f^- \, d\mu < \infty$, and we define
		\[
			\int_X f \, d\mu = \int_X f^+ \, d\mu - \int_X f^- \, d\mu
		\]
	\end{enumerate}
\end{definition}

\begin{definition}[Almost everywhere]\label{def:casi_por_doquier}
	Let $(X, \mathcal{F}, \mu)$ be a measure space. A property $P(x)$ holds \textit{almost everywhere} (a.e.) or \textit{for almost every $x$} if the set of points where it does not hold has measure zero:
	\[
		\mu(\{x \in X : \neg P(x)\}) = 0
	\]
\end{definition}

In the context of the Lebesgue measure on $\mathbb{R}^n$, ``almost everywhere'' means ``except on a set of Lebesgue measure zero''.

\paragraph{Functional-analytic preliminaries}

\begin{definition}[Vector space]\label{def:espacio_vectorial}
	A \textit{vector space} over $\mathbb{R}$ is a set $V$ equipped with two operations, addition $+: V \times V \to V$ and scalar multiplication $\cdot: \mathbb{R} \times V \to V$, satisfying the following axioms for all $u, v, w \in V$ and $\alpha, \beta \in \mathbb{R}$:
	\begin{enumerate}
		\item \textbf{Associativity of addition:} $(u + v) + w = u + (v + w)$
		\item \textbf{Commutativity of addition:} $u + v = v + u$
		\item \textbf{Identity element:} there exists $0 \in V$ such that $v + 0 = v$
		\item \textbf{Inverse element:} for each $v$ there exists $-v \in V$ such that $v + (-v) = 0$
		\item \textbf{Compatibility of scalar multiplication:} $\alpha(\beta v) = (\alpha\beta)v$
		\item \textbf{Unit element:} $1 \cdot v = v$
		\item \textbf{Distributivity with respect to vector addition:} $\alpha(u + v) = \alpha u + \alpha v$
		\item \textbf{Distributivity with respect to scalar addition:} $(\alpha + \beta)v = \alpha v + \beta v$
	\end{enumerate}
\end{definition}

\begin{definition}[Normed vector space]\label{def:espacio_normado}
	A \textit{normed vector space} is a pair $(X, \|\cdot\|)$ where $X$ is a real vector space and $\|\cdot\|: X \to [0, +\infty)$ is a \textit{norm}, that is, a function satisfying:
	\begin{enumerate}
		\item $\|x\| = 0$ if and only if $x = 0$
		\item $\|\alpha x\| = |\alpha| \|x\|$ for all $\alpha \in \mathbb{R}$
		\item $\|x + y\| \leq \|x\| + \|y\|$ (triangle inequality)
	\end{enumerate}
\end{definition}

Fundamental examples include $\mathbb{R}^n$ with the Euclidean norm and the function spaces $L^p(\Omega)$.

\begin{definition}[Cover]\label{def:recubrimiento}
	Let $X$ be a topological space and $K \subseteq X$. A \textit{cover} of $K$ is a family of sets $\{U_\alpha\}_{\alpha \in A}$ such that $K \subseteq \bigcup_{\alpha \in A} U_\alpha$. If each $U_\alpha$ is an open set of $X$, it is called an \textit{open cover}. A \textit{subcover} is a subfamily that is still a cover of $K$.
\end{definition}

\begin{definition}[Compact space]\label{def:espacio_compacto}
	A subset $K$ of a topological space is \textit{compact} if every open cover of $K$ admits a finite subcover.
\end{definition}

Intuitively, compactness generalizes the properties of closed and bounded sets in $\mathbb{R}^n$ to general topological spaces.

\begin{theorem}[Heine-Borel]\label{thm:heine_borel}
	A subset $K \subset \mathbb{R}^n$ is compact if and only if it is closed and bounded \citep[Theorem~2.41]{Rudin1976}.
\end{theorem}

In particular, the unit hypercube $I_m = [0,1]^m$ is compact, which is essential for the universal approximation theorem.

\begin{definition}[Linear functional]\label{def:funcional_lineal}
	Let $X$ be a real vector space. A \textit{linear functional} is a map $f: X \to \mathbb{R}$ satisfying:
	\[
		f(\alpha x + \beta y) = \alpha f(x) + \beta f(y) \quad \forall x, y \in X, \, \forall \alpha, \beta \in \mathbb{R}
	\]
	That is, a linear functional is an element of the algebraic dual space $X^*$.
\end{definition}

\begin{definition}[Continuous linear functional]\label{def:funcional_continuo}
	Let $(X, \|\cdot\|)$ be a normed vector space. A linear functional $f: X \to \mathbb{R}$ is \textit{continuous} (or \textit{bounded}) if there exists a constant $M > 0$ such that:
	\[
		|f(x)| \leq M\|x\| \quad \forall x \in X
	\]
\end{definition}

The set of all continuous linear functionals on $X$ forms the \textit{dual space} $X'$, which is a Banach space with the operator norm $\|f\| = \sup_{\|x\| \leq 1} |f(x)|$.

\begin{definition}[The space $C(K)$]\label{def:espacio_CK}
	Let $K$ be a compact space. The space $C(K)$ is the set of all continuous functions $f: K \to \mathbb{R}$ equipped with the \textit{supremum norm} (or uniform norm):
	\[
		\|f\|_\infty = \sup_{x \in K} |f(x)|
	\]
\end{definition}

With this norm, $C(K)$ is a Banach space (complete normed vector space). In the context of the universal approximation theorem, we work with $C(I_m)$, the space of continuous functions over the unit hypercube.

\begin{theorem}[Lebesgue dominated convergence]\label{thm:convergencia_dominada}
	Let \((X, \mathcal{F}, \mu)\) be a measure space and \(\{f_n\}_{n=1}^\infty\) a sequence of measurable functions \(f_n: X \to \mathbb{R}\) such that:
	\begin{enumerate}
		\item \(f_n(x) \to f(x)\) for almost every \(x \in X\)
		\item There exists an integrable function \(g: X \to \mathbb{R}\) such that \(|f_n(x)| \leq g(x)\) for almost every \(x \in X\) and all \(n \in \mathbb{N}\)
	\end{enumerate}
	Then \(f\) is integrable and:
	\[
		\lim_{n \to \infty} \int_X f_n \, d\mu = \int_X f \, d\mu = \int_X \lim_{n \to \infty} f_n \, d\mu
	\]
	\citep[Theorem~2.24]{Folland1999}.
\end{theorem}

This theorem is fundamental for interchanging limits and integrals, which will be necessary to prove that sigmoidal functions are discriminatory.

\begin{definition}[Convex set]\label{def:conjunto_convexo}
	A subset $C$ of a vector space $V$ is \textit{convex} if for all $x, y \in C$ and $\lambda \in [0,1]$, we have $\lambda x + (1-\lambda)y \in C$. Equivalently, $C$ contains the segment joining any pair of its points.
\end{definition}

\begin{theorem}[Hahn-Banach, geometric form]\label{thm:hahn_banach}
	Let \(X\) be a real topological vector space, and \(A\) and \(B\) disjoint convex subsets of \(X\) with \(A\) open. Then there exist a continuous linear functional \(f: X \to \mathbb{R}\) and a constant \(\alpha \in \mathbb{R}\) such that:
	\[
		f(a) < \alpha \leq f(b) \quad \forall a \in A, b \in B
	\]
	\citep[Theorem~5.5]{Folland1999}.
\end{theorem}

\begin{corollary}[Characterization of density]\label{cor:caracterizacion_densidad}
	Let \(X\) be a normed vector space and \(Y \subset X\) a subspace. Then \(Y\) is dense in \(X\) if and only if the only continuous linear functional that vanishes on \(Y\) is the zero functional.
\end{corollary}

\begin{proof}
	\textbf{(\(\Rightarrow\))} Suppose that $Y$ is dense in $X$. Let $f: X \to \mathbb{R}$ be a continuous linear functional such that $f(y) = 0$ for all $y \in Y$. For any $x \in X$, there exists a sequence $\{y_n\} \subset Y$ with $y_n \to x$. By the continuity of $f$:
	\[
		f(x) = f\!\left(\lim_{n \to \infty} y_n\right) = \lim_{n \to \infty} f(y_n) = 0
	\]
	Therefore $f \equiv 0$.

	\textbf{(\(\Leftarrow\))} Suppose that $Y$ is not dense in $X$. Then $\overline{Y} \neq X$, and there exists $x_0 \in X \setminus \overline{Y}$. Since $\overline{Y}$ is a closed subspace, we take an open ball $B_r(0) \subset X$ and define $A = \overline{Y} + B_r(0)$, which is open and convex, with $A \cap \{x_0\} = \emptyset$ for $r$ sufficiently small (since $x_0 \notin \overline{Y}$). By \autoref{thm:hahn_banach}, there exist a linear functional $f$ and $\alpha \in \mathbb{R}$ such that:
	\[
		f(a) < \alpha \leq f(x_0) \quad \forall a \in A
	\]
	Since $\overline{Y}$ is a subspace, $0 \in \overline{Y} \subset A$, hence $f(0) = 0 < \alpha$. Moreover, for any $y \in Y$ and $t \in \mathbb{R}$, we have $ty \in \overline{Y} \subset A$, so that $tf(y) = f(ty) < \alpha$ for all $t$. This is only possible if $f(y) = 0$. However, $f(x_0) \geq \alpha > 0$, so $f \not\equiv 0$.
\end{proof}

\begin{theorem}[Riesz representation for \(C(K)\)]\label{thm:riesz_representacion}
	Let \(K\) be a compact space and \(C(K)\) the space of real continuous functions on \(K\) with the supremum norm. For every continuous linear functional \(L: C(K) \to \mathbb{R}\), there exists a unique finite signed regular Borel measure \(\mu\) on \(K\) such that:
	\[
		L(f) = \int_K f \, d\mu \quad \forall f \in C(K)
	\]
	Moreover, \(\|L\| = |\mu|(K)\), where \(|\mu|\) is the total variation of \(\mu\) \citep[Theorem~7.17]{Folland1999}.
\end{theorem}

\begin{theorem}[Stone-Weierstrass]\label{thm:stone_weierstrass}
	Let \(K\) be a compact space and \(\mathcal{A} \subset C(K)\) a subalgebra that separates points of \(K\) and does not vanish at any point. Then \(\mathcal{A}\) is dense in \(C(K)\) with respect to the supremum norm \citep[Theorem~7.32]{Rudin1976}.
\end{theorem}

In particular, if \(\mathcal{A}\) contains the constant functions and separates points, then it is dense in \(C(K)\).

There are two main approaches to proving the universal approximation theorem:
\begin{itemize}
	\item \textbf{Approach of \citet{cybenko1989}}: Uses the Hahn-Banach corollary together with the Riesz representation theorem. It proves that sigmoidal functions are ``discriminatory'' (they vanish on a measure only if the measure is zero). This is the proof we present below.
	\item \textbf{Approach of \citet{hornik1989}}: Uses the Stone-Weierstrass theorem, proving that neural networks form a subalgebra that separates points.
\end{itemize}
We include Stone-Weierstrass in the preliminaries for its historical importance and its connection with the alternative approach.

\paragraph{Statement and proof}

The universal approximation theorem is a fundamental theoretical justification for the use of neural networks to approximate functions, proved independently by \citet{cybenko1989} and \citet{hornik1989}:

\begin{definition}[Sigmoidal function]\label{def:funcion_sigmoidal}
	A function \(\phi: \mathbb{R} \to \mathbb{R}\) is \textit{sigmoidal} if it is continuous, bounded and satisfies:
	\[
		\lim_{t \to -\infty} \phi(t) = 0 \quad \text{and} \quad \lim_{t \to +\infty} \phi(t) = 1
	\]
	Classical examples include the logistic function \(\sigma(t) = \frac{1}{1+e^{-t}}\) and the rescaled hyperbolic tangent, both described in detail in \autoref{subsubsec:funciones_activacion_capas}.
\end{definition}

\begin{theorem}[Universal Approximation - Cybenko, 1989]\label{thm:aproximacion_universal}
	Let \(\phi: \mathbb{R} \to \mathbb{R}\) be a sigmoidal function. Let \(I_m = [0,1]^m\) be the \(m\)-dimensional unit hypercube. Then, the set of functions of the form:
	\[
		F(x) = \sum_{i=1}^{N} v_i \phi(w_i^T x + b_i)
	\]
	with \(N \in \mathbb{N}\), \(v_i, b_i \in \mathbb{R}\) and \(w_i \in \mathbb{R}^m\), is dense in \(C(I_m)\) with respect to the uniform norm.

	That is, for any \(f \in C(I_m)\) and any \(\epsilon > 0\), there exist parameters such that:
	\[
		\sup_{x \in I_m} |F(x) - f(x)| < \epsilon
	\]
\end{theorem}

\textbf{Why the unit hypercube.}
The restriction to \(I_m = [0,1]^m\) is not a real limitation:
\begin{itemize}
	\item For any compact set \(K \subset \mathbb{R}^m\), by the Heine-Borel theorem (\autoref{thm:heine_borel}), \(K\) is bounded and therefore contained in a hypercube \([-M, M]^m\). By means of an affine transformation the problem can be translated to \(I_m\), so the result extends to \(C(K)\) for any compact set \(K \subset \mathbb{R}^m\)
	\item Compactness is essential: it guarantees the uniform continuity of \(f\) and the existence of approximations
	\item In practice, data is always bounded, so this restriction is natural
\end{itemize}

\textbf{Extensions of the theorem.}
The theorem has been significantly extended since 1989:
\begin{itemize}
	\item \textbf{More general activation functions}: The theorem holds for any non-polynomial function \citep{leshno1993multilayer}
	\item \textbf{ReLU}: For the \gls{ReLU} function (see \autoref{subsubsec:funciones_activacion_capas}), the result of Leshno et al. implies that one hidden layer is sufficient with arbitrary width. However, with bounded width, sufficient depth is required to achieve universal approximation \citep{lu2017expressive,hanin2019universal}
	\item \textbf{Approximation in \(L^p\)}: Similar results hold for other norms \citep{hornik1991approximation}
	\item \textbf{Approximation rates}: Bounds are known on the number of neurons \(N\) needed as a function of \(\epsilon\) and the regularity of \(f\) \citep{barron1993universal}
\end{itemize}

\begin{proof}
	The proof follows the approach of \citet{cybenko1989} using functional analysis. We proceed in several steps:

	The strategy is indirect: instead of explicitly constructing the approximation, we will prove that no continuous function can ``escape'' being approximated. To this end, we use duality in functional analysis. If neural networks were not dense in \(C(I_m)\), there would exist a nonzero continuous linear functional that vanishes on all of them (Hahn-Banach corollary). Translated into measures via the Riesz theorem, this would mean that there exists a nonzero measure ``invisible'' to all neurons. Proving that such a measure cannot exist (that is, that \(\phi\) is \textit{discriminatory}) is the heart of the proof.

	\textbf{Step 1: Discriminatory functions.}
	We define the set of discriminatory sigmoidal functions. A function \(\phi\) is discriminatory if for any finite signed measure \(\mu\) on \(I_m\):
	\[
		\int_{I_m} \phi(w^T x + b) \, d\mu(x) = 0 \quad \forall w \in \mathbb{R}^m, b \in \mathbb{R} \implies \mu = 0
	\]

	In other words, the family of functions \(\{\phi(w^T x + b)\}_{w,b}\) is rich enough to ``detect'' any nonzero measure: there is no nontrivial mass distribution that is invisible to all neurons simultaneously. The term ``discriminatory'' reflects this ability to discriminate between the zero measure and any other measure, analogously to how test functions in distribution theory distinguish distinct distributions.

	\textbf{Step 2: Sigmoidal functions are discriminatory.}
	Let \(\phi\) be a sigmoidal function (continuous, bounded, \(\lim_{t \to -\infty} \phi(t) = 0\), \(\lim_{t \to +\infty} \phi(t) = 1\)). Suppose that \(\mu\) is a measure such that \(\int_{I_m} \phi(w^T x + b) \, d\mu(x) = 0\) for all \(w, b\). We must prove that \(\mu = 0\).

	For any \(w \in \mathbb{R}^m\), \(b \in \mathbb{R}\) and auxiliary parameter \(\alpha \in \mathbb{R}\), we define the family of functions:
	\[
		\psi_{\lambda}(x) = \phi(\lambda(w^T x + b) + \alpha)
	\]

	The idea is that by multiplying the argument of \(\phi\) by a large \(\lambda\), the smooth transition of the sigmoid is ``compressed'' until it becomes an abrupt jump at the hyperplane \(w^T x + b = 0\). Geometrically, each hyperplane divides \(I_m\) into two half-spaces, and the limit \(\lambda \to \infty\) turns the neuron into an indicator of the positive half-space.

	When \(\lambda \to \infty\), we have pointwise convergence to a step function:
	\[
		\psi_{\lambda}(x) \to \begin{cases}
			                        1 & \text{if } w^T x + b > 0 \\
			                        0 & \text{if } w^T x + b < 0 \\
			                        \phi(\alpha) & \text{if } w^T x + b = 0
		\end{cases}
	\]

	Since \(|\psi_{\lambda}(x)| \leq \sup_{t \in \mathbb{R}} |\phi(t)| < \infty\) uniformly in \(\lambda\) and \(x\), and this bound is integrable over the compact set \(I_m\), we can apply the \textbf{Lebesgue dominated convergence theorem}:
	\[
		\int_{I_m} \psi_{\lambda}(x) \, d\mu(x) \to \mu(H^+_{w,b}) + \phi(\alpha) \cdot \mu(H_{w,b})
	\]
	where \(H^+_{w,b} = \{x \in I_m : w^T x + b > 0\}\) is the open half-space and \(H_{w,b} = \{x \in I_m : w^T x + b = 0\}\) is the hyperplane.

	By hypothesis, the left-hand side is zero for all \(\lambda\). By continuity of the limit, we obtain:
	\[
		\mu(H^+_{w,b}) + \phi(\alpha) \cdot \mu(H_{w,b}) = 0 \quad \text{for all } \alpha \in \mathbb{R}
	\]

	To see why, observe that the equation \(\mu(H^+_{w,b}) + \phi(\alpha) \cdot \mu(H_{w,b}) = 0\) is an affine function of \(\phi(\alpha)\). If \(\mu(H_{w,b}) \neq 0\), this equation would determine a unique value of \(\phi(\alpha)\); but since \(\phi\) takes a continuous range of values (from 0 to 1), the equation could not be satisfied for all \(\alpha\). Thus, necessarily \(\mu(H_{w,b}) = 0\) and, substituting, \(\mu(H^+_{w,b}) = 0\).

	Since this equation holds for all \(\alpha\) and \(\phi\) is not constant (since \(\phi(\alpha) \to 0\) when \(\alpha \to -\infty\) and \(\phi(\alpha) \to 1\) when \(\alpha \to +\infty\)), necessarily \(\mu(H_{w,b}) = 0\) and \(\mu(H^+_{w,b}) = 0\).

	This holds for all \(w\) and \(b\), so \(\mu\) vanishes on all open half-spaces. Since the half-spaces generate the Borel \(\sigma\)-algebra and finite signed measures are determined by their values on a generating class stable under intersections, we conclude that \(\mu = 0\).

	\textbf{Step 3: Density in \(C(I_m)\).}
	Let \(\mathcal{S}\) be the set of functions of the form:
	\[
		G(x) = \sum_{i=1}^{N} v_i \phi(w_i^T x + b_i)
	\]

	The connection between the discriminatory property and density is established via two classical theorems of functional analysis. The Hahn-Banach corollary (\autoref{cor:caracterizacion_densidad}) reduces the density problem to verifying that no nonzero functional vanishes on \(\mathcal{S}\). The Riesz theorem (\autoref{thm:riesz_representacion}) translates functionals into measures, bringing the problem exactly to the condition in the definition of a discriminatory function.

	By Step 2, \(\phi\) is discriminatory. We apply the \textbf{corollary of the Hahn-Banach theorem}: the set \(\mathcal{S}\) is dense in \(C(I_m)\) if and only if the only continuous linear functional that vanishes on \(\mathcal{S}\) is the zero functional.

	Let \(L\) be a bounded linear functional on \(C(I_m)\) such that \(L(G) = 0\) for all \(G \in \mathcal{S}\). By the \textbf{Riesz representation theorem for \(C(K)\)}, there exists a unique finite signed measure \(\mu\) such that:
	\[
		L(h) = \int_{I_m} h(x) \, d\mu(x) \quad \forall h \in C(I_m)
	\]

	Since \(L(G) = 0\) for all \(G \in \mathcal{S}\), in particular:
	\[
		\int_{I_m} \phi(w^T x + b) \, d\mu(x) = 0 \quad \forall w \in \mathbb{R}^m, b \in \mathbb{R}
	\]

	By the definition of a discriminatory function and Step 2, this implies \(\mu = 0\), hence \(L = 0\).

	\textbf{Step 4: Uniform approximation.}
	By the previous step, \(\mathcal{S}\) is dense in \(C(I_m)\) with respect to the supremum norm. By the definition of density, for any \(f \in C(I_m)\) and \(\epsilon > 0\), there exists \(F \in \mathcal{S}\) such that:
	\[
		\|F - f\|_\infty = \sup_{x \in I_m} |F(x) - f(x)| < \epsilon
	\]

	This completes the proof. \(\square\)
\end{proof}

This theorem guarantees that a neural network with a single hidden layer can approximate any continuous function with arbitrary precision. However, it is a result of \textbf{existence}, not of \textbf{construction}: it proves that suitable weights exist, but it does not say how to find them. Moreover, the number of neurons needed can be prohibitively large, and deep networks (multiple layers) can represent complex functions more efficiently, allowing hierarchical representations of the data to be learned.

\subsection{Convergence theory}\label{subsec:teoria_convergencia}

The universal approximation theorem guarantees the existence of a weight configuration that approximates any target function. In \autoref{subsec:descenso_gradiente} we presented gradient descent and backpropagation as the mechanism for searching for those weights. We now formalize the conditions under which this process effectively converges and establish its convergence rates. For an in-depth treatment, see \citet{bottou2018optimization}.

\subsubsection{Regularity conditions}

The convergence theorems for gradient descent do not hold for arbitrary loss functions: they require certain regularity hypotheses. In this section we formalize the two fundamental conditions: \textit{$L$-smoothness}, which controls how much the gradient can vary and allows bounding the progress per iteration, and \textit{convexity}, which guarantees the absence of spurious local minima. These properties provide the analytic framework needed for the convergence results that we present below.

\begin{definition}[$L$-smooth function]\label{def:L_smooth}
	A differentiable function \(f: \mathbb{R}^n \to \mathbb{R}\) is \textit{\(L\)-smooth} (or has \(L\)-Lipschitz gradient) if there exists a constant \(L > 0\) such that:
	\[
		\|\nabla f(x) - \nabla f(y)\| \leq L\|x - y\| \quad \forall x, y \in \mathbb{R}^n
	\]
\end{definition}

This condition is analogous to Lipschitz continuity (\autoref{prop:lipschitz_continuidad}), but applied to the gradient rather than to the function itself.

\begin{proposition}[Equivalent characterization of $L$-smoothness]\label{prop:L_smooth_equiv}
	Let \(f: \mathbb{R}^n \to \mathbb{R}\) be differentiable. The following statements are equivalent:
	\begin{enumerate}
		\item \(f\) is \(L\)-smooth (\(L\)-Lipschitz gradient)
		\item For all \(x, y \in \mathbb{R}^n\): \(f(y) \leq f(x) + \langle \nabla f(x), y - x \rangle + \frac{L}{2}\|y - x\|^2\)
	\end{enumerate}
\end{proposition}

\begin{proof}
	\textbf{(1) \(\Rightarrow\) (2):} We define \(g(t) = f(x + t(y-x))\) for \(t \in [0,1]\). By the fundamental theorem of calculus:
	\[
		f(y) - f(x) = g(1) - g(0) = \int_0^1 g'(t) \, dt = \int_0^1 \langle \nabla f(x + t(y-x)), y-x \rangle \, dt
	\]
	Subtracting \(\langle \nabla f(x), y-x \rangle\):
	\begin{align*}
		f(y) - f(x) - \langle \nabla f(x), y-x \rangle &= \int_0^1 \langle \nabla f(x + t(y-x)) - \nabla f(x), y-x \rangle \, dt \\
		&\leq \int_0^1 \|\nabla f(x + t(y-x)) - \nabla f(x)\| \cdot \|y-x\| \, dt \\
		&\leq \int_0^1 L \cdot t\|y-x\| \cdot \|y-x\| \, dt = \frac{L}{2}\|y-x\|^2
	\end{align*}

	\textbf{(2) \(\Rightarrow\) (1):} Applying (2) twice (interchanging \(x\) and \(y\)) and adding:
	\begin{align*}
		f(y) &\leq f(x) + \langle \nabla f(x), y-x \rangle + \frac{L}{2}\|y-x\|^2 \\
		f(x) &\leq f(y) + \langle \nabla f(y), x-y \rangle + \frac{L}{2}\|x-y\|^2
	\end{align*}
	Adding: \(0 \leq \langle \nabla f(x) - \nabla f(y), y-x \rangle + L\|y-x\|^2\), that is:
	\[
		\langle \nabla f(x) - \nabla f(y), x-y \rangle \leq L\|x-y\|^2
	\]
	By Cauchy-Schwarz: \(\|\nabla f(x) - \nabla f(y)\| \cdot \|x-y\| \geq \langle \nabla f(x) - \nabla f(y), x-y \rangle\), but this inequality alone is not enough. To complete it rigorously, it is required that \(f\) be convex: in that case, the above inequality implies co-coercivity of the gradient, which in turn yields the Lipschitz condition. See \citet{Nesterov2018}, Theorem~2.1.5, for the details.
\end{proof}

\begin{remark}
	The direction (1) \(\Rightarrow\) (2) holds for every differentiable function, without the need for convexity. The direction (2) \(\Rightarrow\) (1), by contrast, requires that \(f\) be convex \citep[Theorem~2.1.5]{Nesterov2018}. Since in the subsequent convergence theorems (\autoref{thm:gd_convergence_convex}, \autoref{thm:gd_nonconvex}) only the direction (1) \(\Rightarrow\) (2) is used, this restriction does not affect the results in the rest of the chapter.
\end{remark}

The constant \(L\) controls how much the gradient can vary locally. Intuitively, functions with small \(L\) are ``flatter'' and allow larger optimization steps.

The rate \(\eta = 1/L\) is optimal in the following sense. The \(L\)-smoothness bound
(\autoref{prop:L_smooth_equiv}) guarantees that for every pair of points:
\[
	f(y) \leq f(x) + \langle \nabla f(x), y - x \rangle + \frac{L}{2}\|y - x\|^2
\]
Substituting the gradient descent step with generic rate \(\eta > 0\),
that is, \(x = \theta_t\) and \(y = \theta_{t+1} = \theta_t - \eta\nabla f(\theta_t)\),
so that \(y - x = -\eta\nabla f(\theta_t)\):
\begin{align*}
	f(\theta_{t+1})
		&\leq f(\theta_t) + \langle \nabla f(\theta_t),\, -\eta\nabla f(\theta_t) \rangle
			+ \frac{L}{2}\|\!-\eta\nabla f(\theta_t)\|^2 \\
		&= f(\theta_t) - \eta\|\nabla f(\theta_t)\|^2
			+ \frac{L\eta^2}{2}\|\nabla f(\theta_t)\|^2 \\
		&= f(\theta_t)
			- \underbrace{\left(\eta - \frac{L\eta^2}{2}\right)}_{g(\eta)}
			\|\nabla f(\theta_t)\|^2
\end{align*}
The coefficient \(g(\eta) = \eta - \frac{L\eta^2}{2}\) is a concave parabola that attains its maximum when \(g'(\eta) = 1 - L\eta = 0\), that is, \(\eta^* = 1/L\), with value \(g(1/L) = 1/(2L)\). Moreover, \(g(\eta) > 0\) requires \(\eta < 2/L\): for rates \(\eta \geq 2/L\) the bound no longer guarantees descent and the algorithm may diverge. For \(0 < \eta < 1/L\) the algorithm converges, but with suboptimal guaranteed decrease.

\begin{definition}[Convex function]\label{def:convex_function}
	A function \(f: \mathbb{R}^n \to \mathbb{R}\) is \textit{convex} if for all \(x, y \in \mathbb{R}^n\) and \(\lambda \in [0,1]\):
	\[
		f(\lambda x + (1-\lambda)y) \leq \lambda f(x) + (1-\lambda)f(y)
	\]
\end{definition}

This definition is intimately related to the notion of a convex set (\autoref{def:conjunto_convexo}): a function is convex if and only if its epigraph $\{(x,y) \in \mathbb{R}^{n+1} : y \geq f(x)\}$ is a convex set.

\begin{proposition}[Characterization of convexity for differentiable functions]\label{prop:convex_equiv}
	Let \(f: \mathbb{R}^n \to \mathbb{R}\) be differentiable. Then \(f\) is convex if and only if:
	\[
		f(y) \geq f(x) + \langle \nabla f(x), y - x \rangle \quad \forall x, y \in \mathbb{R}^n
	\]
	That is, the graph of \(f\) lies above all of its tangent hyperplanes.
\end{proposition}

\begin{proof}
	\textbf{(\(\Rightarrow\)):} Suppose \(f\) is convex. For any \(x, y\) and \(\lambda \in (0,1]\):
	\[
		f(x + \lambda(y-x)) \leq (1-\lambda)f(x) + \lambda f(y)
	\]
	Rearranging: \(\frac{f(x + \lambda(y-x)) - f(x)}{\lambda} \leq f(y) - f(x)\).

	Taking \(\lambda \to 0^+\), the left-hand side converges to the directional derivative \(\langle \nabla f(x), y-x \rangle\), obtaining \(\langle \nabla f(x), y-x \rangle \leq f(y) - f(x)\).

	\textbf{(\(\Leftarrow\)):} Suppose the gradient condition holds. Let \(x, y \in \mathbb{R}^n\) and \(\lambda \in [0,1]\). We define \(z = \lambda x + (1-\lambda)y\). Applying the hypothesis:
	\begin{align*}
		f(x) &\geq f(z) + \langle \nabla f(z), x - z \rangle = f(z) + (1-\lambda)\langle \nabla f(z), x - y \rangle \\
		f(y) &\geq f(z) + \langle \nabla f(z), y - z \rangle = f(z) - \lambda\langle \nabla f(z), x - y \rangle
	\end{align*}
	Multiplying the first by \(\lambda\) and the second by \((1-\lambda)\) and adding:
	\[
		\lambda f(x) + (1-\lambda)f(y) \geq f(z) = f(\lambda x + (1-\lambda)y)
	\]
\end{proof}

\subsubsection{Convergence of gradient descent}

\begin{theorem}[Convergence of \gls{GD} for convex functions]\label{thm:gd_convergence_convex}
	Let \(f: \mathbb{R}^n \to \mathbb{R}\) be convex and \(L\)-smooth with global minimum \(f^* = f(\theta^*)\). If we apply gradient descent with learning rate \(\eta = 1/L\), then:
	\[
		f(\theta_T) - f^* \leq \frac{L\|\theta_0 - \theta^*\|^2}{2T}
	\]
	That is, gradient descent converges with rate \(O(1/T)\) to the optimum.
\end{theorem}

\begin{proof}
	We proceed in several steps.

	\textbf{Step 1: Descent inequality.}
	By \autoref{prop:L_smooth_equiv}, for any pair of points:
	\[
		f(y) \leq f(x) + \langle \nabla f(x), y - x \rangle + \frac{L}{2}\|y - x\|^2
	\]
	Applying this inequality with \(x = \theta_t\) and \(y = \theta_{t+1} = \theta_t - \frac{1}{L}\nabla f(\theta_t)\):
	\begin{align*}
		f(\theta_{t+1}) &\leq f(\theta_t) + \left\langle \nabla f(\theta_t), -\frac{1}{L}\nabla f(\theta_t) \right\rangle + \frac{L}{2}\left\|-\frac{1}{L}\nabla f(\theta_t)\right\|^2 \\
		&= f(\theta_t) - \frac{1}{L}\|\nabla f(\theta_t)\|^2 + \frac{1}{2L}\|\nabla f(\theta_t)\|^2 \\
		&= f(\theta_t) - \frac{1}{2L}\|\nabla f(\theta_t)\|^2
	\end{align*}
	This inequality guarantees that the function value decreases at each iteration.

	\textbf{Step 2: Bound on the distance to the optimum.}
	We analyze how the distance \(\|\theta_{t+1} - \theta^*\|^2\) evolves:
	\begin{align*}
		\|\theta_{t+1} - \theta^*\|^2 &= \left\|\theta_t - \frac{1}{L}\nabla f(\theta_t) - \theta^*\right\|^2 \\
		&= \|\theta_t - \theta^*\|^2 - \frac{2}{L}\langle \nabla f(\theta_t), \theta_t - \theta^* \rangle + \frac{1}{L^2}\|\nabla f(\theta_t)\|^2
	\end{align*}

	\textbf{Step 3: Use of convexity.}
	By \autoref{prop:convex_equiv}:
	\[
		f(\theta^*) \geq f(\theta_t) + \langle \nabla f(\theta_t), \theta^* - \theta_t \rangle
	\]
	Rearranging:
	\[
		\langle \nabla f(\theta_t), \theta_t - \theta^* \rangle \geq f(\theta_t) - f^*
	\]

	\textbf{Step 4: Combination of the bounds.}
	Substituting the inequality from Step 3 into the one from Step 2:
	\[
		\|\theta_{t+1} - \theta^*\|^2 \leq \|\theta_t - \theta^*\|^2 - \frac{2}{L}(f(\theta_t) - f^*) + \frac{1}{L^2}\|\nabla f(\theta_t)\|^2
	\]

	From Step 1, we have \(\|\nabla f(\theta_t)\|^2 \leq 2L(f(\theta_t) - f(\theta_{t+1}))\). Substituting:
	\begin{align*}
		\|\theta_{t+1} - \theta^*\|^2 &\leq \|\theta_t - \theta^*\|^2 - \frac{2}{L}(f(\theta_t) - f^*) + \frac{2}{L}(f(\theta_t) - f(\theta_{t+1})) \\
		&= \|\theta_t - \theta^*\|^2 - \frac{2}{L}(f(\theta_{t+1}) - f^*)
	\end{align*}

	Rearranging:
	\begin{equation}\label{eq:gd_telescopic}
		f(\theta_{t+1}) - f^* \leq \frac{L}{2}\left(\|\theta_t - \theta^*\|^2 - \|\theta_{t+1} - \theta^*\|^2\right)
	\end{equation}

	\textbf{Step 5: Telescoping sum.}
	Summing inequality~\eqref{eq:gd_telescopic} for \(t = 0, 1, \ldots, T-1\):
	\[
		\sum_{t=0}^{T-1} (f(\theta_{t+1}) - f^*) \leq \frac{L}{2} \sum_{t=0}^{T-1} \left(\|\theta_t - \theta^*\|^2 - \|\theta_{t+1} - \theta^*\|^2\right)
	\]

	The right-hand side is a telescoping sum:
	\[
		\sum_{t=0}^{T-1} (f(\theta_{t+1}) - f^*) \leq \frac{L}{2}\left(\|\theta_0 - \theta^*\|^2 - \|\theta_T - \theta^*\|^2\right) \leq \frac{L}{2}\|\theta_0 - \theta^*\|^2
	\]

	\textbf{Step 6: Conclusion.}
	By Step 1, the sequence \(\{f(\theta_t)\}\) is decreasing, so \(f(\theta_T) \leq f(\theta_{t+1})\) for all \(t < T\). Then:
	\[
		T(f(\theta_T) - f^*) \leq \sum_{t=0}^{T-1} (f(\theta_{t+1}) - f^*) \leq \frac{L}{2}\|\theta_0 - \theta^*\|^2
	\]

	Dividing by \(T\):
	\[
		f(\theta_T) - f^* \leq \frac{L\|\theta_0 - \theta^*\|^2}{2T}
	\]

	This proves convergence with rate \(O(1/T)\).
\end{proof}

\begin{theorem}[Convergence of SGD]\label{thm:sgd_convergence}
	Let \(f\) be convex and \(L\)-smooth, and let \(g_t\) be an unbiased estimator of the gradient with bounded variance:
	\[
		\mathbb{E}[g_t] = \nabla f(\theta_t), \quad \mathbb{E}[\|g_t - \nabla f(\theta_t)\|^2] \leq \sigma^2
	\]
	Suppose moreover that the gradients are uniformly bounded: \(\|\nabla f(\theta_t)\| \leq G\) for all \(t\). With decreasing learning rate \(\eta_t = \eta_0/\sqrt{t}\), SGD satisfies:
	\[
		\mathbb{E}[f(\bar{\theta}_T)] - f^* = O\!\left(\frac{(\sigma^2 + G^2)\log T}{\sqrt{T}} + \frac{\|\theta_0 - \theta^*\|^2}{\eta_0\sqrt{T}}\right)
	\]
	where \(\bar{\theta}_T = \frac{1}{T}\sum_{t=1}^T \theta_t\) is the average of the iterates.
\end{theorem}

\begin{proof}
	We present a sketch of the proof following~\citet{bottou2018optimization}.

	\textbf{Step 1: Decomposition of the stochastic gradient.}
	We write \(g_t = \nabla f(\theta_t) + \xi_t\) where \(\xi_t = g_t - \nabla f(\theta_t)\) is the stochastic noise with:
	\[
		\mathbb{E}[\xi_t | \theta_t] = 0, \quad \mathbb{E}[\|\xi_t\|^2 | \theta_t] \leq \sigma^2
	\]

	\textbf{Step 2: Analysis of the expected progress.}
	The update is \(\theta_{t+1} = \theta_t - \eta_t g_t\). By convexity:
	\[
		f(\theta^*) \geq f(\theta_t) + \langle \nabla f(\theta_t), \theta^* - \theta_t \rangle
	\]
	Rearranging: \(\langle \nabla f(\theta_t), \theta_t - \theta^* \rangle \geq f(\theta_t) - f^*\).

	We analyze the distance to the optimum:
	\begin{align*}
		\|\theta_{t+1} - \theta^*\|^2 &= \|\theta_t - \eta_t g_t - \theta^*\|^2 \\
		&= \|\theta_t - \theta^*\|^2 - 2\eta_t\langle g_t, \theta_t - \theta^* \rangle + \eta_t^2\|g_t\|^2
	\end{align*}

	Taking conditional expectation and using that \(\mathbb{E}[g_t|\theta_t] = \nabla f(\theta_t)\):
	\[
		\mathbb{E}[\|\theta_{t+1} - \theta^*\|^2 | \theta_t] \leq \|\theta_t - \theta^*\|^2 - 2\eta_t(f(\theta_t) - f^*) + \eta_t^2(\|\nabla f(\theta_t)\|^2 + \sigma^2)
	\]

	\textbf{Step 3: Telescoping sum and averaging.}
	We rearrange the inequality from Step~2 to isolate the progress term:
	\[
		2\eta_t(f(\theta_t) - f^*) \leq \|\theta_t - \theta^*\|^2 - \mathbb{E}[\|\theta_{t+1} - \theta^*\|^2 | \theta_t] + \eta_t^2(\|\nabla f(\theta_t)\|^2 + \sigma^2)
	\]
	Taking total expectation and summing over \(t = 1, \ldots, T\):
	\[
		\sum_{t=1}^T 2\eta_t \mathbb{E}[f(\theta_t) - f^*] \leq \sum_{t=1}^T \Big(\mathbb{E}[\|\theta_t - \theta^*\|^2] - \mathbb{E}[\|\theta_{t+1} - \theta^*\|^2]\Big) + \sum_{t=1}^T \eta_t^2(\|\nabla f(\theta_t)\|^2 + \sigma^2)
	\]
	The first sum is telescoping: consecutive terms cancel, leaving only the endpoints,
	\[
		\sum_{t=1}^T \Big(\mathbb{E}[\|\theta_t - \theta^*\|^2] - \mathbb{E}[\|\theta_{t+1} - \theta^*\|^2]\Big) = \mathbb{E}[\|\theta_1 - \theta^*\|^2] - \mathbb{E}[\|\theta_{T+1} - \theta^*\|^2] \leq \|\theta_0 - \theta^*\|^2
	\]
	where the last inequality uses that \(\mathbb{E}[\|\theta_1 - \theta^*\|^2] \leq \|\theta_0 - \theta^*\|^2\) (by the result of Step~2 with \(t=0\) and \(f(\theta_0) \geq f^*\)) and that \(-\mathbb{E}[\|\theta_{T+1} - \theta^*\|^2] \leq 0\). Substituting:
	\[
		\sum_{t=1}^T 2\eta_t \mathbb{E}[f(\theta_t) - f^*] \leq \|\theta_0 - \theta^*\|^2 + \sum_{t=1}^T \eta_t^2(\|\nabla f(\theta_t)\|^2 + \sigma^2)
	\]

	\textbf{Step 4: Choice of decreasing rate.}
	We choose \(\eta_t = \eta_0/\sqrt{t}\). Using the hypothesis \(\|\nabla f(\theta_t)\| \leq G\) and substituting into the inequality from Step~3:
	\[
		\sum_{t=1}^T \frac{2\eta_0}{\sqrt{t}}\, \mathbb{E}[f(\theta_t) - f^*] \leq \|\theta_0 - \theta^*\|^2 + \sum_{t=1}^T \frac{\eta_0^2}{t}(G^2 + \sigma^2)
	\]

	To bound the right-hand side, we use that the harmonic series satisfies \(\sum_{t=1}^T \frac{1}{t} \leq 1 + \log T\):
	\[
		\sum_{t=1}^T \frac{\eta_0^2}{t}(G^2 + \sigma^2) \leq \eta_0^2(G^2 + \sigma^2)(1 + \log T)
	\]

	To bound the left-hand side from below, we observe that \(\eta_t = \eta_0/\sqrt{t}\) is decreasing, so that \(\eta_t \geq \eta_T = \eta_0/\sqrt{T}\) for all \(t \leq T\). Since each term \(\mathbb{E}[f(\theta_t) - f^*] \geq 0\) by convexity:
	\[
		\sum_{t=1}^T \frac{2\eta_0}{\sqrt{t}}\, \mathbb{E}[f(\theta_t) - f^*] \geq \frac{2\eta_0}{\sqrt{T}} \sum_{t=1}^T \mathbb{E}[f(\theta_t) - f^*]
	\]

	Combining both bounds and dividing by \(T\):
	\[
		\frac{1}{T}\sum_{t=1}^T \mathbb{E}[f(\theta_t) - f^*] \leq \frac{\|\theta_0 - \theta^*\|^2 + \eta_0^2(G^2 + \sigma^2)(1 + \log T)}{2\eta_0\sqrt{T}}
	\]

	Finally, by convexity of \(f\) and Jensen's inequality applied to the average \(\bar{\theta}_T = \frac{1}{T}\sum_{t=1}^T \theta_t\):
	\[
		\mathbb{E}[f(\bar{\theta}_T)] \leq \frac{1}{T}\sum_{t=1}^T \mathbb{E}[f(\theta_t)]
	\]
	which, combined with the previous bound, gives:
	\[
		\mathbb{E}[f(\bar{\theta}_T)] - f^* \leq \frac{\|\theta_0 - \theta^*\|^2 + \eta_0^2(\sigma^2 + G^2)(1 + \log T)}{2\eta_0\sqrt{T}}
	\]
	This establishes the rate \(O(\log(T)/\sqrt{T})\).
\end{proof}

The \(O(\log(T)/\sqrt{T})\) rate of \gls{SGD} is slower than the \(O(1/T)\) rate of deterministic \gls{GD}, but \gls{SGD} processes \(T\) times less data per iteration, which can result in better total computational efficiency.

\begin{theorem}[Convergence to stationary points (non-convex case)]\label{thm:gd_nonconvex}
	Let \(f: \mathbb{R}^n \to \mathbb{R}\) be an \(L\)-smooth function (not necessarily convex) and bounded below by \(f^*\). With learning rate \(\eta \leq 1/L\), gradient descent satisfies:
	\[
		\min_{t=0,\ldots,T-1} \|\nabla f(\theta_t)\|^2 \leq \frac{2(f(\theta_0) - f^*)}{\eta T}
	\]
	That is, \gls{GD} finds an \(\epsilon\)-stationary point (with \(\|\nabla f\| \leq \epsilon\)) in \(O(1/\epsilon^2)\) iterations.
\end{theorem}

\begin{proof}
	We use the characterization of \(L\)-smoothness from \autoref{def:L_smooth}.

	\textbf{Step 1: Descent inequality.}
	By \(L\)-smoothness, for the update \(\theta_{t+1} = \theta_t - \eta \nabla f(\theta_t)\):
	\begin{align*}
		f(\theta_{t+1}) &\leq f(\theta_t) + \langle \nabla f(\theta_t), \theta_{t+1} - \theta_t \rangle + \frac{L}{2}\|\theta_{t+1} - \theta_t\|^2 \\
		&= f(\theta_t) - \eta \|\nabla f(\theta_t)\|^2 + \frac{L\eta^2}{2}\|\nabla f(\theta_t)\|^2 \\
		&= f(\theta_t) - \eta\left(1 - \frac{L\eta}{2}\right)\|\nabla f(\theta_t)\|^2
	\end{align*}
	If \(\eta \leq 1/L\), then \(1 - \frac{L\eta}{2} \geq \frac{1}{2}\), so:
	\[
		f(\theta_{t+1}) \leq f(\theta_t) - \frac{\eta}{2}\|\nabla f(\theta_t)\|^2
	\]

	\textbf{Step 2: Telescoping sum.}
	Summing over \(t = 0, \ldots, T-1\):
	\[
		\sum_{t=0}^{T-1} \frac{\eta}{2}\|\nabla f(\theta_t)\|^2 \leq \sum_{t=0}^{T-1} \left(f(\theta_t) - f(\theta_{t+1})\right) = f(\theta_0) - f(\theta_T) \leq f(\theta_0) - f^*
	\]
	where we use that \(f(\theta_T) \geq f^*\) by the lower bound.

	\textbf{Step 3: Conclusion.}
	Dividing both sides by \(\eta/2\):
	\[
		\sum_{t=0}^{T-1} \|\nabla f(\theta_t)\|^2 \leq \frac{2(f(\theta_0) - f^*)}{\eta}
	\]
	Since each term \(\|\nabla f(\theta_t)\|^2 \geq 0\), the minimum value of the collection cannot exceed its arithmetic mean (otherwise, all terms would exceed the mean, and their sum would exceed \(T\) times the mean, contradicting its definition):
	\[
		\min_{t=0,\ldots,T-1} \|\nabla f(\theta_t)\|^2 \leq \frac{1}{T}\sum_{t=0}^{T-1} \|\nabla f(\theta_t)\|^2
	\]
	Combining with the previous bound, divided by \(T\):
	\[
		\min_{t=0,\ldots,T-1} \|\nabla f(\theta_t)\|^2 \leq \frac{1}{T}\sum_{t=0}^{T-1} \|\nabla f(\theta_t)\|^2 \leq \frac{2(f(\theta_0) - f^*)}{\eta T}
	\]
\end{proof}

The loss functions of neural networks are typically non-convex due to the nonlinearities. \autoref{thm:gd_nonconvex} guarantees convergence to stationary points, but not to global minima. However, in practice the overparameterization of modern networks creates many equivalent minima, the noise of \gls{SGD} helps to escape saddle points, and the local minima found tend to generalize well.

\subsection{Advanced optimizers}\label{subsec:optimizadores_avanzados}

Modern optimizers improve upon basic \gls{SGD} by incorporating historical information about the gradients~\citep{kingma2014adam, duchi2011adagrad}. In this section, all operations on parameter and gradient vectors are applied element-wise: we use \(\odot\) for the Hadamard product and \(\oslash\) for element-wise division, while the square-root and scalar-addition operations are also applied component-wise.

\subsubsection{SGD with Momentum}

\begin{definition}[SGD with Momentum]\label{def:sgd_momentum}
	\textit{SGD with momentum}~\citep{polyak1964methods} introduces a velocity variable \(v_t\) that accumulates past gradients, with \(v_0 = 0\):
	\begin{align*}
		v_t &= \beta v_{t-1} + \nabla_\theta \mathcal{L}(\theta_{t-1}) \\
		\theta_t &= \theta_{t-1} - \eta v_t
	\end{align*}
	where \(\beta \in [0,1)\) is the \textit{momentum coefficient} (typically \(\beta = 0.9\)).
\end{definition}

Momentum accumulates exponentially weighted past gradients, smoothing the oscillations in directions of high curvature and accelerating convergence in consistent directions. Note that the momentum formula is \textit{not} an exponential moving average (EMA), since it lacks the factor \((1-\beta)\): whereas an EMA normalizes so that the weights sum to one, the momentum accumulator grows by a factor \(1/(1-\beta)\) in steady state, and it is the learning rate \(\eta\) that absorbs this scale.

\subsubsection{Adaptive methods}

\begin{definition}[Adagrad]\label{def:adagrad}
	The \textit{Adagrad}~\citep{duchi2011adagrad} optimizer adapts the learning rate per parameter according to the history of gradients, with \(G_0 = 0\):
	\begin{align*}
		G_t &= G_{t-1} + g_t \odot g_t \\
		\theta_t &= \theta_{t-1} - \eta \cdot g_t \oslash (\sqrt{G_t} + \epsilon)
	\end{align*}
	where \(g_t = \nabla_\theta \mathcal{L}(\theta_{t-1})\) and \(\epsilon \approx 10^{-8}\) avoids division by zero. Parameters with historically large gradients receive smaller learning rates.
\end{definition}

\begin{definition}[RMSProp]\label{def:rmsprop}
	\textit{RMSProp}~\citep{hinton2012rmsprop} corrects the monotonic decay of Adagrad by using an exponential moving average of the square of the gradients. We denote by \(v_t\) the accumulator of the second moment (reusing the symbol for the velocity in momentum, following the usual convention in the literature), with \(v_0 = 0\):
	\begin{align*}
		v_t &= \beta v_{t-1} + (1-\beta) (g_t \odot g_t) \\
		\theta_t &= \theta_{t-1} - \eta \cdot g_t \oslash (\sqrt{v_t} + \epsilon)
	\end{align*}
	where typically \(\beta = 0.9\). Unlike Adagrad, the denominator does not grow indefinitely, allowing sustained learning in non-stationary problems.
\end{definition}

\begin{definition}[Adam]\label{def:adam}
	The \textit{Adam}~\citep{kingma2014adam} optimizer (Adaptive Moment Estimation) combines momentum with adaptive rates, maintaining estimates of the first two moments of the gradient, with \(m_0 = v_0 = 0\):
	\begin{align*}
		m_t &= \beta_1 m_{t-1} + (1-\beta_1) g_t & \text{(first moment)} \\
		v_t &= \beta_2 v_{t-1} + (1-\beta_2) (g_t \odot g_t) & \text{(second moment)} \\
		\hat{m}_t &= m_t / (1 - \beta_1^t) & \text{(bias correction)} \\
		\hat{v}_t &= v_t / (1 - \beta_2^t) & \\
		\theta_t &= \theta_{t-1} - \eta \cdot \hat{m}_t \oslash (\sqrt{\hat{v}_t} + \epsilon)
	\end{align*}
	where typically \(\beta_1 = 0.9\), \(\beta_2 = 0.999\), \(\epsilon = 10^{-8}\).
\end{definition}

The corrections \(\hat{m}_t\) and \(\hat{v}_t\) compensate for the bias toward zero of the exponential moving averages when they are initialized at zero; for large \(t\), \((1 - \beta^t) \approx 1\) and the correction is negligible. It should be clarified that the terms ``first moment'' and ``second moment'' in \gls{Adam} refer to \textit{exponential moving averages} of the gradient and of its square, not to probabilistic moments in the strict sense. The terminology comes from the analogy with the estimation of the mean (\(\mathbb{E}[g_t]\)) and the uncentered variance (\(\mathbb{E}[g_t^2]\)) of the stochastic gradient, but \(m_t\) and \(v_t\) are biased estimators that require correction.

The effective update in \gls{Adam} for each parameter \(\theta_i\) is approximately \(\Delta\theta_i \approx \eta \cdot \hat{m}_{t,i} / (\sqrt{\hat{v}_{t,i}} + \epsilon)\), which provides smaller steps for parameters with historically large gradients and relatively larger steps for parameters with small but consistent gradients.

\section{Deep Learning}

\subsection{From the multilayer perceptron to deep learning}

Interest in neural models began when Rosenblatt showed that the perceptron could learn to classify patterns from data~\citep{rosenblatt1958perceptron}. However, Minsky and Papert~\citep{minsky1969perceptrons} proved that the simple perceptron can only represent linearly separable functions, which revealed fundamental limitations of the model and triggered a prolonged period of reduced interest and funding in this field.

The introduction of the backpropagation algorithm~\citep{backpropagation} provided an efficient method for training multilayer networks via gradient descent. Nevertheless, this advance proved insufficient on its own: the computational barriers and the problems inherent in training networks with many layers persisted for years.

In the previous sections we have seen mathematical tools for training neural networks. However, for decades deep networks remained impractical due to the vanishing gradient problem analyzed in \autoref{subsubsec:vanishing_gradient}: as layers are added, the product of derivatives decreases exponentially, preventing the first layers from learning. The alternative of increasing the width of the network (see \autoref{thm:aproximacion_universal}) was computationally too costly to truly see the usefulness of \glspl{ANN} in moderately complex problems.

\gls{DL}, or deep learning, emerged as the convergence of multiple advances that overcame these limitations:
\begin{itemize}
	\item Initialization techniques that preserve the variance of the activations across the layers.
	\item Activation functions such as \gls{ReLU} (see \autoref{subsubsec:funciones_activacion_capas}) that avoid saturation and maintain nonzero gradients.
	\item Regularization methods such as Dropout (\autoref{def:dropout}) that control overfitting in high-capacity networks.
	\item Specialized architectures that drastically reduce the computational cost by exploiting the structure of the data. For example, \glspl{CNN} exploit spatial locality to require orders of magnitude fewer parameters than equivalent dense networks.
\end{itemize}

The crystallization of these advances and their practical impact are documented in \autoref{ch:cnn}, where it is analyzed how convolutional networks materialized the deep learning revolution.

These deeper networks have greater expressive capacity, but they require more data and computational resources. A common technique to compensate for the scarcity of data is \textit{data augmentation}, which artificially increases the training set by applying transformations that preserve the labels.

\begin{definition}[Data augmentation]\label{def:data_augmentation}
	Let $\{(\mathbf{x}_i, y_i)\}_{i=1}^N$ be a training set and $G$ a group of transformations acting on the input space $\mathcal{X}$ preserving the labels, that is, such that the correct label of $g \cdot \mathbf{x}$ coincides with that of $\mathbf{x}$ for all $g \in G$. \textit{Data augmentation} consists of replacing the original set by the augmented set
	\[
		\{(g \cdot \mathbf{x}_i,\; y_i) : (\mathbf{x}_i, y_i) \in \{(\mathbf{x}_j, y_j)\}_{j=1}^N,\; g \in G\},
	\]
	whose cardinality is $|G| \cdot N$. In practice, a finite subset of $G$ is sampled in each training epoch.
\end{definition}

Data augmentation acts as a form of implicit regularization: it forces the model to learn representations invariant with respect to the group $G$. However, as will be analyzed in \autoref{subsec:equivarianza_invarianza}, this strategy is inefficient compared to the equivariant architectures of \gls{GDL}, which incorporate the symmetries directly into their structure.

We next analyze the theoretical advantages that justify the use of deep architectures.

\subsubsection{Advantages of depth}

Deep networks offer several fundamental advantages over wide shallow networks, supported by both theoretical and empirical results:

\begin{enumerate}
	\item \textbf{Representation efficiency}:

	There exist functions that exhibit an exponential separation between deep and shallow networks. A classic example is the parity function:

	\begin{example}[Parity function]
		The parity function \(f: \{0,1\}^n \to \{0,1\}\) defined as:
		\[
			f(x_1, \ldots, x_n) = \bigoplus_{i=1}^n x_i = \sum_{i=1}^n x_i \mod 2
		\]
		requires \(O(2^n)\) neurons in a 2-layer network, but only \(O(n)\) neurons in a network of \(O(\log n)\) layers.
	\end{example}

	\begin{theorem}[Exponential separation~\citep{telgarsky2016benefitsdepthneuralnetworks}]
		For any integer \(k \geq 1\), there exist functions computable by networks of depth \(\Theta(k^3)\) with \(\Theta(1)\) neurons per layer, that any network of depth \(O(k)\) would need \(\Omega(2^k)\) neurons to approximate with constant error.
	\end{theorem}

	\begin{proof}[Sketch of the proof]
		We define the base triangular function using ReLU neurons:
		\[
			\varphi(x) = 2\operatorname{ReLU}(x) - 4\operatorname{ReLU}\!\left(x - \tfrac{1}{2}\right) + 2\operatorname{ReLU}(x - 1),
		\]
		which forms a ``tent'' on $[0,1]$ with a peak at $x = \tfrac{1}{2}$. We consider the iterated compositions $\varphi_1 = \varphi$, $\varphi_{k+1} = \varphi \circ \varphi_k$. The key observation is that each composition doubles the number of oscillations: $\varphi_k$ has exactly $2^k$ oscillations on $[0,1]$, since $\varphi$ maps $[0,1] \to [0,1]$ and each arc of $\varphi_k$ generates two arcs in $\varphi_{k+1}$.

		Since each composition requires $O(1)$ additional neurons, the function $\varphi_k$ is computable by a network of depth $O(k)$ and constant width. However, a network with ReLU activation and $p$ total neurons produces a piecewise-linear function with at most $O(p^L)$ segments for depth $L$. For a network of depth $O(k)$, approximating $2^k$ oscillations requires at least $\Omega(2^k)$ neurons. The construction extends to the $\Theta(k^3)$ vs $O(k)$ case of the statement by means of more refined techniques for counting linear segments. See \citet{telgarsky2016benefitsdepthneuralnetworks} for the complete proof.
	\end{proof}

	\autoref{fig:profundidad_vs_anchura} visually illustrates this exponential separation between shallow and deep networks.

	\begin{figure}[h]
	\centering
	\begin{tikzpicture}[
	    neuron/.style={circle, draw, minimum size=0.3cm, fill=gdlblue!20},
	    arrow/.style={->, thick, >=stealth}
	]

	% Red superficial (izquierda)
	\node[font=\bfseries] at (-3, 3) {Shallow network (2 layers)};
	\node[font=\footnotesize] at (-3, 2.5) {$O(2^n)$ neurons};

	% Entrada
	\foreach \i in {1,2,3} {
	    \node[neuron, fill=gdlgreen!30] (si\i) at (-4.5, 1.5-\i*0.5) {};
	}
	\node at (-4.5, -0.7) {$\vdots$};

	% Capa oculta (muchas neuronas)
	\foreach \i in {1,...,8} {
	    \node[neuron] (sh\i) at (-3, 2-\i*0.4) {};
	}
	\node at (-3, -1.5) {$\vdots$};
	\node[font=\scriptsize] at (-3, -2) {$2^n$};

	% Salida
	\node[neuron, fill=gdlred!30] (so) at (-1.5, 0) {};

	% Conexiones (simplificadas)
	\draw[gdlgray, thin] (si1) -- (sh1);
	\draw[gdlgray, thin] (si2) -- (sh4);
	\draw[gdlgray, thin] (si3) -- (sh7);
	\draw[gdlgray, thin] (sh1) -- (so);
	\draw[gdlgray, thin] (sh8) -- (so);

	% Red profunda (derecha)
	\node[font=\bfseries] at (3, 3) {Deep network ($\log n$ layers)};
	\node[font=\footnotesize] at (3, 2.5) {$O(n)$ neurons};

	% Capas
	\foreach \l in {0,...,4} {
	    \foreach \i in {1,2,3} {
	        \pgfmathsetmacro{\col}{ifthenelse(\l==0,"gdlgreen!30",ifthenelse(\l==4,"gdlred!30","gdlblue!20"))}
	        \node[neuron, fill=\col] (d\l\i) at (1.5+\l*0.8, 1-\i*0.6) {};
	    }
	}

	% Conexiones entre capas
	\foreach \l in {0,...,3} {
	    \pgfmathtruncatemacro{\nextl}{\l+1}
	    \foreach \i in {1,2,3} {
	        \foreach \j in {1,2,3} {
	            \draw[gdlgray, very thin] (d\l\i) -- (d\nextl\j);
	        }
	    }
	}

	% Etiquetas de complejidad
	\node[font=\small, align=center] at (-3, -3) {Complexity:\\exponential};
	\node[font=\small, align=center] at (3, -3) {Complexity:\\polynomial};

	\end{tikzpicture}
	\caption[Depth vs width in neural networks]{Comparison between a wide shallow network and a narrow deep network for computing the same function. The shallow network requires exponentially more neurons, while the deep network achieves the same expressiveness with polynomial growth.}
	\label{fig:profundidad_vs_anchura}
	\end{figure}

	\item \textbf{Hierarchical learning and compositionality}:

	Deep networks exploit the natural compositional structure of many problems:

	\begin{itemize}
		\item \textbf{Vision}: pixels \(\to\) edges \(\to\) textures \(\to\) parts \(\to\) objects
		\item \textbf{Language}: characters \(\to\) morphemes \(\to\) words \(\to\) phrases \(\to\) semantics
		\item \textbf{Audio}: samples \(\to\) phonemes \(\to\) words \(\to\) concepts
	\end{itemize}

	This hierarchy is reflected mathematically in the composition of functions:
	\[
		f(x) = f_L \circ f_{L-1} \circ \cdots \circ f_1(x)
	\]
	where each \(f_i\) learns incrementally more abstract transformations.

	\item \textbf{Feature reuse and parameter efficiency}:

	Depth allows sharing computations between different processing paths. Consider a network that must compute \(m\) different functions \(\{g_1, \ldots, g_m\}\):

	\begin{itemize}
		\item Shallow network: \(O(m \cdot n)\) parameters (each function independent)
		\item Deep network: \(O(n + m)\) parameters (shared features + specific heads)
	\end{itemize}

	\item \textbf{Implicit inductive bias}:

	Deep networks have an implicit bias toward ``simple'' functions in the sense of Kolmogorov complexity:

	\begin{definition}[Function complexity]\label{def:complejidad_funcion}
		For a neural network \(f_\theta\) with fixed architecture, we define its complexity as:
		\[
			\mathcal{C}(f_\theta) = \|\theta\|_0 + \sum_{l=1}^L \log(\text{fan-in}_l \cdot \text{fan-out}_l)
		\]
	\end{definition}

	Deep networks favor solutions with lower \(\mathcal{C}(f_\theta)\), which corresponds to more ``natural'' functions according to the distribution of real-world data.

	\item \textbf{Optimization and geometry of the loss function}:

	Counterintuitively, deeper networks can be easier to optimize:

	\begin{itemize}
		\item The local minima in overparameterized networks tend to have error similar to the global one
		\item Depth induces a hierarchical structure in the geometry of the loss function
		\item ``Lazy training'' phenomenon: deep networks find simple solutions first
	\end{itemize}
\end{enumerate}

\subsection{Modern architectures}

The development of deep learning has produced specialized architectures, each designed to exploit specific structures in the data. Many of these architectures have already been introduced in previous sections of this chapter; here we offer a comparative view focused on their distinctive characteristics and on how they relate to one another. In later chapters we will explore in depth how geometric ideas transform some of these architectures.

\autoref{fig:arquitecturas_modernas} presents a visual comparison of the four main deep learning architectures, showing their distinctive characteristics in terms of information flow and connectivity structure.

\begin{figure}[h]
\centering
\begin{tikzpicture}[scale=1.05, every node/.style={scale=1.0}]
% Estilos comunes
\tikzset{
    neuron/.style={circle, draw, minimum size=0.4cm, fill=#1!40},
    layer/.style={rectangle, draw, rounded corners, minimum width=0.8cm, minimum height=1.8cm, fill=#1!20},
    kernel/.style={rectangle, draw, minimum width=0.3cm, minimum height=0.3cm, fill=gdlorange!50},
    arrow/.style={->, thick, >=stealth},
    title/.style={font=\small\bfseries, above}
}

% === FEEDFORWARD ===
\begin{scope}[shift={(-6,0)}]
    \node[title] at (0.6,2.5) {Feedforward};
    % Capas
    \foreach \y in {0,0.6,1.2,1.8} {
        \foreach \x in {0,0.4,0.8,1.2} {
            \node[neuron=gdlblue] at (\x,\y) {};
        }
    }
    % Conexiones (algunas representativas)
    \foreach \x in {0,0.4,0.8,1.2} {
        \draw[gdlgray!50, thin] (\x,0.15) -- (0.2,0.45);
        \draw[gdlgray!50, thin] (\x,0.15) -- (0.6,0.45);
        \draw[gdlgray!50, thin] (\x,0.15) -- (1.0,0.45);
    }
    \draw[arrow, gdlblue!70] (0.6,-0.3) -- (0.6,0) node[below, pos=0, font=\tiny] {input};
    \draw[arrow, gdlblue!70] (0.6,1.95) -- (0.6,2.2) node[above, pos=1, font=\tiny] {output};
    \node[below, font=\tiny, text=gdlgray] at (0.6,-0.6) {Unidirectional flow};
\end{scope}

% === CNN ===
\begin{scope}[shift={(-2,0)}]
    \node[title] at (0.8,2.5) {CNN};
    % Feature maps (rectángulos apilados)
    \draw[fill=gdlgreen!20, draw=gdlgreen!60] (0,0) rectangle (1.2,1.6);
    \draw[fill=gdlgreen!30, draw=gdlgreen!60] (0.15,0.15) rectangle (1.05,1.35);
    \draw[fill=gdlgreen!40, draw=gdlgreen!60] (0.3,0.3) rectangle (0.9,1.1);
    % Kernel
    \draw[fill=gdlorange!50, draw=gdlorange!80] (0.4,0.5) rectangle (0.7,0.8);
    \node[font=\tiny] at (0.55,0.65) {$k$};
    % Flecha de convolución
    \draw[arrow, gdlorange!80, dashed] (0.55,0.9) -- (0.55,1.5);
    \draw[arrow, gdlgreen!70] (0.6,-0.3) -- (0.6,0) node[below, pos=0, font=\tiny] {image};
    \draw[arrow, gdlgreen!70] (0.6,1.7) -- (0.6,2.2) node[above, pos=1, font=\tiny] {features};
    \node[below, font=\tiny, text=gdlgray] at (0.6,-0.6) {Weight sharing};
\end{scope}

% === RNN ===
\begin{scope}[shift={(2,0)}]
    \node[title] at (0.6,2.5) {RNN};
    % Celdas en secuencia
    \foreach \i/\x/\t in {0/0/$t{-}1$, 1/0.8/$t$, 2/1.6/$t{+}1$} {
        \node[layer=gdlred, minimum height=1cm, minimum width=0.5cm] (rnn\i) at (\x,0.8) {};
        \node[font=\tiny] at (\x,0.8) {\t};
    }
    % Conexiones recurrentes
    \draw[arrow, gdlred!70] (0.3,0.8) -- (0.5,0.8);
    \draw[arrow, gdlred!70] (1.1,0.8) -- (1.3,0.8);
    % Loop
    \draw[arrow, gdlred!70, thick] (rnn1.north) .. controls (0.8,1.6) and (1.2,1.6) .. (rnn2.north);
    \node[font=\tiny, gdlred!70] at (1.2,1.8) {$h_t$};
    % Entradas
    \foreach \x in {0,0.8,1.6} {
        \draw[arrow, gdlgray] (\x,0) -- (\x,0.3);
    }
    \node[below, font=\tiny] at (0.8,-0.3) {sequence};
    \node[below, font=\tiny, text=gdlgray] at (0.8,-0.6) {Hidden state};
\end{scope}

% === TRANSFORMER ===
\begin{scope}[shift={(6,0)}]
    \node[title] at (0.8,2.5) {Transformer};
    % Tokens
    \foreach \x in {0,0.5,1,1.5} {
        \node[neuron=gdlpurple, minimum size=0.3cm] (t\x) at (\x,0) {};
    }
    % Attention layer
    \draw[fill=gdlpurple!20, draw=gdlpurple!60] (-0.2,0.6) rectangle (1.7,1.2);
    \node[font=\tiny] at (0.75,0.9) {Attention};
    % Attention connections (all-to-all)
    \foreach \x in {0,0.5,1,1.5} {
        \draw[gdlpurple!40, thin] (\x,0.15) -- (0.75,0.55);
    }
    % Output
    \foreach \x in {0,0.5,1,1.5} {
        \node[neuron=gdlpurple, minimum size=0.3cm] at (\x,1.6) {};
        \draw[gdlpurple!40, thin] (0.75,1.25) -- (\x,1.45);
    }
    \draw[arrow, gdlpurple!70] (0.75,-0.5) -- (0.75,-0.15) node[below, pos=0, font=\tiny] {tokens};
    \draw[arrow, gdlpurple!70] (0.75,1.75) -- (0.75,2.2) node[above, pos=1, font=\tiny] {output};
    \node[below, font=\tiny, text=gdlgray] at (0.75,-0.8) {Global attention};
\end{scope}

\end{tikzpicture}
\caption[Comparison of deep learning architectures]{Comparison of the four main deep learning architectures. \textbf{Feedforward}: fully connected layers with unidirectional flow. \textbf{CNN}: convolutions with shared kernels that preserve spatial structure. \textbf{RNN}: recurrent connections that maintain a temporal hidden state. \textbf{Transformer}: attention mechanisms that connect all elements simultaneously.}
\label{fig:arquitecturas_modernas}
\end{figure}

\paragraph{Feedforward networks.}
Feedforward networks, based on dense layers (\autoref{def:capa_densa}), represent the most basic architecture where information flows unidirectionally from the input to the output. A network with $L$ layers is defined as the composition $f(x) = f_L \circ f_{L-1} \circ \cdots \circ f_1(x)$, where each $f_i(x) = \sigma(W_i x + b_i)$ is an affine transformation followed by a nonlinearity. The absence of cycles guarantees that the computation is deterministic and efficient.

\paragraph{Convolutional networks.}
\glspl{CNN}, which we will study in detail in \autoref{ch:cnn}, exploit local structure and translational invariance by means of convolution operations. For an input $x$ and a kernel $k$, the discrete convolution (\autoref{def:discrete_convolution}) is defined as:
\[
	(x * k)[i] = \sum_j x[j] \cdot k[i-j].
\]
Convolutional layers preserve spatial locality and drastically reduce the number of parameters through weight sharing, resulting in translational equivariance: $T_a[f(x)] = f(T_a[x])$ where $T_a$ is a translation (\autoref{theorem:convolution_translation_equivariance}).

\paragraph{Recurrent networks.}
\glspl{RNN} (\autoref{subsubsec:rnn}) process sequences by means of a hidden state $h_t$ that evolves according to:
\[
	h_t = \sigma(W_{hh} h_{t-1} + W_{xh} x_t + b_h).
\]
This recurrence allows capturing temporal dependencies, although it suffers from the vanishing gradient problem (\autoref{subsubsec:vanishing_gradient}). Variants such as \gls{LSTM} (\autoref{subsubsec:lstm}) and \gls{GRU} introduce gating mechanisms that allow learning long-term dependencies.

\paragraph{Attention mechanisms.}\label{par:atencion_arquitecturas}
Attention mechanisms, which we formalize below and extend in \autoref{ch:attention}, allow modeling long-range dependencies without recurrence. The central idea is to dynamically weight the contributions of different parts of the input according to their relevance:
\[
	\text{Attention}(q, K, V) = \sum_{i=1}^n \alpha_i v_i, \quad \alpha_i = \frac{\exp(s(q, k_i))}{\sum_{j=1}^n \exp(s(q, k_j))}
\]
where $s(q, k_i)$ is a scoring function that measures the relevance between the query $q$ and the key $k_i$ (see \autoref{def:generalized_attention}). This formulation is invariant to permutations of the inputs and generalizes operations such as pooling and convolution. The quadratic complexity $O(n^2)$ in sequence length has motivated the development of efficient variants. Attention mechanisms form the basis of the \gls{Transformer} architecture.

\begin{definition}[Self-Attention layer]\label{def:self_attention_basica}
Let $\{h_1, \ldots, h_n\} \subset \mathbb{R}^d$ be a sequence of input vectors. A \textit{self-attention layer} with key dimension $d_k$ is defined by means of three linear projections $Q, K, V \in \mathbb{R}^{d \times d_k}$ (query, key, value) and computes:
\[
\mathrm{SelfAttention}(h_i) = \sum_{j=1}^n \alpha_{ij}\, V h_j, \quad \alpha_{ij} = \frac{\exp\!\left(\frac{(Qh_i)^T(Kh_j)}{\sqrt{d_k}}\right)}{\sum_{m=1}^n \exp\!\left(\frac{(Qh_i)^T(Kh_m)}{\sqrt{d_k}}\right)}
\]
where the coefficients $\alpha_{ij} \in (0,1)$ with $\sum_j \alpha_{ij} = 1$ are called \textit{attention weights}.
\end{definition}

\paragraph{Generative adversarial networks.}
\glspl{GAN} consist of two competing networks: a generator $G: \mathcal{Z} \to \mathcal{X}$ and a discriminator $D: \mathcal{X} \to [0,1]$. The minimax objective is:
\[
	\min_G \max_D\; \mathbb{E}_{x \sim p_{\text{data}}}[\log D(x)] + \mathbb{E}_{z \sim p_z}[\log(1 - D(G(z)))].
\]
At the Nash equilibrium, $G$ generates samples indistinguishable from real data.

\paragraph{Graph networks.}
\glspl{GNN} generalize neural operations to non-Euclidean data, a central topic of \autoref{ch:gdl}. Message passing (\autoref{def:mpnn_unificador}) in a graph $\mathcal{G} = (\mathcal{V}, \mathcal{E})$ is expressed as:
\[
	h_v^{(k+1)} = \sigma\!\left(W_{\text{self}}h_v^{(k)} + \sum_{u \in \mathcal{N}(v)} W_{\text{msg}}h_u^{(k)}\right),
\]
where $\mathcal{N}(v)$ denotes the neighbors of node $v$. This formulation preserves the permutation invariance of the graph.

\paragraph{Diffusion models.}
Diffusion models learn to reverse a process of gradual degradation. Given a forward process that adds Gaussian noise:
\[
	q(x_t \mid x_{t-1}) = \mathcal{N}\!\left(x_t;\, \sqrt{1-\beta_t}\, x_{t-1},\, \beta_t \mathbf{I}\right),
\]
the neural network learns the inverse process $p_\theta(x_{t-1} \mid x_t)$ to generate samples from pure noise. The objective function is related to the estimation of the score $\nabla_x \log p(x)$.
